\documentclass{article}
\usepackage[english]{babel}
\usepackage{geometry,amsmath,amssymb,enumerate,hyperref,xcolor,latexsym,theorem,bbm}
\geometry{a4paper}

%%%%%%%%%% Start TeXmacs macros
\newcommand{\Tau}{\mathrm{T}}
\newcommand{\assign}{:=}
\newcommand{\asterisk}{\mathord{*}}
\newcommand{\backassign}{=:}
\newcommand{\llangle}{{\langle\!\langle}}
\newcommand{\mathd}{\mathrm{d}}
\newcommand{\mathi}{\mathrm{i}}
\newcommand{\nobracket}{}
\newcommand{\tmaffiliation}[1]{\\ #1}
\newcommand{\rrangle}{{\rangle\!\rangle}}
\newcommand{\tmcolor}[2]{{\color{#1}{#2}}}
\newcommand{\tmop}[1]{\ensuremath{\operatorname{#1}}}
\newcommand{\tmscript}[1]{\text{\scriptsize{$#1$}}}
\newcommand{\tmtextbf}[1]{{\bfseries{#1}}}
\newcommand{\tmtextit}[1]{{\itshape{#1}}}
\newenvironment{enumerateroman}{\begin{enumerate}[i.] }{\end{enumerate}}
\newenvironment{proof}{\noindent\textbf{Proof\ }}{\hspace*{\fill}$\Box$\medskip}
\newenvironment{proof*}[1]{\noindent\textbf{#1\ }}{\hspace*{\fill}$\Box$\medskip}
\newtheorem{theorem}{Theorem}
\newtheorem{corollary}[theorem]{Corollary}
\newtheorem{definition}[theorem]{Definition}
\newtheorem{lemma}[theorem]{Lemma}
\newtheorem{proposition}[theorem]{Proposition}
{\theorembodyfont{\rmfamily}\newtheorem{remark}[theorem]{Remark}}
\newcommand{\tmkeywords}{\textbf{Keywords:} }
%%%%%%%%%% End TeXmacs macros

\newcommand{\nbmax}[1]{\tmcolor{blue}{[Max] #1}}
\newcommand{\nbluigi}[1]{\tmcolor{blue}{[Luigi] #1}}
\newcommand{\wip}[1]{\tmcolor{red}{[WIP] #1}}
%


\begin{document}

\title{
  Grassmannian stochastic analysis\\
  and the stochastic quantization\\
  of Euclidean Fermions
}

\author{Sergio Albeverio, Luigi Borasi,\\ Francesco C. De Vecchi, Massimiliano Gubinelli\\[.5em]
  \tmaffiliation{Hausdorff Center of Mathematics \&\\
  Institute of Applied Mathematics\\
  University of Bonn, Germany}\\ 
  \\
  \tt{albeverio@iam.uni-bonn.de}\\
 \tt{ borasi@iam.uni-bonn.de}\\
  \tt{francesco.devecchi@uni-bonn.de}\\
 \tt{ gubinelli@iam.uni-bonn.de}
}
%\date{April 2020}

\maketitle

\begin{abstract}
  We introduce a stochastic analysis of Grassmann random variables suitable
  for the stochastic quantization of Euclidean fermionic quantum field
  theories. Analysis on Grassmann algebras is developed from the point of view
  of quantum probability: a Grassmann random variable is a homomorphism of an
  abstract Grassmann algebra into a quantum probability space, i.e. the von
  Neumann algebra of bounded operators on a Hilbert space endowed with a
  suitable state. On the way we define the notion of Gaussian processes,
  Brownian motion and stochastic differential equations taking values in
  Grassmann algebras. We use them to study the long time behavior of Langevin
  Grassmann stochastic differential equations driven by white noise and
  describe their invariant measures. As an application we give a proof of the
  stochastic quantization and of the removal of the infrared cutoff for the
  Euclidean Yukawa model.\\[.4em]

{\small 
\noindent \tmkeywords{Grassmann algebra, Euclidean Fermions, stochastic quantization,
non-commutative probability, Yukawa model.}}

\end{abstract}

{\tableofcontents}

\section{Introduction}

The aim of this work is to develop a stochastic quantization of Euclidean
fermionic quantum field theories (QFTs). This involves defining and studying
the long time behavior of stochastic equations in Grassmann algebras.

\

Euclidean fermion fields form a Grassmann
algebra~{\cite{osterwalder_euclidean_1973,osterwalder_euclidean_1973-2}}.
Their correlation functions (Schwinger functions) are typically represented by
Berezin integrals~{\cite{kirillov_introduction_1987}} of the type
\begin{equation}
  \langle O (\psi, \bar{\psi}) \rangle = \frac{\int \mathd \psi \mathd
  \bar{\psi} O (\psi, \bar{\psi}) e^{- S_E (\psi, \bar{\psi})}}{\int \mathd
  \psi \mathd \bar{\psi} e^{- S_E (\psi, \bar{\psi})}} \label{eq:eqft-berezin}
\end{equation}
where the fields $\psi, \bar{\psi}$ form a Grassmann algebra, $O$ is a
functional of $\psi$ and $\bar{\psi}$ generating the \ Schwinger functions,
and $S_E (\psi, \bar{\psi})$ is the Euclidean action, usually the sum of a
quadratic part and a polynomial interaction (possibility involving Boson
fields, which will have to be averaged according to an appropriate probability
measure)~{\cite{feldman_renormalizable_1986,disertori_continuous_1998,mastropietro_non-perturbative_2008}}.
Without entering too much into the precise details, let us point out here that
the formula~{\eqref{eq:eqft-berezin}} is a purely algebraic expression, devoid
of any analytic meaning and is used as a convenient bookkeeping for various
kind of expansions used to define and study the properties of the Schwinger
functions. This is due to the fact that Grassmann algebras do not have a
natural notion of positivity, nor a norm, associated with them.

\

Stochastic quantization (SQ) is an approach to the construction of the
correlation functions of an Euclidean QFT that has been introduced by Parisi
and Wu~{\cite{parisi_perturbation_1981}}. The basic idea is to introduce an
additional variable (usually a fictitious time) and interpret the Euclidean
fields $\varphi$, to be constructed, as solutions to stochastic differential
equations involving the mentioned additional variable and a source of noise
$\xi$, most commonly Gaussian. One of the advantages of this approach is that
the fields of the theory are expressed via the stochastic evolutions as
functionals of the noise, i.e. $\varphi = \varphi (\xi)$, which then allows to
study them as functionals on $\xi$ exploiting detailed information about the
law of $\xi$. In this way even interacting fields can be studied
non-perturbatively exploiting the Gaussian structure of $\xi$. Moreover,
another conceptual advantage of stochastic quantization is that it is
insensitive to questions of absolute continuity with respect to a reference
measure. This, on the other hand, is usually the main difficulty of the
Gibbsian point of view as expressed in equation~{\eqref{eq:eqft-berezin}}.

\

This strategy has been employed to construct bosonic theories, starting with
the work~{\cite{jona_lasinio_stochastic_1985}} on the stochastic quantization
of the $\Phi^4_2$ model for a scalar particle in two dimension with quartic
interaction. More recently the work of Hairer on regularity
structures~{\cite{hairer_theory_2014}} opened the way to the study of the
three dimensional $\Phi^4$ model, see
also~{\cite{catellier_paracontrolled_2013,kupiainen_renormalization_2016}}.
These new perspectives on this and related problems led to a series of new
results on the global space-time control of the stochastic
dynamics~{\cite{MWcomedown,gubinelli_global_2019,albeverio_invariant_2017,moinat_space_time_2018}}
and to a novel proof of the constructions of non-Gaussian bosonic Euclidean
quantum field theories in three dimensions~{\cite{gubinelli_pde_2018}}.

\

One of main differences between the bosonic and the fermionic case is the
necessity, in the fermionic setting, of considering stochastic processes in a
non-commutative framework.

\

The interest in studying time evolutions and stochastic processes with values
in non commutative algebras arose in connection with the development of
quantum mechanics in the 20s--30s of last century and later in quantum field
theory, particularly around the 50s. In these theories one associates physical
observables with operators evolving in time and acting on a, in general
infinite dimensional, complex Hilbert space. Mathematically, this lead to the
representation theory of certain operator algebras (CCR-algebras resp.
CAR-algebras in the case of fermions) that are particular cases of more
general operator algebras.

\

In this line quantum fields require the study of certain infinite dimensional
(non-commutative) algebras, the study of which has been particularly
stimulated by developments in the period 1940-1960, where the principles of
systems with an infinite number of degrees of freedom, in particular quantum
fields and statistical mechanics has been worked out mathematically. Let us
mention the first works on von Neumann algebras by Murray and
von~Neumann~{\cite{Murray1,Murray2,Murray3,VonNeumann}} and the work on
$C^{\asterisk}$-algebras developed originally and quite independently by I. M.
Gelfand and his school (see, e.g.,~{\cite{Gelfand}}), since the beginnings of
the 40s, and by I. E. Segal and his school (see~{\cite{Segal1,Segalbook}}).
Both types of algebras are associative and equipped with an involution (the
adjoint operation). Whereas $C^{\ast}$-algebras are looked upon as abstract
normed algebras, without necessarily being realized as algebras of operators
acting in Hilbert spaces , the study of von Neumann algebras (also called
$W^{\ast}$-algebras) rather involves looking at them as self-adjoint algebras
of bounded operators in Hilbert spaces,closed in the weak operator topology.
Both algebras can be looked upon as particular cases of normed algebras, that
in the case of closedness are called Banach algebras. More precisely,
$C^{\ast}$-algebras are normed algebras with a unit $\mathbb{I}$ and an
involution $\cdot^{\ast}$, any element $A$ of which satisfies $\| A A^{\ast}
\| = \| A \|^2$. Although in principle every von~Neumann algebra can be looked
upon as a $C^{\ast}$-algebra, structural and conceptual developments in the
study of von Neumann algebras involve particularly their realizations in
Hilbert space. The mathematical literature on normed, Banach, $C^{\ast}$ and
$W^{\ast}$ algebras is very rich and there are a number of excellent
monographs on them, see, e.g., Naimark~{\cite{Naimark}} and
Zelazko~{\cite{MR0448079}} for normed spaces and algebras, and Kadison and
Ringrose~{\cite{Kadison1,Kadison2}} (and references therein) for the other
type of algebras.

\

For applications in mathematical physics, in particular quantum mechanics,
quantum field theory and statistical mechanics, there are monographs by
Bratteli and Robinson~{\cite{Bratteli1,Bratteli2}}, Emch~{\cite{Emch}},
Baumg{\"a}rtel~{\cite{Baumgartel1}}, as well as Baumg{\"a}rtel and
Wollenberg~{\cite{Baumgartel2}}, Segal~{\cite{Segalbook2}}; we refer
particularly to the books by Baez, Segal and Zhou~{\cite{Segalbook}}, and
Haag~{\cite{Haag1}}, and also the articles by
Summers~{\cite{summers2012perspective}} and Doplicher~{\cite{Doplicher}}, that
surveys particularly some recent developments related to such algebras in
quantum field theory.

Before describing the approach we use in this article, which is based on
stochastic processes forming a Grassmann algebra, let us mention some works
about stochastic quantization of fermionic fields not using Grassmann algebras
but more connected with ideas of the theory of noncommutative processes and
Markov semigroups. They involve noncommutative Brownian motion and stochastic
analysis based on on it, in particular also SDEs but with no space variable
dependence, hence no SPDEs. Pioneering work in this direction has been done by
Accardi in the 70s, see e.g.~{\cite{accardi_outline_2015}}. Work extending
SDEs with values in algebras other than Grassmann algebras, like $C^{\ast}$
and von Neumann algebras are by authors like Streater and
coworkers~{\cite{Streater1,Streater2,Streater3}}, Applebaum and
Hudson~{\cite{Applebaum1,Applebaum2}}, Belavkin~{\cite{Belavkin}}, Carlen and
Kree~{\cite{Carlen1}}, Gordina~{\cite{Gordina}}, Sinha and
Goswami~{\cite{Sinha}}, K{\"u}mmer {\cite{Kummer}} see also the monographs by
Meyer~{\cite{Meyerbook}}, Holevo~{\cite{Holevo}}, and the recent survey
article by Cipriani~{\cite{Cipriani}} and the references therein.

\

In this paper we focus on stochastic quantization of interacting fermions,
realizing the fermionic fields as non-commutative operator solution of a
suitable stochastic differential equation taking values in a Grassmann
sub-algebra of the Banach algebra of all bounded operators acting on a Hilbert
space (see the precise definition in Section~\ref{sec:grassmann-rv}). Here let
us mention that Grassmann algebras constitute a class of associative non
commutative algebras that were discovered by Grassmann in the second half of
the 19th century. They are connected with Clifford algebras, that in turn are
the underlying structure to the study of CCR and CAR algebras. All such
algebras are nowadays part of the general theory of algebra as pioneered by
E.~Noether in the 20s-40s, see e.g.
Van~der~Waerden~{\cite{van2003algebra1,van2003algebra2}}.

\

In our paper we use stochastic processes forming a Grassmann algebra, rather
than Clifford, $C^{\ast}$ or $W^{\ast}$ algebras, since they better relate to
work done in the context of interesting models of constructive quantum field
theory involving fermions, in which methods based on integrals on non
commutative structures like Berezin integrals are involved. Such methods are
quite essential in most oft the physics literature on fermionic quantum
fields, including renormalization theory, see, e.g., Feldman, Kn{\"o}rrer and
Trubowitz~{\cite{MR1905424}}, Salmhofer~{\cite{salmhofer_renormalization}} or
Zinn-Justin~{\cite{ZinnJustin}}. As for fermions (i.e fields associated with
particles of half integer spins, like e.g. electrons) stochastic quantization
ideas have been discussed heuristically, after Parisi and Wu seminal paper, in
the physics literature, starting with Kakudo et al.~{\cite{kakudo1983gauge}},
Fukai et al.~{\cite{fukai_stochastic_1983}}, Damgaard and
Tsokos~{\cite{damgaard_stochastic_1984}} and She-Sheng and
Ting-Chang~{\cite{she-sheng_stochastic_1985}} \ up to the recent paper of
Efremov~{\cite{efremov_stochastic_2019}}. These authors stress particularly
the relations with Grassmann valued random variables, solving partial
differential equations with Grassmann valued Gaussian white noise. The
invariant measures of physical interest are here described by averages with
respect to Berezin integrals, see~{\cite{Berezin1}} (further developed into a
analysis on Grassmann algebras by the same author, in the form of
super-analysis~{\cite{Berezin2}}; for super-algebras see also, e.g.,
Rogers~{\cite{Rogers2007}} and Le\u{\i}tes~{\cite{Leites}}, on a more heuristic
level also DeWitt~{\cite{DeWitt}}). Berezin integrals have been useful in
calculating fermionic systems, see e.g. Zinn-Justin~{\cite{ZinnJustin}},
Izkyson and Zuber~{\cite{Itzykson}}.

\

Rigorous attempts to define Grassmann stochastic quantization were undertaken
in a series of papers by Scherbakov et
al.~{\cite{shcherbakov_elements_1993,shcherbakov_elements_1993-1,shcherbakov_elements_1993-2}}
using the locally convex topology generated by considering all the possible
correlation functions of the fields and the noise to generate a family of
semi-norms. Convergence in this setting corresponds to convergence of all the
correlation functions. \ This approach has the disadvantage of not allowing
the introduction of norms strong enough to study differential equations via
standard approaches. However it should be said that these papers define and
successfully solve equations with finitely many degrees of freedom (i.e. with
small scale and large scale cut-offs) and then address the convergence when
the large scale cut-off is removed (infinite volume limit) using cluster
expansion techniques. Other results in this direction were obtained later on
by Ignatyuk et al.~{\cite{ignatyuk_convergence_1993}}, with applications in
statistical mechanics.

\

Our approach here is different and, in our opinion, both simpler and more
effective. We consider this problem in the context of non-commutative
probability (sometimes referred to as \tmtextit{quantum} probability), namely
we will define the relevant Grassmannian objects as non-commutative random
variables. A non-commutative probability space $(\mathcal{A}, \omega)$ is
given by a $C^{\ast}$ algebra $\mathcal{A}$ and a state $\omega$, which is a
linear positive functional, on $\mathcal{A}$. Inspired by the general approach
of Accardi et al.~{\cite{accardi_quantum_1982}} (see
also~{\cite{accardi_outline_2015}}) we will define random variables with
values in a Grassmann algebra $\Lambda$ \ as homomorphisms of $\Lambda$ into
$\mathcal{A}$ (note however that we will not require the homomorphisms to
respect the involution, since there is no canonical candidate for that in
Grassmann algebras).

\

It is useful to keep in mind how classical commutative random variables fit
in this algebraic approach. Consider a random variable $X : \Omega \rightarrow
\mathcal{M}$ taking values in a manifold $\mathcal{M}$ and let $\mathcal{B}$
be the algebra of all measurable, complex bounded function on $\mathcal{M}$.
Then all the relevant probabilistic information about $X$ is encoded in the
homomorphisms of algebras $f \in \mathcal{B} \mapsto f (X) \in \mathcal{A}$
where $\mathcal{A}$ is the algebra of bounded measurable complex random
variables on $\Omega$ endowed with the linear functional given by the
expectation $\mathbb{E}$. In particular this is the point of view with which
SDEs on manifolds are usually defined and
studied~{\cite{hsu_stochastic_2002}}.

\

The embedding of $\Lambda$ into $\mathcal{A}$ allows to use the norm of
$\mathcal{A}$ to do analysis on Grassmann algebras. Of course the associated
analysis will not be \tmtextit{canonical} from the point of view of the
Grassmann algebra but seems enough to allow us to obtain a satisfactory theory
of stochastic quantization. Again, an analogy can be helpful to understand
this setting: the standard approach to study Gaussian processes in Hilbert
spaces~{\cite{Gross1967}} (see also, e.g.,~{\cite{janson_gaussian_1997}})
requires to have a non-canonical embedding of the Hilbert space into a larger
Banach space which supports the Gaussian measure. In our case, the role of
``bigger'' Banach space is played by a suitable CAR sub-algebra of
$\mathcal{A}$ which provides us with enough Grassmann generators to realize
our homomorphisms and has a canonical notion of norm. It will be for us just
``a place where to hang our (analytic) hat on''. \ \



Another approach to infinite dimensional Grassmann algebras (and the related
probability theory) we are aware of is that of
Rogers~{\cite{rogers_graded_1986}} which essentially consists in looking at
all the finite dimensional sub-algebras, use the Berezin integral (see,
e.g,~{\cite{kirillov_introduction_1987}}) to compute averages and require
certain natural consistency conditions. It is interesting to note however that
our approach follows Osterwalder and
Schrader~{\cite{osterwalder_euclidean_1973,osterwalder_euclidean_1973-2}}
which were the first to do rigorous analysis of Euclidean Fermions: they were
using the language of Fock space operators and not of Berezin calculus, and
this allowed them to consider rigorously expressions involving infinite
dimensional Grassmann algebras. The Fock space approach has also been used in
discussing supersymmetry in the context of Parisi--Sourlas dimensional
reduction~{\cite{klein_supersymmetry_1983,klein_supersymmetry_1985}} and
Euclidean stochastic
quantization~{\cite{albeverio_elliptic_2018,albeverio_elliptic_2019,de_vecchi_note_2019}}.

\

{\noindent}\tmtextbf{Plan of the paper.} In Section~\ref{sec:grassmann-rv} we
define Grassmann random variables and various related notions. We then
introduce Gaussian variables, including white noise and Brownian motion.
Section~\ref{sec:finite-dim} studies finite dimensional Grassmann differential
equations with additive noise and prove certain properties of their
correlation functions, showing in particular that we can construct Grassmann
Langevin diffusions which have prescribed stationary correlations.
Section~\ref{sec:infinite-dim} is devoted to extend the previous results to
SDE in infinite dimensions of the kind relevant for stochastic quantization.
This allows us to use these processes to represent the correlation functions
of Euclidean fermionic theories like~{\eqref{eq:eqft-berezin}}. As an
application we provide in Section~\ref{sec:yukawa} a construction of the
fermionic sector of the massive Yukawa Euclidean quantum field theory in two
dimensions, in the weak coupling region and with an ultraviolet cutoff. In
Appendix~\ref{appendix_Besov} we collect various technical results about
functional spaces of Banach-valued functions, while in
Appendix~\ref{app:series} we include some side results about the convergence
of series expansions for the solution to finite dimensional SDEs.

\

{\noindent}\tmtextbf{Acknowledgements.} This work is supported by DFG via the
grant AL 214/50-1 \ ``Invariant measures for SPDEs and Asymptotics''.

\

{\noindent}\tmtextbf{Notations.} We denote by $\mathcal{L} (A, B)$ the space
of linear maps between vector spaces $A$ and $B$, if both $A, B$ have
topologies we consider all the maps to be continuous. We let $\mathcal{L} (A)
=\mathcal{L} (A, A)$. With $\tmop{Hom} (A, B)$ we denote homomorphisms between
algebras $A$ and $B$.

\section{Grassmann random variables}\label{sec:grassmann-rv}

\tmtextit{Algebraic probability.} We consider a complex Hilbert space
$\mathcal{H}$ (to be fixed later) and denote by $\mathcal{A}=\mathcal{L}
(\mathcal{H})$ the von Neumann algebra of bounded operators with the operator
norm. Moreover we assume to have a state $\omega$ on $\mathcal{A}$, that is a
positive definite functional normalized by the condition that $\omega
(\mathbb{I}) = 1$, where $\mathbb{I}$ is the identity in $\mathcal{L}
(\mathcal{H})$. The pair $(\mathcal{A}, \omega)$ is a (non-commutative or
algebraic) probability space. We do not require the state $\omega$ to be
either faithful, or tracial.

\

\tmtextit{Grassmann algebras.} Let $V$ be a (finite dimensional or infinite
dimensional and separable) real Hilbert space. Denote by $\Lambda V$ the
Grassmann algebra generated by $V$, i.e. the exterior algebra obtained by
quotienting the tensor algebra $T (V) = \oplus_{n \geqslant 0} V^{\otimes n}$
by the two-sided ideal generated by elements of the form $x \otimes x$ for $x
\in V$. We denote the product of two elements $f, g \in \Lambda V$ by $f
\wedge g$ (or if there is no confusion with other products simply by $f g$).

\

Let $(v_{\alpha})_{\alpha}$ be a fixed basis of $V$. As a consequence
$\Lambda V$ is spanned by elements of the form $v_A \assign v_{\alpha_1}
\wedge \cdots \wedge v_{\alpha_n}$ and $A = (\alpha_1, \ldots, \alpha_n)$ is a
multi-index with the convention that $v_{\emptyset} = 1$ is the unit element
in $\Lambda V$. When it is clear from the context we will denote the product
in $\Lambda V$ simply by $v_A = v_{\alpha_1} \ldots v_{\alpha_n}$. We recall
that $\Lambda V$ is $\mathbb{Z}_2$ graded in the sense that it splits into odd
and even parts $\Lambda V = \Lambda_{\tmop{odd}} V \oplus
\Lambda_{\tmop{even}} V$. On $\Lambda V$ there is an Hopf algebra structure
with coproduct $\Delta : \Lambda V \rightarrow \Lambda V \otimes \Lambda V$
where $\Lambda V \otimes \Lambda V$ is the $\mathbb{Z}_2$-graded tensor
product algebra which satisfies
\[ (f \otimes g) (h \otimes k) = (- 1)^{| g |  | h |} f h \otimes g k, \qquad
   f, g, h, k \in \Lambda V, \]
with $| \cdot | : \Lambda V \rightarrow \{ 0, 1 \}$ the even/odd grading. The
coproduct $\Delta$ is the algebra homomorphism such that $\Delta v = 1 \otimes
v + v \otimes 1$ for all $v \in V \subseteq \Lambda V$ and the counit
$\varepsilon : \Lambda V \rightarrow \mathbb{C}$ given by $\varepsilon (v_A)
=\mathbbm{1}_{A = \emptyset}$.

\

\tmtextit{Random variables.}

\begin{definition}
  A $V$-Grassmann random variable $\Psi$ is an algebra homomorphism from the
  Grassmann algebra $\Lambda V$ into $\mathcal{A}$. We denote by $\mathcal{G}
  (V) = \tmop{Hom} (\Lambda V, \mathcal{A})$ the set of all such
  homomorphisms. We will always assume that $\Psi \in \mathcal{G} (V)$ comes
  with a unitary element $U_{\Psi}$ of $\mathcal{A}$ such $U_{\Psi} \Psi (v)
  U_{\Psi}^{\ast} = \Psi (- v)$ for all $v \in V$. We call the \tmtextit{law}
  of $\Psi \in \mathcal{G} (V)$ the family of its moments $\omega^{\Psi} (F) :
  = \omega (\Psi (F))$ for all $F \in \Lambda V$, also represented by the
  linear functional $\omega^{\Psi} : \Lambda V \rightarrow \mathbb{R}$.
\end{definition}

Note that $\Psi \in \mathcal{G} (V)$ cannot be assumed a to be
$\ast$-algebra-homomorphism since $\Lambda V$ has no (natural)
$\ast$-operation. If $F \in \Lambda V$ has representation $F = \sum_A F_A v_A$
we shall employ the following dual notation
\[ F (\Psi) \assign \Psi (F) = \sum_A F_A \Psi^A, \]
where $\Psi^A \assign \Psi (v_A)$ and $F_A \in \mathbb{R}$. Since $\Psi$ is
assumed to be an algebra homomorphism, we have e.g. $\Psi^{\alpha}
\Psi^{\beta} = - \Psi^{\beta} \Psi^{\alpha}$, where $\Psi^{\alpha} = \Psi
(v_{\alpha})$ and $(v_{\alpha})_{\alpha}$ is a fixed basis of $V$. Moreover,
$\Psi (\Lambda V)$ is a Grassmann sub-algebra of $\mathcal{A}$ and $\Psi (F)$
has the same degree (even or odd) as $F$. It is important to note that, as
shown by arguments in \ {\cite{van2003algebra1,van2003algebra2}}, even if some
arguments are formulated in a basis dependent way the definition of $\Lambda
V$ and its characterization by anti-commutation relations is basis
independent.

\

\tmtextit{Compatibility.} When the context is clear we will abbreviate $X \in
\mathcal{G} (V)$ as $X \in \mathcal{G}$. Let $X \in \mathcal{G} (V)$ and $Y
\in \mathcal{G} (W)$, we say that they are \tmtextit{compatible} if the linear
map $Z : V \oplus W \rightarrow \mathcal{A}$ given by $Z (v) = X (v)$ if $v
\in V$ and $Z (w) = Y (w)$ if $w \in W$, extends to a homomorphism $Z :
\Lambda (V \oplus W) \rightarrow \mathcal{A}$. In this case we write $Z \in
\mathcal{G} (V \oplus W)$ or briefly $(X, Y) \in \mathcal{G}$. Note that we
have a super- algebra isomorphism $\Lambda (V \oplus W) \approx \Lambda (V)
\otimes \Lambda (W)$ where tensor product is in the sense of super-algebras.
From this isomorphism we get in particular $Z (F) = m_{\mathcal{A}} [(X
\otimes Y) (F)]$ for all $F \in \Lambda (V \oplus W)$ where $m_{\mathcal{A}} :
\mathcal{A} \otimes \mathcal{A} \rightarrow \mathcal{A}$ denotes the
multiplication of $\mathcal{A}$. Compatibility can of course be defined for
$X_1, \ldots, X_n \in \mathcal{G}$ in a corresponding way. We shall express
that $X_1, \ldots, X_n \in \mathcal{G}$ are compatible by writing $(X_1,
\ldots, X_n) \in \mathcal{G}$.

\begin{remark}
  The notion of compatibility which we introduce here in not standard in
  algebraic probability. A related concept is that of kinematic independence
  (see e.g. the review of Accardi~{\cite{accardi_outline_2015}}).
\end{remark}

Given $X, Y \in \mathcal{G} (V)$ which are compatible we define $X + Y \in
\mathcal{G} (V)$ as
\[ (X + Y) (F) = F (X + Y) \assign m_{\mathcal{A}} [(X \otimes Y) \Delta F],
   \qquad F \in \Lambda V, \]
Note indeed that by compatibility $Z = m_{\mathcal{A}} \circ (X \otimes Y) :
\Lambda (V \oplus V) \approx \Lambda V \otimes \Lambda V \rightarrow
\mathcal{A}$ is an algebra homomorphism and therefore \
\[ (X + Y) (F G) = Z (\Delta (F G)) = Z (\Delta F \Delta G) = Z (\Delta F) Z
   (\Delta G) = (X + Y) (F) (X + Y) (G) . \]
The notation is justified from the fact that $(X + Y) (v) = X (v) + Y (v)$ for
all $v \in V$.

\

\tmtextit{Independence.} If $(X_1, \ldots, X_n) \in \mathcal{G} (V_1 \oplus
\cdots \oplus V_n)$ are compatible Grassmann variables with values in the
probability space $(\mathcal{A}, \omega)$, then we say that $X_1, \ldots, X_n$
are (tensor) independent (with respect to the state $\omega$) if, for all $F_j
\in \Lambda V_j$, we have that~{\cite{dirksen_noncommutative_2010}}
\[ \omega \left( \prod_{j = 1}^k \vphantom{\frac{}{}} X_j (F_j) \right) =
   \prod_{j = 1}^k \omega (X_j (F_j)) . \]


For example if $(X, Y) \in \mathcal{G} (V \oplus W)$ are (compatible and)
independent Grassmann random variables with values in $(\mathcal{A}, \omega)$,
then for all $v \in V \approx V \oplus 0, w \in W \approx 0 \oplus W$ we have
$v w \approx ((v \oplus 0) \wedge (0 \oplus w)) \in \Lambda (V \oplus W)
\approx \Lambda V \otimes \Lambda W$ and
\[ \omega ((X \otimes Y) (v w)) = \omega (X (v) Y (w)) = \omega (X (v)) \omega
   (Y (w)) . \]
Let $X_1 \in \mathcal{G} (V_1)$, $X_2 \in \mathcal{G} (V_2)$ be two Grassmann
random variables. We want to show how to construct two compatible and
independent random variables $\tilde{X}_1, \tilde{X}_2$ with the same law as
$X_1, X_2$. The random variables $X_1$ and $X_2 $ have values in the
non-commutative probability spaces $(\mathcal{L} (\mathcal{H}_1), \omega_1)$
and $(\mathcal{L} (\mathcal{H}_2), \omega_2)$ respectively. By a GNS argument
we can assume that $\omega_1$ is the state associated to a vector $\Omega_1$
in $\mathcal{H}_1$. Moreover we also assume that $\mathcal{H}_1 =
\overline{X_1 (\Lambda V) \Omega_1}$ (the bar denotes the closure with respect
to the Hilbert space topology) since this enough to compute all the quantities
$\omega_1 (X_1 (F))$ for all $F \in \Lambda V$. Then on $\mathcal{H}_1$ we can
introduce an involution $R_1$ which act as $R_1 X_1 (F) \Omega_1 = X_1 ((-
1)^{| F |} F) \Omega_1$ where $| \cdot |$ is the grading on $\Lambda V$. A
similar construction can be done for the other space $\mathcal{H}_2$. Now we
take $\mathcal{H}=\mathcal{H}_1 \otimes \mathcal{H}_2$, $\tilde{X}_1 (v) = X_1
(v) \otimes R_2$ for $v \in V_1$ and $\tilde{X}_2 (w)
=\mathbbm{1}_{\mathcal{H}_1} \otimes X_2 (w)$ for $w \in V_2$. Then
$\tilde{X}_1 (v) \tilde{X}_2 (w) = - \tilde{X}_2 (w) \tilde{X}_1 (v)$ and
since $\Lambda (V_1 \oplus V_2) \approx \Lambda V_1 \otimes \Lambda V_2$ we
have also that the map $X (v \oplus w) = X_1 (v) + X_2 (w)$ with $v \in V_1, w
\in V_2$ can be extended to a homomorphism from $\Lambda (V_1 \oplus V_2)$ to
$\mathcal{A}$ in such a way that $\omega (\tilde{X}_i (F)) = \omega_i (X_i
(F))$ for all $F \in \Lambda V_i$, $i = 1, 2$.

\subsection{Topology and calculus on Grassmann variables}

\

Before beginning we want to specify the topology we consider on $\mathcal{G}
(V)$. We consider here only the case where $V$ is finite dimensional, since
there is a more or less unique natural topology on $\mathcal{G} (V)$. Some
topologies of $\mathcal{G} (V)$ when $V$ is infinite dimensional are discussed
in Section \ref{sec:infinite-dim}.

\

When $V$ is finite dimensional all norms on $V$ are equivalent. For this
reason we choose the norm induced by the pre-Hilbert space inner product
$\langle \cdot, \cdot \rangle$ (that in the case of $V$ finite dimensional is
a Hilbert space inner product on $V$) related to the construction of
$V$-valued Gaussian random variables $X$ (see
Section~\ref{sec:gaussians}).\quad In this case $\mathcal{G} (V)$ has a
natural metric topology given by the distance
\begin{equation}
  d_{\mathcal{G} (V)} (X, Y) \assign \| X - Y \|_{\mathcal{G} (V)} = \sup_{v
  \in V, | v |_V = 1} \| X (v) - Y (v) \|_{\mathcal{A}},
  \label{eq:distrance-g}
\end{equation}
where $\| \cdot \|_{\mathcal{A}}$ is the natural norm in the
$\asterisk$-algebra $\mathcal{A}$.

\begin{remark}
  \label{remark_infinitenorm}In principle the definition of the distance
  {\eqref{eq:distrance-g}} and the related norm $\| \cdot \|_{\mathcal{G}
  (V)}$ does not use in any way the fact that $V$ is finite dimensional. For
  this reason in the following, when $V$ is a (in general infinite
  dimensional) per-Hilbert space we will use the notation $\| X
  \|_{\mathcal{G} (V)}$ for the quantity $\sup_{v \in V, | v |_V = 1} \| X (v)
  - Y (v) \|_{\mathcal{A}}$.
\end{remark}

When $V$ is finite dimensional $\mathcal{G} (V)$ is a complete metric space
with respect to $d_{\mathcal{G} (V)}$, in fact we have:

\begin{lemma}
  \label{lemma_completeness}The metric $d_{\mathcal{G} (V)}$ makes
  $\mathcal{G} (V)$ a complete metric space.
\end{lemma}

\begin{proof}
  We let $\| \cdot \|_{\mathcal{G} (V)}$ to be the restriction to $\mathcal{G}
  (V)$ of the natural norm of the Banach space $\mathcal{L} (V, \mathcal{A})$
  of the continuous linear maps between $V$ and $\mathcal{A}$, then the
  quantity
  \[ d_{\mathcal{G} (V)} (X, Y) = \| X - Y \|_{\mathcal{G} (V)} \]
  where $(X - Y) (v) = X (v) - Y (v)$, satisfies the usual properties of a
  metric and if $X_n$ is a Cauchy sequence with respect to $d_{\mathcal{G}
  (V)}$, there is an element $X \in \mathcal{L} (V, \mathcal{A})$ such that
  $X_n$ converges to $X$ in $\mathcal{L} (V, \mathcal{A})$. The only thing
  that we have to prove is that the linear map $X \in \mathcal{L} (V,
  \mathcal{A})$ can be extended to a homomorphism from $\Lambda V$ to
  $\mathcal{A}$. This is equivalent to prove that for any $v, w \in V$ we have
  $X (v) X (w) = - X (w) X (v)$ where the product is the natural one in
  $\mathcal{A}$. On the other hand by the continuity of the product of
  $\mathcal{A}$ with respect to $\| \cdot \|_{\mathcal{A}}$ we have
  \[ X (v) X (w) = \lim_{n \rightarrow + \infty} X_n (v) X_n (w) = - \lim_{n
     \rightarrow + \infty} X_n (w) X_n (v) = - X (w) X (v) . \]
\end{proof}

\begin{remark}
  The distance $d_{\mathcal{G} (V)}$ is not the unique reasonable choice in
  $\mathcal{G} (V)$ and it depends on the topology chosen on $\mathcal{A}$ (in
  this case we choose the topology of the uniform converge of operators). An
  essential choice to preserve the non-linear structure of $\mathcal{G} (V)$,
  exploited in the proof of Lemma~\ref{lemma_completeness}, seems to be that
  the product on $\mathcal{A}$ is continuous with respect to this topology. 
\end{remark}

In the following we will denote $\| Y \|_{\mathcal{G} (V)}$ for $Y \in
\mathcal{G} (V)$ the positive number
\[ \| Y \|_{\mathcal{G} (V)} \assign \| Y \|_{\mathcal{L} (V, \mathcal{A})} =
   \sup_{v \in V, \| v \| = 1} \| X (v) \|_{\mathcal{A}}, \]
i.e. the norm of the restriction of $Y$ to $\mathcal{L} (V, \mathcal{A})$.
This is useful to formulate a Taylor formula on $\mathcal{G} (V)$. We will sue
the simpler notation $m$ for the multiplication $m_{\mathcal{A}}$ in
$\mathcal{A}$.

\

Define the \tmtextit{right derivative} $\partial_R : \Lambda V \rightarrow
\Lambda V \otimes V$ by
\begin{equation}
  \partial_R F \assign (\mathbbm{1} \otimes \Pi_V) (\Delta F) : \Lambda V
  \rightarrow \Lambda V \otimes V \label{eq:derivativeR}
\end{equation}
where $\Pi_V : \Lambda V \rightarrow V$ is the projection from the tensor
algebra $\Lambda V$ onto $V$. Then
\[ F (X + Y) - F (X) - m [(X \otimes Y) (\partial_R F)] = m [(X \otimes Y)
   (\mathbbm{1}-\mathbbm{1} \otimes \varepsilon -\mathbbm{1} \otimes \Pi_V)
   \Delta F] \]
for $F \in \Lambda V$, where we used that $F (X) = m [(X \otimes Y)
(\mathbbm{1} \otimes \varepsilon) \Delta F]$. Given that
\[ (\mathbbm{1}-\mathbbm{1} \otimes \varepsilon -\mathbbm{1} \otimes \Pi_V)
   \Delta F \in \Lambda V \otimes \Lambda_{\geqslant 2} V \]
where $\Lambda_{\geqslant 2} V$ denotes the subspace of $\Lambda V$ of
elements of degree $\geqslant 2$. We can define recursively the $k$+1-th
derivative as
\[ \partial_R^{k + 1} = (\partial_R \otimes 1 \otimes \cdots \otimes 1)
   \partial_R^k : \Lambda V \rightarrow \Lambda V \otimes V \otimes \cdots
   \otimes V \]
where $\partial_R^k$ is the $k$-th order derivative. Note that the right
derivative $\partial_R : \Lambda V \rightarrow \Lambda V \otimes V$ satisfies
\[ \partial_R (f_1 \cdots f_n) = \sum_{k = 1}^n (- 1)^{n - k} \left( f_1
   \cdots \not{f_k} \cdots f_n \right) \otimes f_k, \qquad f_1, \ldots, f_n
   \in V. \]
Of course one can define also a left derivative $\partial_L = (\Pi_V \otimes
\mathbbm{1}) \circ \Delta : \Lambda V \rightarrow V \otimes \Lambda V$ with
similar properties and
\[ \partial_L (f_1 \cdots f_n) = \sum_{k = 1}^n (- 1)^{k - 1} \left( f_1
   \cdots \not{f_k} \cdots f_n \right) \otimes f_k, \qquad f_1, \ldots, f_n
   \in V. \]
We will consider $\Lambda V \otimes V$ as a $\Lambda V$-bimodule and in
particular we define the bilinear form $\langle \cdot, \cdot \rangle :
(\Lambda V \otimes V) \otimes V \rightarrow \Lambda V$ by
\begin{equation}
  \langle f \otimes v, w \rangle = f \langle v, w \rangle, \qquad f \in
  \Lambda V, v, w \in V. \label{eq:bimodule}
\end{equation}
Note that there is no ambiguity on whether $\langle \cdot, \cdot \rangle$
denotes this bilinear form or whether it denotes the scalar product of $V$.

\begin{remark}
  \label{remark_infinitederivative}The definition of derivative $\partial_R^k$
  and the bilinear for {\eqref{eq:bimodule}} do not depend on the fact that
  $V$ is finite dimensional. For this reason we can define $\partial_R^k$ and
  the bilinear for {\eqref{eq:bimodule}} also when $V$ is a infinite
  dimensional pre-Hilbert space.
\end{remark}

\begin{remark}
  Note that if $F_{\tmop{odd}} \in \Lambda_{\tmop{odd}} V$ then $\partial_R
  F_{\tmop{odd}} = \partial_L F_{\tmop{odd}}$. On the other hand, if
  $F_{\tmop{even}} \in \Lambda_{\tmop{even}} V$ then $\partial_R
  F_{\tmop{even}} = - \partial_L F_{\tmop{even}}$. In general, if we denote by
  $\mathbb{P}_{\tmop{odd}}, \mathbb{P}_{\tmop{even}}$ respectively the
  projection onto the odd and even parts of $\Lambda V$, then $\partial_R =
  (\mathbb{P}_{\tmop{even}} -\mathbb{P}_{\tmop{odd}}) \partial_L$.
\end{remark}

We can now state the version of Taylor formula which we will use below.

\begin{lemma}
  \label{lemma:taylor}Consider $F \in \Lambda V$ and let $X, Y \in \mathcal{G}
  (V)$ be two compatible Grassmann random variables such that $\| Y
  \|_{\mathcal{G} (V)} \leqslant 1$, then
  \begin{equation}
    F (X + Y) = F (X) + \sum_{k = 1}^n \frac{1}{k!} m [(X \otimes Y \otimes
    \cdots \otimes Y) (\partial_R^k F)] + O (\| Y \|^{n + 1}_{\mathcal{G}
    (V)}) \label{eq:taylor}
  \end{equation}
\end{lemma}

\begin{proof}
  We have
  \[ F (X + Y) = m [(X \otimes Y) (\Delta F)] = \sum_{k \geqslant 0} m [(X
     \otimes Y) (1 \otimes \Pi_{\Lambda^k V}) (\Delta F)], \]
  where the sum is finite since $F \in \Lambda V$ is a finite polynomial.
  Equation~{\eqref{eq:taylor}} easily follows as soon as we prove that
  \[ m [(X \otimes Y \otimes \cdots \otimes Y) (\partial_R^k F)] = (k!) m [(X
     \otimes Y) (1 \otimes \Pi_{\Lambda^k V}) (\Delta F)] . \]
  We have
  \begin{eqnarray*}
    \partial^k_R F & = & (\partial_R \otimes 1 \otimes \cdots \otimes 1)
    \cdots (\partial_R \otimes 1) \partial_R F\\
    & = & \left( \vphantom{\tfrac{a}{b}} (1 \otimes \Pi_V) \Delta \otimes 1
    \otimes \cdots \otimes 1 \right) \cdots \left( \vphantom{\tfrac{a}{b}} (1
    \otimes \Pi_V) \Delta \otimes 1 \right) \left( \vphantom{\tfrac{a}{b}} 1
    \otimes \Pi_V \right) \Delta F\\
    & = & \left( \vphantom{\tfrac{a}{b}} (1 \otimes \Pi_V) \otimes 1 \otimes
    \cdots \otimes 1 \right) (\Delta \otimes 1 \otimes \cdots \otimes 1)
    \cdots \left( \vphantom{\tfrac{a}{b}} (1 \otimes \Pi_V) \otimes 1 \right)
    (\Delta \otimes 1) \left( \vphantom{\tfrac{a}{b}} 1 \otimes \Pi_V \right)
    \Delta F\\
    & = & \left( \vphantom{\tfrac{a}{b}} (1 \otimes \Pi_V) \otimes 1 \otimes
    \cdots \otimes 1 \right) \cdots \left( \vphantom{\tfrac{a}{b}} (1 \otimes
    \Pi_V) \otimes 1 \right) \left( \vphantom{\tfrac{a}{b}} 1 \otimes \Pi_V
    \right)  (\Delta \otimes 1 \otimes \cdots \otimes 1) \cdots (\Delta
    \otimes 1) \Delta F\\
    & = & \left( \vphantom{\tfrac{a}{b}} 1 \otimes \Pi_V \otimes \Pi_V
    \otimes \cdots \otimes \Pi_V \right)  (\Delta \otimes 1 \otimes \cdots
    \otimes 1) \cdots (\Delta \otimes 1) \Delta F.
  \end{eqnarray*}
  Let $N \assign \{ 1, \ldots, n \}$, $\#N = n$. For $I = \{ j_1, \ldots, j_k
  \} \subset N$, we put $v_I \assign v_{j_1} \cdots v_{j_k}$ for $v_j \in V, j
  \in N$. Let $\tmop{sgn} (N\backslash I)$ be a sign such that $v_N =
  \tmop{sgn} (N\backslash I) v_{N\backslash I} v_I$ By definition of the
  coproduct $\Delta$ in the exterior algebra $\Lambda V$, we have
  \[ \begin{array}{c}
       (1 \otimes \Pi_V \otimes \cdots \otimes \Pi_V) (\Delta \otimes 1
       \otimes \cdots \otimes 1) \cdots (\Delta \otimes 1) (\Delta) (v_1
       \cdots v_n) =\\
       = \sum_{j_1, \ldots j_n \in N} \tmop{sgn} (N\backslash I)
       v_{N\backslash \{ j_1, \ldots, j_n \}} \otimes v_{j_1} \otimes \cdots
       \otimes v_{j_n} .
     \end{array} \]
  On the other hand
  \[ \begin{array}{lll}
       (1 \otimes \Pi_{\Lambda^k V}) \Delta v_1 \cdots v_n & = & (1 \otimes
       \Pi_{\Lambda^k V}) \sum_{I \subset N} \tmop{sgn} (N\backslash I) v_{N
       \setminus I} \otimes v_I\\
       & = & \frac{1}{k!} \sum_{j_1, \ldots, j_k} \tmop{sgn} (N\backslash I)
       v_{N\backslash \{ j_1, \ldots, j_k \}} \otimes v_{\{ j_1, \ldots, j_k
       \}} .
     \end{array} \]
  By linearity, we therefore have
  \[ (1 \otimes m^{(k)}) (\partial^k_R F) = k! (1 \otimes \Pi_{\Lambda^k V})
     (\Delta F) . \]
  where $m^{(k)} : V^{\otimes k} \rightarrow \Lambda V$ is the product. Since
  $X, Y$ are algebra homomorphisms they commute with the product hence we have
  \begin{eqnarray*}
    m [(X \otimes Y \otimes \cdots \otimes Y) (\partial_R^k F)] & = & m [(X
    \otimes Y) (1 \otimes m^{(k)}) (\partial_R^k F)]\\
    & = & k!m [(X \otimes Y) (1 \otimes \Pi_{\Lambda^k V}) (\Delta F)] .
  \end{eqnarray*}
  
  
  \
  
  This concludes the proof.
\end{proof}

\subsection{Grassmann Gaussian variables}\label{sec:gaussians}

Let now $V$ be a real pre-Hilbert space with scalar product $\langle \cdot,
\cdot \rangle$ and with an antisymmetric bounded operator $C : V \rightarrow
V$. We recall that, by Remark \ref{remark_infinitenorm} and Remark
\ref{remark_infinitederivative}, we can extend the definitions of $\| \cdot
\|_{\mathcal{G} (V)}$, $\partial_R^k$ and the bilinear form
{\eqref{eq:bimodule}} from the finite dimensional case to the generic
(infinite) pre-Hilbert space $V$.

\begin{definition}
  A ($V$-)Grassmann (centered) Gaussian variable with correlation $C$ is a
  random variable $X \in \mathcal{G} (V)$ such that
  \begin{equation}
    \omega (X (F) X (f)) = \omega (X (\langle \partial_R F, C f \rangle)),
    \qquad F \in \Lambda V, f \in V. \label{eq:ibp-r}
  \end{equation}
  We also require that $\| X \| < \infty$, i.e that the map $X : V \rightarrow
  \mathcal{A}$ must be continuous with respect the topology induced on $V$ by
  the pre-Hilbert product structure and the uniform operator topology on
  $\mathcal{A}$.
\end{definition}

If we define $\partial_R X (F) = X (\partial_R F)$, with the understanding
that $X (g \otimes w) = X (g) \otimes w$, $g \in V$ and $w \in \Lambda V$,
then we have the more familiar expression (similar to the bosonic counterpart)
\[ \omega (F (X) X (h)) = \omega (\langle \partial_R F (X), C h \rangle) . \]
Note that the integration by parts formula determine all the moments of the
Gaussians. In particular~{\eqref{eq:ibp-r}} implies that $\omega (X (f_1)
\cdots X (f_n)) = 0$ if $n$ is odd and if $n = 2 k$ is even
\begin{equation}
  \omega (X (f_1) \cdots X (f_{2 k})) = \sum_{\sigma} (- 1)^{\sigma} \prod_{i
  = 1}^k \langle f_{\sigma (2 i - 1)}, C f_{\sigma (2 i)} \rangle_V
  \label{eq:wick}
\end{equation}
where the sum runs over all the pairings $\sigma$ of $\{ 1, \ldots, 2 k \}$
and $(- 1)^{\sigma}$ is an appropriate sign. The right hand side of
{\eqref{eq:wick}} can be written as a Pfaffian:
\[ \omega (X (f_1) \cdots X (f_n)) = \underset{1 \leqslant i, j \leqslant
   n}{\tmop{Pf}} \langle f_i, C f_j \rangle_V, \]
where the Pfaffian is defined, for an anti-symmetric $n \times n$ matrix $M$,
by the relation
\[ (\underset{1 \leqslant i, j \leqslant n}{\tmop{Pf}} M_{i j})^2 =
   \underset{1 \leqslant i, j \leqslant n}{\det} M_{i j} . \]
The existence of a Gaussian variable, in the above sense, implies the
inequality
\[ \underset{1 \leqslant i, j \leqslant n}{\det} \langle f_i, C f_j \rangle_V
   = [\omega (X (f_1) \cdots X (f_n))]^2 \leqslant \| X (f_1) \|^2 \cdots \| X
   (f_n) \|^2 \leqslant \| X \|^{2 n} \| f_1 \|^2 \cdots \| f_n \|^2 \]
well known in the mathematical physics literature relative to fermionic
expansions e.g. see~{\cite{mastropietro_non-perturbative_2008}}.

\begin{remark}
  The averages of Gaussian variables depend only on the quadratic form $(f, g)
  \mapsto \langle f, C g \rangle$, however analysis on the Grassmann algebra
  relies on the scalar product $\langle \cdot, \cdot \rangle$ also via the
  requirement that $\| X (f) \|_{\mathcal{A}} \leqslant \| f \|_V$. In
  particular the realization of the Grassmann algebra as a family of bounded
  operators is not canonical. This is the reason we need to require the above
  continuity of the map $X : V \rightarrow \mathcal{A}$.
\end{remark}

In order to construct Grassmann Gaussian variables, we need ``a place to hang
the hat on'', this place will be a (canonical commuting relations) CAR algebra
endowed with its vacuum state. \ To allow for arbitrary correlations $C$ we
can use a standard trick~{\cite{osterwalder_euclidean_1973-2}} which consists
in doubling the generators of the CAR algebra with respect to the generators
of the Grassmann algebra.

\begin{lemma}
  For every antisymmetric and bounded $C : V \rightarrow V$ there exists a
  ($V$-)Grassmann Gaussian variable $X$ with correlation $C$ (on a suitable
  probability space $(\mathcal{A}, \omega)$).
\end{lemma}

\begin{proof}
  Let us consider $\Lambda V$ itself as a real pre-Hilbert space with respect
  to the scalar product $\llangle \cdot, \cdot \rrangle$ on $\Lambda V$ given
  by
  \[ \llangle v_1 \wedge \cdots \wedge v_n, w_1 \wedge \cdots \wedge w_m
     \rrangle = \delta_{n m} \underset{1 \leqslant j, k \leqslant n}{\det}
     \langle v_j, w_k \rangle, \]
  for $v_j, w_k \in V$, $j = 1, \ldots, n$, $k = 1, \ldots, m$, where $\langle
  \cdot, \cdot \rangle$ denotes as usual the scalar product on $V$. Note that
  $\llangle \cdot, \cdot \rrangle$ on $\Lambda^n V$ is simply the restriction
  of the Hilbert scalar product on the tensor product $V^{\otimes n}$. Let
  $\mathcal{H}$ be the completion of $\Lambda V$ with respect to $\llangle
  \cdot, \cdot \rrangle$ and denote by $\Omega$ the element in $\mathcal{H}$
  which corresponds to $1 \in \Lambda V$. Denote by $\lambda : \Lambda V
  \rightarrow \tmop{End} (\Lambda V)$ the left action of $\Lambda V$ on itself
  given by $\lambda (F) G = F \wedge G$, where $F, G \in \Lambda V$. We show
  that this action extends to a representation of $\Lambda V$ on
  $\mathcal{H}$. Indeed $\lambda (v)$, $v \in V$ corresponds to a creation
  operator. Let us denote by $\lambda (v)^{\Tau}$ the adjoint of the operator
  $\lambda (v)$ with respect to the scalar product $\llangle \cdot, \cdot
  \rrangle$. A standard computation shows that for $v \in V$ we have
  \[ \lambda (v)^{\Tau} w_1 \wedge \cdots \wedge w_n = \sum_{\ell = 1}^n (-
     1)^{\ell - 1} \langle v, w_{\ell} \rangle w_1 \wedge \cdots \wedge
     \not{w_{\ell}} \wedge \cdots \wedge w_n . \]
  Namely $\lambda^T (v) x$ can be expressed in terms of the left derivative as
  $\lambda^T (v) x = \langle v, \partial_L x \rangle$ (recall the definition
  given in equation~{\eqref{eq:bimodule}}) for all $x \in \Lambda V$. In
  particular we have
  \begin{equation}
    \{ \lambda (v), \lambda (w) \} = \{ \lambda (v)^{\Tau}, \lambda (w)^{\Tau}
    \} = 0, \qquad \{ \lambda (v)^{\Tau}, \lambda (w) \} = \langle v, w
    \rangle \mathbb{I}, \qquad v, w \in V, \label{eq:car}
  \end{equation}
  where as usual $\{ \cdot, \cdot \}$ denotes the anticommutator. As a
  consequence we obtain, for any $x \in \Lambda V$
  \[ \llangle \lambda (v) x, \lambda (v) x \rrangle + \llangle \lambda
     (v)^{\Tau} x, \lambda (v)^{\Tau} x \rrangle = \llangle x, \{ \lambda
     (v)^{\Tau}, \lambda (v) \} x \rrangle = \langle v, v \rangle \llangle x,
     x \rrangle . \]
  In particular $\lambda (v), \lambda (v)^{\Tau}, v \in V$, extend to bounded
  operators on $\mathcal{H}$ (still denoted by the same symbol). We now define
  $X : \Lambda V \rightarrow \mathcal{L} (\mathcal{H})$ to be the algebra
  homomorphism defined first on $V$ by
  \[ X (v) = \lambda (C v) + \lambda (v)^{\Tau}, \qquad v \in V. \]
  Then $X$ is extended in a natural way to $\Lambda V$ (still denoted by the
  same symbol). Note that $X (x)$, $x \in \Lambda V$, is indeed a bounded
  operator since it is the sum of products of bounded operators. Finally we
  define $\omega : \mathcal{L} (\mathcal{H}) \rightarrow \mathbb{R}$ to be the
  state on $\mathcal{A} =\mathcal{L} (\mathcal{H})$ defined by $\Omega$, that
  is
  \[ \omega (T) \assign \llangle \Omega, T \Omega \rrangle, \qquad T \in
     \mathcal{L} (\mathcal{H}) . \]
  We claim that $X$ is a Gaussian Grassmann random variable on $(\mathcal{L}
  (\mathcal{H}), \omega)$. We first show that it is Grassmann:
  \[ \begin{array}{lll}
       \{ X (v), X (w) \} & = & \{ \lambda (C v), \lambda (w)^{\Tau} \} + \{
       \lambda (v)^{\Tau}, \lambda (C w) \}\\
       & = & \langle C v, w \rangle + \langle v, C w \rangle = \langle v,
       C^{\Tau} w \rangle + \langle v, C w \rangle = 0,
     \end{array} \]
  where $v, w \in \Lambda$, because by assumption $C^{\Tau} = - C$. Now we
  have also that
  \begin{eqnarray*}
    \omega (X (v) X (w)) & = & \llangle \Omega,^{} \lambda (v)^{\Tau} \lambda
    (C w) \Omega \rrangle\\
    & = & \llangle \lambda (v) \Omega, \lambda (C w) \Omega \rrangle =
    \llangle v, C w \rrangle,
  \end{eqnarray*}
  where the first equality follows from the fact that $\lambda (v)^{\Tau}
  \Omega = 0$. More generally using that $X (v) \Omega = \lambda (C v) \Omega$
  and the commutation relation~{\eqref{eq:car}} we have that, for all $f_1,
  \ldots, f_{2 k} \in V$
  \[ \omega (X (f_1) \cdots X (f_{2 k})) = \omega (X (f_1) \cdots X (f_{2 k -
     1}) \lambda (C f_{2 k})) \]
  \[ = - \omega (X (f_1) \cdots \lambda (C f_{2 k}) X (f_{2 k - 1})) + \omega
     (X (f_1) \cdots X (f_{2 k - 2})) \langle f_{2 k - 1}, C f_{2 k} \rangle
  \]
  \[ = \cdots = \sum_{\ell = 1}^{2 k - 1} (- 1)^{2 k - 1 - \ell} \omega \left(
     X (f_1) \cdots \not{X (f_{\ell})} \cdots X (f_{2 k - 1}) \right) \langle
     f_{\ell}, C f_{2 k} \rangle \]
  Therefore
  \[ \omega (X (F) X (f)) = \omega (X (\langle \partial_R F, C f \rangle)),
     \qquad F \in \Lambda V, f \in V. \]
  Note that, similarly
  \begin{equation}
    \omega (X (f) X (F)) = \omega (X (\langle f, C \partial_L F \rangle)) .
    \label{eq:ibp-l}
  \end{equation}
  
\end{proof}

\tmtextit{Complex Gaussians.} Later on we will need also complex Grassmann
Gaussian variables. As in the commutative setting, their definition poses no
particular problems. Let $V_{\mathbb{C}}$ be a complex pre-Hilbert space with
a notion of complex conjugation denoted $v \mapsto \bar{v}$ and an Hermitian
scalar product $(\cdot, \cdot)$ which will be assumed to be anti-linear in the
left variable. Consider a complex bilinear antisymmetric form $B_{\mathbb{C}}$
on $V_{\mathbb{C}}$
\[ B_{\mathbb{C}} : V_{\mathbb{C}} \times V_{\mathbb{C}} \rightarrow
   \mathbb{C}, \quad B_{\mathbb{C}} : v \times w \mapsto B_{\mathbb{C}} (v,
   w), \quad v, w \in V_{\mathbb{C}} . \]
Antisymmetric here means simply that $B_{\mathbb{C}} (v, w) = - B_{\mathbb{C}}
(w, v)$, $v, w \in V_{\mathbb{C}}$. Define the standard bilinear form of
$V_{\mathbb{C}}$ to be
\[ \langle \cdot, \cdot \rangle_{\mathbb{C}} : v \times w \mapsto \langle v, w
   \rangle_{\mathbb{C}} \assign (\bar{v}, w), \quad v, w \in V_{\mathbb{C}} .
\]
The suffix $\mathbb{C}$ in $\langle \cdot, \cdot \rangle_{\mathbb{C}}$ is a
reminder that this is a bilinear form with complex values. The antisymmetric
complex bilinear form $B_{\mathbb{C}}$ can be written in terms of a linear
operator $C_{\mathbb{C}}$ in the following way
\[ B_{\mathbb{C}} (v, w) = \langle v, C_{\mathbb{C}} w \rangle_{\mathbb{C}} =
   (\bar{v}, C_{\mathbb{C}} w), \]
where $C_{\mathbb{C}}$ is a complex operator which is antisymmetric with
respect to $\langle \cdot, \cdot \rangle_{\mathbb{C}}$, that is, we have
\[ \langle v, C_{\mathbb{C}} w \rangle = \langle C_{\mathbb{C}}^{\Tau} v, w
   \rangle = - \langle C_{\mathbb{C}} v, w \rangle, \quad v, w \in
   V_{\mathbb{C}} . \]
\begin{definition}
  A $V_{\mathbb{C}}$-complex Grassmann (centered) Gaussian variable with
  complex correlation $C_{\mathbb{C}}$ is a random variable $X \in \mathcal{G}
  (V_{\mathbb{C}})$ such that
  \[ \omega (X (F) X (h)) = \omega (X (\langle \partial_R F, C_{\mathbb{C}} h
     \rangle_{\mathbb{C}})), \qquad F \in \Lambda V_{\mathbb{C}}, h \in
     V_{\mathbb{C}} . \]
  We require that $\| X \|_{\mathcal{G} (V_{\mathbb{C}})} < \infty$, i.e that
  the map $X : V_{\mathbb{C}} \rightarrow \mathcal{A}$ is continuous.
\end{definition}

\tmtextit{White noise.} A relevant example of a random variable taking values
in an infinite dimensional Grassmann algebra is $d$-dimensional (Gaussian)
white noise.

\begin{definition}
  \label{definition_Brownian}A $V$-valued $d$-dimensional (Gaussian) white
  noise with correlation $C : V \rightarrow V$ is the Grassmann Gaussian
  variable $\Xi \in \mathcal{G} (L^2 (\mathbb{R}^d) \otimes V)$ with
  correlation $\tilde{C}$ given by
  \[ (\tilde{C} f) (x) = C f (x) \]
  for all $x \in \mathbb{R}^d$ and $f \in L^2 (\mathbb{R}^d) \otimes V \approx
  L^2 (\mathbb{R}^d ; V)$.
\end{definition}

Take now a one dimensional white noise $\Xi$ with values in $V$ and let $B_t
(v) = \Xi (\mathbbm{1}_{[0, t]} \otimes v)$ for $v \in V$. For fixed $t
\geqslant 0$, $B_t$ extends as homomorphism to all $\Lambda V$ and therefore
$B_t \in \mathcal{G} (V)$. Note also that
\begin{equation}
  \| B_t (v) - B_s (v) \|_{\mathcal{A}} \leqslant | t - s |^{1 / 2} \| v \|_V,
  \qquad t, s \geqslant 0, v \in V, \label{eq:regularity-B}
\end{equation}
so $B (v) \in C (\mathbb{R}_+, \mathcal{A})$ and now note that $B \in
\tmop{Hom} (\Lambda V, C (\mathbb{R}_+, \mathcal{A}))$. This in particular
implies that $B_{t_1}, \ldots, B_{t_n}$ is a compatible family and that $B =
(B_t)_{t \in \mathbb{R}_+}$ is a Gaussian process with continuous
trajectories. We have that $B_0 (v) = 0$,
\[ \omega (B_t (v)) = 0, \qquad \omega (B_t (v) B_s (w)) = \langle v, C w
   \rangle  (t \wedge s), \qquad t, s \geqslant 0, \quad v, w \in V, \]
where $C$ is the correlation of the Grassmann Gaussian random variable $\Xi$.
Increments of $B$ are independent and higher order moments can be computed via
Wick's rule~{\eqref{eq:wick}}. Note in particular that
\begin{equation}
  \sup_{0 \leqslant t \leqslant T} \| B_t \|_{\mathcal{G} (V)} < \infty
  \label{eq:boundedness-B}
\end{equation}
for all $T > 0$. Let us remark here that properties~{\eqref{eq:regularity-B}}
and~{\eqref{eq:boundedness-B}} are very different from the pathwise properties
of the commutative (bosonic) Brownian motion.

\section{Finite dimensional SDEs}\label{sec:finite-dim}

We want to study simple SDEs taking values in $\Lambda V$ as far as it is
needed for the purpose of stochastic quantization, that is with additive white
noise. We refrain to undertake here a general study of Grassmann SDEs, in
particular no stochastic calculus will be needed below. It seems possible to
devise such a calculus but we leave it for a future work. More precisely in
this section we want to study SDEs driven by an additive Brownian motion $B_t$
when the linear space $V$ is finite dimensional.

\begin{definition}
  Let $F \in \mathcal{L} (V, \Lambda V)$ and assume that $F (v)$ is odd for
  all $v \in V$, and let $\Psi_0 \in \mathcal{G} (V)$ be a random variable
  compatible and independent of the Brownian motion $B \in \mathcal{G} (L^2
  (\mathbb{R}_+, V))$, then $\Psi \in C^0 ([0, T] ; \mathcal{G} (V))$ is a
  solution in $[0, T]$ to the (additive) SDE driven by $B$ with drift $F$ and
  initial condition $\Psi_0$, if,
  \begin{equation}
    \Psi_t (v) = \Psi_0 (v) + \int_0^t \Psi_s (F (v)) \mathd s + B_t (v),
    \qquad t \in [0, T], v \in V, \label{eq:sde}
  \end{equation}
  where the integral with respect to the variable $s$ is understood in
  Bochner's sense (relative to the norm on $\mathcal{G} (V)$ introduced in
  equation {\eqref{eq:distrance-g}}).
\end{definition}

\begin{remark}
  Note that we do not require that $\Psi_t, \Psi_s$ are compatible for $t \neq
  s$, in particular
  \[ C^0 ([0, T] ; \mathcal{G} (V)) \neq \tmop{Hom} (\Lambda V ; C ([0, T],
     \mathcal{A})), \]
  nor we require any compatibility of $\Psi$ and $B$. It will turn out that
  such compatibility holds in fact for the unique solution
  of~{\eqref{eq:sde}}, but is not necessary to put such restriction to
  formulate the notion of solution.
\end{remark}

\subsection{The Ornstein--Uhlenbeck motion}

In this section we introduce the Grassmann analog of the Ornstein--Uhlenbeck
process, that is when $F (v) = A v$ with $A : V \rightarrow V$ is a linear
operator on $V$. In this case we can write down an explicit formula for the
solution to equation~{\eqref{eq:sde}}.

\begin{proposition}
  \label{proposition_linear}Suppose that $A : V \rightarrow V$ and $F (v) = A
  v$. The unique solution to equation~{\eqref{eq:sde}} is given by
  \begin{equation}
    \Psi_t (v) = \Psi_0 (e^{A t} v) + \Xi (\mathbbm{1}_{[0, t]} (\cdot) e^{A (t
    - \cdot)} v), \qquad t \in \mathbb{R}_+, \label{eq:linear}
  \end{equation}
  where $\Xi$ is the Grassmann Gaussian noise related with the Brownian motion
  $B_t$ (see Definition~\ref{definition_Brownian} and the discussion that
  follows).
\end{proposition}

\begin{proof}
  Let $h (t, v) \in L^2 (\mathbb{R}; V)$ be given by
  \[ h (t, v) (s) \assign \mathbbm{1}_{[0, t]} (s) e^{A (t - s)} v, \qquad s
     \in \mathbb{R}. \]
  Then {\eqref{eq:linear}} reads $\Psi_t (v) = \Psi_0 (e^{A t} v) + \Xi (h (t,
  v))$. An explicit computation using that
  \[ h (t, v) (s) -\mathbbm{1}_{[0, t]} (s) v = \int_0^t A h (r, v) (s) \mathd
     r = \int_0^t h (r, A v) (s) \mathd r, \qquad s \in \mathbb{R}, \]
  in $L^2 (\mathbb{R}; V)$ gives
  \[ \Xi (h (t, v)) = B_t (v) + \int_0^t \Xi (h (r, A v)) \mathd r. \]
  Moreover, by the definition of exponential of a matrix, $\partial_t \Psi_0
  (e^{A t} v) = \Psi_0 (e^{A t} A v)$ so
  \[ \Psi_t (v) = \Psi_0 (v) + \int_0^t \Psi_0 (e^{A r} A v) \mathd r + B_t
     (v) + \int_0^t \Xi (h (r, A v)) \mathd r \]
  \[ = \Psi_0 (v) + B_t (v) + \int_0^t \Psi_r (A v) \mathd r \]
  as required by equation {\eqref{eq:linear}} for our assumption on $F$.
\end{proof}

\begin{remark}
  In particular this shows that if $(\Psi_0, \Xi)$ is a Grassmann Gaussian
  process then also the solution $\Psi$ to the SDE~{\eqref{eq:sde}}, is a
  Grassmann Gaussian process compatible with $(\Psi_0, \Xi)$. In order to
  prove this fact it is sufficient to prove that any product of the form
  \begin{equation}
    \omega (\Psi_{t_1} (v_1) \cdots \Psi_{t_{2 k}} (v_2)) \label{eq:product}
  \end{equation}
  can be computed using the products (weighted with a suitable sign due to the
  anti-commutation of $\Psi$) of the covariance $\omega (\Psi_{t_i} (v_i)
  \Psi_{t_j} (v_j))$ for any $i, j = 1, \ldots, 2 k$. This properties of the
  expectation~{\eqref{eq:product}} follows from the fact that
  \[ \begin{array}{c}
       \omega ((\Psi_{t_1} (v_1) - \Psi_0) \cdots (\Psi_{t_{2 k}} (v_n) -
       \Psi_0)) = \omega (\Xi (h (t_1, v_1)) \cdots \Xi (h (t_{2 k}, v_{2
       k}))) =\\
       = \sum_{\sigma} (- 1)^{\sigma} \prod_{i = 1}^k \int_0^{\min (t_i, t_j)}
       \langle e^{A (t_i - s)} v_{\sigma (2 i - 1)}, C e^{A (t_j - s)}
       v_{\sigma (2 i)} \rangle_V \mathd s
     \end{array} \]
  where $C$ is the covariance of $\Xi$. Using the fact that $\Psi_0$ is a
  Gaussian random variable independent of $\Xi$, and so $\Psi_0 (e^{A (t - s)}
  v)$ is also a Gaussian random variable independent of $\Xi$, we obtain the
  Gaussian behavior of $\Psi_t$. The compatibility of $\Psi_t$ with $(\Psi_0,
  \Xi)$ follows from the fact that $\Psi$ is a linear function of $(\Psi_0,
  \Xi) .$
\end{remark}

Let us now study the family of random variables
\[ \Psi_t^s (v) = \Xi (e^{A (t - \cdot)} v), \qquad t \in \mathbb{R}. \]
In the next proposition, we will show that this represents the stationary
solution to the SDE {\eqref{eq:linear}} (this motivates our putting the upper
index $s$ to $\Psi^s_t$).

\begin{proposition}
  \label{proposition_linearlimit}Assume that all eigenvalues of $A$ have
  strictly negative real part less or equal than $- \lambda_A$, where
  $\lambda_A > 0$. Then, we have $\sup_{t \in \mathbb{R}} \| \Psi^s_t
  \|_{\mathcal{G} (V)} < + \infty$, furthermore for any $G \in \Lambda V$ and
  any $t \in \mathbb{R}$ we get
  \begin{equation}
    \omega (G (\Psi^s_t)) = \omega (G (X^A)), \label{eq:G}
  \end{equation}
  where $X^A \in \mathcal{G} (V)$ is a Gaussian random variable with
  covariance $C_A$ given by
  \begin{equation}
    C_A \assign \int^{\infty}_0 e^{A^{\Tau} s} C e^{A s} \mathd s.
    \label{eq:ca}
  \end{equation}
  where $A^{\Tau}$ denote the transpose matrix.
\end{proposition}

\begin{remark}
  Expression {\eqref{eq:G}} in particular show that the independence of $t$ of
  the expectation $\omega (G (\Psi^s_t))$, expressing the stationarity of
  $\Psi^s_t$. 
\end{remark}

\begin{proof*}{Proof of Proposition \ref{proposition_linearlimit}}
  We have easily by the definition of $\Psi^s_t$:
  \[ \| \Psi^s_t \|_{\mathcal{G} (V)} \lesssim \| \mathbb{I}_{(- \infty, t]}
     (\cdot) e^{A (t - \cdot)} \|_{L^2 (\mathbb{R})} \lesssim
     \frac{1}{\lambda_A}, \]
  where $\lesssim$ stands for inequality modulo some appropriate positive
  constant. The random variable $\Psi^s_t$ is Gaussian so it is completely
  characterized (in term of moments) by its covariance. Note that the
  covariance of $\Psi_t^s$ can be easily computed as follows
  \[ \omega (\Psi_t^s (v) \Psi_t^s (w)) = \omega (\Xi (\mathbb{I}_{(- \infty,
     t]} (\cdot) e^{A (t - \cdot)} v) \Xi (\mathbb{I}_{(- \infty, t]} (\cdot)
     e^{A (t - \cdot)} w)) = \]
  \[ = \langle \mathbb{I}_{(- \infty, t]} (\cdot) e^{A (t - \cdot)} v,
     C\mathbb{I}_{(- \infty, t]} (\cdot) e^{A (t - \cdot)} w \rangle =
     \int^0_{- \infty} \langle e^{- A^{\Tau} \tau} v, C e^{- A \tau} w \rangle
     \mathd \tau = \langle v, C_A w \rangle \]
  with $C_A$ given by~{\eqref{eq:ca}}. The appearance of the transposition is
  due to the fact that $V$ is a \tmtextit{real} pre-Hilbert space.
\end{proof*}

Let us point out that
\begin{equation}
  A^{\Tau} C_A + C_A A = - C \label{eq:relation-ca}
\end{equation}
by a simple integration by parts applied to~{\eqref{eq:ca}}.

\begin{remark}
  Suppose that $A$ commutes with $C$, then we have that $\Psi_t^s$ is for all
  $t \in \mathbb{R}$ a Grassmann Gaussian process with correlation $C_A = - (A
  + A^{\Tau})^{- 1} C$ (the operator $A$ is invertible since all its
  eigenvalues has strictly negative real part).
\end{remark}

\subsection{Existence and uniqueness for general drift}

\begin{theorem}
  \label{theorem_picard}For any $\Psi_0 \in \mathcal{G} (V)$ compatible with
  $B$, there exists $T > 0$ such that equation~{\eqref{eq:sde}} admits a
  unique solution $\Psi \in C^0 ([0, T] ; \mathcal{G} (V))$. Moreover $\Psi$
  is compatible with $\Psi_0 \oplus B$.
\end{theorem}

\begin{proof}
  We are going to construct a solution via Picard's iteration. Let $\Psi^0_t =
  \Psi_0$ for all $t \geqslant 0$ and define the $n + 1$-th Picard's iteration
  by
  \[ \Psi_t^{n + 1} (v_{\alpha}) = \Psi_0 (v_{\alpha}) + \int_0^t \Psi_s^n (F
     (v_{\alpha})) \mathd s + B_t (v_{\alpha}), \qquad t \in [0, T], \]
  for any basis element $v_{\alpha} \in V$. Then $\Psi^0 \in C^0 ([0, T] ;
  \mathcal{G} (V))$ and since $\Psi_0$ is compatible with $B$ we have that,
  for all $t \geqslant 0$, $\Psi^1_t$ belong to the Grassmann algebra
  generated by $\Psi_0$ and $B$ in $\mathcal{A}$ and there $\Psi^1_t
  (v_{\alpha})$ is an odd element for all $\alpha$. Therefore $\Psi^1_t \in
  \mathcal{G} (V)$ for all $t \geqslant 0$ and also $\Psi^1_t$ is compatible
  with $\Psi_0 \oplus B$. Since it is true for all $t$, by approximation in
  the operator norm we see that $\Psi^1 \in C^0 ([0, T] ; \mathcal{G} (V))$
  and that it is compatible with $\Psi_0 \oplus B$. By induction we prove then
  that $\Psi^n \in C^0 ([0, T] ; \mathcal{G} (V))$ for all $n \geqslant 0$ and
  that it is compatible with $\Psi_0 \oplus B$.
  
  Now observe that $\Psi_s^n (F (v_{\alpha}))$ is a polynomial function of
  $(\Psi^{n, \alpha}_s)_{\alpha}$ and therefore by a standard contraction
  argument in the Banach space $\mathcal{A}$ we obtain that $(\Psi^n)_n$
  converges, as $n \rightarrow + \infty$, in $C^0 ([0, T] ; \mathcal{A})$ for
  some positive $T$ which depends only on $F$, on $\sup_{\alpha} \|
  \Psi_0^{\alpha} \|_{\mathcal{A}}$ and on $\sup_{0 \leqslant t \leqslant T,
  \alpha} \| B_t^{\alpha} \|_{\mathcal{A}}$. If we call $\Psi$ the limit we
  have that $\Psi \in C ([0, T] ; \mathcal{G} (V))$ by the first part of the
  proof and we deduce easily that $\Psi$ is compatible with $\Psi_0 \oplus B$.
\end{proof}

\begin{theorem}
  \label{th:global-sol}The solution to equation~{\eqref{eq:sde}} exists for
  all times.
\end{theorem}

\begin{proof}
  We will have global existence as soon as we can rule out explosions, that is
  prove that for all $t \geqslant 0$ we have $\sup_{\alpha} \| \Psi_t^{\alpha}
  \|_{\mathcal{A}} < \infty$ Let $\Theta_t (v) = \Psi_t (v) - B_t (v)$ and
  extend this map to a homeomorphism of $\Lambda V$ into $\mathcal{A}$.
  Observe that
  \[ \Theta_t (v) = \Theta_0 (v) + \int_0^t m_{\mathcal{A}} [(\Theta_s \otimes
     B_s) (\Delta F (v_{\alpha}))] \mathd s. \]
  In particular, if we consider a (linear) basis of $\Lambda V$ denoted by
  $(e_A)_A$ and let $\Theta^A_t \assign \Theta_t (e_A)$ and $B^A_t \assign B_t
  (e_A)$ we have
  \[ \Theta_t^A = \Theta_0^A + \int_0^t \sum_B h_{A, B, C} \Theta_s^B B_s^C
     \mathd s, \]
  for real coefficients $(h_{A, B, C})_{A, B, C}$. This is a
  finite-dimensional, linear system of non-autonomous ODEs in $\mathcal{L}
  (H)$ for $(\Theta_t^A)_A$ ($\mathcal{H}$ being the Hilbert space the element
  of $\mathcal{A}$ act, see Section \ref{sec:grassmann-rv}), and
  \[ \sum_A \| \Theta_t^A \|_{\mathcal{A}} \leqslant \sum_A \| \Theta_0^A
     \|_{\mathcal{A}} + | h | \int_0^t \sum_B \| \Theta_s^B \|_{\mathcal{A}}
     \sum_C \| B_s^C \|_{\mathcal{A}} \mathd s, \]
  with $| h | \assign \sup_{A, B, C} | h_{A, B, C} |$. Therefore by Gronwall
  inequality
  \[ \sum_A \| \Theta_t^A \|_{\mathcal{A}} \leqslant \sum_A \| \Theta_0^A
     \|_{\mathcal{A}} \exp \left( | h | \int_0^t \sum_C \| B_s^C
     \|_{\mathcal{A}} \mathd s \right) < \infty \]
  for all $t \geqslant 0$. In particular $\sup_{\alpha} \| \Psi_t (v_{\alpha})
  \|_{\mathcal{A}} \leqslant \sup_{\alpha} \| \Theta_t (v_{\alpha})
  \|_{\mathcal{A}} + \sup_{\alpha} \| B_t (v_{\alpha}) \|_{\mathcal{A}} <
  \infty$ for all $t \geqslant 0$ and the solution to
  equation~{\eqref{eq:sde}} exists for all times.
\end{proof}

\subsection{An It{\^o} formula for solutions of SDEs}

We want to prove a kind of It{\^o} formula for the solution to
equation~{\eqref{eq:sde}}. In order to do that let $Q_C \in \mathcal{L}
(\Lambda V \otimes V \otimes V, \Lambda V)$ be given by
\begin{equation}
  Q_C (f \otimes v \otimes w) = \langle v, C w \rangle f, \qquad f \in \Lambda
  V, v, w \in V. \label{eq:QC}
\end{equation}
For $F \in \mathcal{L} (V, \Lambda V)$ and $G \in \Lambda V \otimes V$ we
define $G \cdot F \in \Lambda V$ by extending linearly
\[ (f \otimes v) \cdot F = f F (v), \qquad f \in \Lambda V, v \in V, \]
where in both sides we have multiplication of elements in $\Lambda V$.

\begin{theorem}
  For the global solution to equation~{\eqref{eq:sde}} we have
  \begin{equation}
    \omega (\Psi_t (G)) = \omega (\Psi_0 (G)) + \int_0^t \omega \left( \Psi_s
    \left( \partial_R G \cdot F + \frac{1}{2} Q_C (\partial_R^2 G) \right)
    \right) \mathd s \label{eq:ito-3}
  \end{equation}
  for all $G \in \Lambda V$ and $t \geqslant 0$.
\end{theorem}

\begin{proof}
  Fix any $T \geqslant 0$. Note that
  \[ \| B_t - B_s \|_{\mathcal{G} (V)} \lesssim | t - s |^{1 / 2}, \]
  for all $0 \leqslant s < t \leqslant T$. From that it follows that, if $F$
  in {\eqref{eq:sde}} is a polynomial of degree $\deg (F)$:
  \begin{equation}
    \| (\Psi_t - \Psi_s) - (B_t - B_s) \|_{\mathcal{G} (V)} = \left\| \int_s^t
    F (\Psi_s) \mathd s \right\|_{\mathcal{G} (V)} \lesssim (1 + \| \Psi
    \|_{[0, T]})^{\deg (F)} | t - s |, \label{eq:holderreg}
  \end{equation}
  where $\| \Psi \|_{[0, T]} \assign \sup_{t \in [0, T]} \| \Psi_s
  \|_{\mathcal{G} (V)}$ which, by Theorem~\ref{th:global-sol} and the
  continuity of $\Psi$ with respect to the time $t$, is finite for every $T
  \in \mathbb{R}_+$. By Taylor formula (see Lemma \ref{lemma:taylor}) applied
  to polynomial $G$ of degree $\deg (G)$ we have
  \begin{equation}
    \begin{array}{c}
      \nobracket \| G (\Psi_r) - G (\Psi_s) - (\Psi_s \otimes (\Psi_r -
      \Psi_s)) \partial_R G + \nobracket \nobracket\\
      \left. \left. + \frac{1}{2} (\Psi_s \otimes (\Psi_r - \Psi_s) \otimes
      (\Psi_r - \Psi_s)) (\partial_R \partial_R G) \right\|
      \right._{\mathcal{A}} \lesssim (1 + \| \Psi \|_{[0, T]})^{\deg (F) (\deg
      (G) - 2)} | r - s |^{3 / 2}
    \end{array} \label{eq:holderreg2}
  \end{equation}
  since by {\eqref{eq:holderreg}} $\| \Psi_t - \Psi_s \|_{\mathcal{G} (V)}
  \lesssim (1 + \| \Psi \|_{[0, T]})^{\deg (F)} | t - s |^{1 / 2}$. In a
  similar way it is possible to obtain
  \begin{equation}
    \| \langle \partial_R G (\Psi_u), F (\Psi_u) \rangle - \langle \partial_R
    G (\Psi_s), F (\Psi_u) \rangle \|_{\mathcal{A}} \lesssim (1 + \| \Psi
    \|_{[0, T]})^{\deg (F) (\deg (G) - 2)} | u - s |^{1 / 2}
    \label{eq:holderreg3}
  \end{equation}
  On the other hand, by {\eqref{eq:holderreg}}, we have
  \begin{equation}
    \begin{array}{lll}
      \omega ((\Psi_s \otimes (\Psi_r - \Psi_s)) \partial_R G) & = & \omega
      \left( \int_s^r \langle \partial_R G (\Psi_s), F (\Psi_u) \rangle \mathd
      u \right)\\
      &  & + \omega ((\Psi_s \otimes (B_r - B_s)) (\partial_R G))\\
      & = & \omega \left( \int_s^r \langle \partial_R G (\Psi_u), F (\Psi_u)
      \rangle \mathd u \right)\\
      &  & + \omega (\langle \partial_R G (\Psi_s), (B_r - B_s) \rangle)\\
      &  & + O ((1 + \| \Psi \|_{[0, T]})^{\deg (F) (\deg (G) - 2)} | r - s
      |^{3 / 2})
    \end{array} \label{eq:holderreg4}
  \end{equation}
  In a similar way, using inequality {\eqref{eq:holderreg}} we obtain
  \begin{equation}
    \begin{array}{c}
      \omega ((\Psi_s \otimes (\Psi_r - \Psi_s) \otimes (\Psi_r - \Psi_s))
      (\partial_R^2 G)) = \omega ((\Psi_s \otimes (B_r - B_s) \otimes (B_r -
      B_s)) (\partial_R^2 G))\\
      + O ((1 + \| \Psi \|_{[0, T]})^{\deg (F) (\deg (G) - 2)} | r - s |^{3 /
      2}) .
    \end{array}
  \end{equation}
  Furthermore using the fact that $\partial_R G (\Psi_s)$ is independent of
  $(B_r - B_s)$ and $((B_r - B_s) \otimes (B_r - B_s))$ is independent of
  $\partial_R^2 G (\Psi_s)$ (this is due to the fact that \ $\Psi_s$ is a
  function of $\nobracket \{ \Psi_0, B_k | \nobracket \nobracket k \in [0, s]
  \}$ and $B_t$ is independent of $\Psi_0$ and has independent increments, we
  have:
  \[ \omega (\langle \partial_R G (\Psi_s), (B_r - B_s) \rangle) = \langle
     \omega ((B_r - B_s)), \omega (\partial_R G (\Psi_s)) \rangle = 0 \]
  (where on the r.h.s. we understand that averages w.r.t. $\omega$ are taken
  componentwise) and
  \[ \omega [\langle \partial_R \partial_R G (\Psi_s), (B_r - B_s) \otimes
     (B_r - B_s) \rangle_{V \otimes V}] \]
  \[ = \langle \omega [(\partial_R \partial_R G (\Psi_s))], \omega [((B_r -
     B_s) \otimes (B_r - B_s))] \rangle \]
  \[ = (r - s) \omega [Q_C (\partial_R^2 G (\Psi_s))] . \]
  By taking a partition $\{ t_i \}_{i \in \{ 0, \ldots, n \}}$, of diameter
  $\rho_n \rightarrow 0$ as $n \rightarrow + \infty$, of $[0, t]$ and
  exploiting inequality {\eqref{eq:holderreg2}} we obtain
  \[ \begin{array}{lll}
       \omega (G (\Psi_t)) & = & \omega (G (\Psi_0)) + \sum_{i = 0}^{n + 1}
       \omega (G (\Psi_{t_i})) - \omega (G (\Psi_{t_{i - 1}}))\\
       & = & \omega (G (\Psi_0)) + \sum_{i = 1}^{n + 1} \omega ((\Psi_{t_{i -
       1}} \otimes (\Psi_{t_i} - \Psi_{t_{i - 1}})) \partial_R G)\\
       &  & + \frac{1}{2} \omega ((\Psi_s \otimes (\Psi_r - \Psi_s) \otimes
       (\Psi_r - \Psi_s)) (\partial_R^2 G)) + O (\rho_n^{1 / 2})\\
       & = & \omega (G (\Psi_0)) + \omega \left( \int_0^t \partial_R G
       (\Psi_u) F (\Psi_u) \mathd u \right)\\
       &  & + \frac{1}{2} \sum_{i = 1}^{n + 1} (t_i - t_{i - 1}) \omega [Q_C
       (\partial_R^2 G (\Psi_{t_i}))] + O (\rho_n^{1 / 2}),
     \end{array} \]
  where the constants in $O$ are proportional to $(1 + \| \Psi \|_{[0,
  T]})^{\deg (F) \deg (G)}$ and do not depend on the partition. Taking the
  limit $n \rightarrow + \infty$, and so $\rho_n \rightarrow 0$, and using the
  fact that $\omega [Q_C (\partial_R^2 G (\Psi_s))]$ is continuous in $s$ ($G
  (\Psi_s)$ being a continuous function from $\mathbb{R}_+$ to $\mathcal{A}$
  since $\Psi \in C^0 ([0, T], \mathcal{G} (V))$) we obtain the thesis.
\end{proof}

\subsection{Invariant measures}

Consider an equation of the following kind
\begin{equation}
  \Psi_t (v) = \Psi_0 (v) + \int_0^t (\Psi_s (A v) + \Psi_s (F (v))) \mathd s
  + B_t (v), \qquad t \geqslant 0, v \in V, \label{eq:sde-langevin-0}
\end{equation}
for $\Psi \in C^0 ([0, T] ; \tmop{Hom} (V, \mathcal{A}))$, $F \in \tmop{Hom}
(V, \Lambda V)$, $A \in \mathcal{L} (V)$, $B$ a Grassmann Brownian motion with
correlation $C$ and with initial condition $\Psi_0$ distributed as a Grassmann
Gaussian independent of $B$ with correlation $C_A$ (compare to
equation~{\eqref{eq:ca}}). Fix $U \in \Lambda V$ even, and consider the linear
functional $\tilde{\omega} : \mathcal{A} \rightarrow \mathbb{C}$ given by
\begin{equation}
  \tilde{\omega} (\cdot) \assign \omega (\cdot e^{- 2 U (\Psi_0)}),
  \label{eq:gibbs-U}
\end{equation}
where $\omega$ is, as before, is the positive state on the $\asterisk$-algebra
$\mathcal{A} = \mathcal{L} (\mathcal{H})$.

\begin{definition}
  \label{definition_invariant}We say that the functional $\tilde{\omega}$ is
  an invariant measure for the equation~{\eqref{eq:sde-langevin-0}} if for any
  $G \in \Lambda V$ and any $t \in \mathbb{R}_+$ we have
  \begin{equation}
    \tilde{\omega} (G (\Psi_t)) = \tilde{\omega} (G (\Psi_0)) .
    \label{eq:constant}
  \end{equation}
\end{definition}

\begin{remark}
  The functional is invariant in the sense of
  Definition~\ref{definition_invariant}, only if the expectations of the
  form~{\eqref{eq:constant}} are constant when expressions of the form $G
  (\Psi_t)$ are evaluated. This means in particular that only polynomials in
  $\Psi_t$ are considered.
\end{remark}

There is a simple condition, analogous to the Fokker--Planck equation for
commutative SDEs, for checking the invariance of the functional
$\tilde{\omega}$. Here it is:

\begin{lemma}
  \label{lemma_invariant}If $\Psi$ is a solution
  to~{\eqref{eq:sde-langevin-0}} and
  \begin{equation}
    \tilde{\omega} (\mathcal{L}G (\Psi_0)) = 0, \qquad G \in \Lambda V
    \label{eq:inf-invariance},
  \end{equation}
  for
  \[ \mathcal{L}G \assign \partial_R G \cdot F + \frac{1}{2} Q_C (\partial_R^2
     G), \]
  then $\tilde{\omega}$ is an invariant measure in the sense of
  Definition~\ref{definition_invariant}.
\end{lemma}

\begin{proof}
  Consider $G = G_A = v_A$ where $(v_A)_A$ is a (finite) linear basis of
  $\Lambda V$. Then we have $\mathcal{L} (G_A) = \sum_B \kappa_A^B v_B$ for a
  suitable family $(\kappa_A^B)_{A, B}$ of constants in $\mathbb{R}$. Write
  $\mathfrak{P}_A (t) \assign \tilde{\omega} (G_A (\Psi_t)) = \omega (G_A
  (\Psi_t) e^{- 2 U (\Psi_0)})$, due to definition {\eqref{eq:gibbs-U}}. By
  It{\^o} formula~{\eqref{eq:ito-3}} we have that $\mathfrak{P}_A$ is the
  unique solution of the system of ODEs
  \begin{eqnarray}
    \partial_t \mathfrak{P}_A (t) & = & \partial_t \omega (G_A (\Psi_t) e^{- 2
    U (\Psi_0)}) = \omega ((\mathcal{L} G_A) (\Psi_t) e^{- 2 U (\Psi_0)})
    \nonumber\\
    & = & \sum_B \kappa_A^B \omega (G_B (\Psi_t) e^{- 2 U (\Psi_0)}) = \sum_B
    \kappa^B_A \mathfrak{P}_B (t), \label{eq:ODE1} 
  \end{eqnarray}
  with initial condition $\mathfrak{P}_A (0) = \omega (G_A (\Psi_0) e^{- 2 U
  (\Psi_0)})$. On the other hand condition~{\eqref{eq:inf-invariance}} implies
  that
  \[ \sum_B \kappa^B_A \mathfrak{P}_B (0) = \sum_B \kappa^B_A \omega (G_B
     (\Psi_0) e^{- 2 U (\Psi_0)}) = \omega (\mathcal{L}G (\Psi_0) e^{- 2 U
     (\Psi_0)}) = 0 \]
  which means that $(\mathfrak{P}_A (t))_A$ is constantly zero by uniqueness
  of solution to the ODEs system~{\eqref{eq:ODE1}}. Since the expectations of
  the form~{\eqref{eq:constant}} are linear combinations of $(\mathfrak{P}_A
  (t))_A$ the proof is complete.
\end{proof}

The following theorem gives a sufficient condition on the
SDE~{\eqref{eq:sde-langevin-0}} to have the invariant
measure~{\eqref{eq:gibbs-U}}.

\begin{theorem}
  \label{th:sq-finite}For any even $U \in \Lambda_{\tmop{even}} V$, the SDE
  \begin{equation}
    \Psi_t (v) = \Psi_0 (v) + \int_0^t (\Psi_s (A v) + \Psi_s (\langle C
    \partial_R U, v \rangle) \nobracket \mathd s + B_t (v) \qquad t \geqslant
    0, v \in V, \label{eq:sde-langevin}
  \end{equation}
  with $B$ a Brownian motion with correlation $C$ and $\Psi_0$ an independent
  Gaussian initial condition with correlation $C_A$ (defined in equation
  {\eqref{eq:ca}}) has $\tilde{\omega}$ as invariant measure provided
  \begin{equation}
    A^{\Tau} C_A - C_A A = 0. \label{eq:gradient}
  \end{equation}
\end{theorem}

\begin{proof}
  The proof is a consequence of Lemma~\ref{lemma_invariant}. Denote
  $\omega^{\Psi_0} (G) \assign \omega (G (\Psi_0))$ the law of $\Psi_0$, and
  observe that for $H \in \Lambda V \otimes V$ we have the integration by
  parts formula
  \[ \omega ((\Psi_0 \otimes \Psi_0) (H)) = \omega (\Psi_0 (Q_{C_A}
     (\partial_R H))), \]
  where $Q_C$ is defined in {\eqref{eq:QC}}. Then for all $G \in \Lambda V$
  and $U \in \Lambda_{\tmop{even}} V$ we have
  \[ \omega ((\Psi_0 \otimes \Psi_0) (A (\partial_R G) e^{- 2 U (\Psi_0)})) =
     \omega^{\Psi_0} (Q_{C_A} (\partial_R (A \partial_R G e^{- 2 U
     (\Psi_0)}))) \]
  \[ = \omega^{\Psi_0} (Q_{C_A} ((A \otimes 1) \partial_R \partial_R G) e^{- 2
     U (\Psi_0)}) - 2 \omega^{\Psi_0} (Q_{C_A} ((A \partial_R G) \otimes
     \partial_R U) e^{- 2 U (\Psi_0)}) . \]
  Observe that using the equality~{\eqref{eq:relation-ca}} we have
  \[ Q_{C_A} ((A \otimes 1) \partial_R \partial_R G) = \tmop{Tr} ((A \otimes
     C_A) \partial_R \partial_R G) \]
  \[ = \frac{1}{2} \tmop{Tr} ((C_A A + A^{\Tau} C_A) \partial_R \partial_R G)
     = - \frac{1}{2} Q_C (\partial_R \partial_R G), \]
  where we have used the trace $\tmop{Tr}$ on $V$ given by $\tmop{Tr} (v
  \otimes w) = \langle v, w \rangle$. Moreover, from~{\eqref{eq:gradient}} we
  have also
  \[ Q_{C_A} ((A \partial_R G) \otimes \partial_R U) = \tmop{Tr} ((A
     \partial_R G) \otimes (C_A \partial_R U)) = \tmop{Tr} (\partial_R G
     \otimes (A^{\Tau} C_A \partial_R U)) \]
  \[ = \tmop{Tr} \left( \partial_R G \otimes \left( \frac{A^{\Tau} C_A + C_A
     A}{2} \partial_R U \right) \right) + \tmop{Tr} \left( \partial_R G
     \otimes \left( \frac{A^{\Tau} C_A - C_A A}{2} \partial_R U \right)
     \right) \]
  \[ = - \frac{1}{2} \tmop{Tr} (\partial_R G \otimes (C \partial_R U)) = -
     \frac{1}{2} \langle \partial_R G, C \partial_R U \rangle . \]
  since we assumed that $A^{\Tau} C_A - C_A A = 0$ and
  used~{\eqref{eq:relation-ca}} again. Therefore
  \[ \omega ((\Psi_0 \otimes \Psi_0) (A (\partial_R G) e^{- 2 U (\Psi_0)})) =
     - \frac{1}{2} \omega^{\Psi_0} (Q_C (\partial_R^2 G) e^{- 2 U (\Psi_0)}) +
     \omega^{\Psi_0} (\langle \partial_R G, C \partial_R U \rangle e^{- 2 U
     (\Psi_0)}) . \]
  Now observing that
  \[ \mathcal{L}G (\Psi_0) = (\Psi_0 \otimes \Psi_0) (A \partial_R G) + \Psi_0
     \left( \langle \partial_R G, C \partial_R U \rangle + \frac{1}{2} Q_C
     (\partial_R \partial_R G) \right) \]
  we can deduce
  \[ \omega (\mathcal{L}G (\Psi_0) e^{- 2 U (\Psi_0)}) = \omega ((\Psi_0
     \otimes \Psi_0) (A \partial_R G) e^{- 2 U (\Psi_0)}) + \]
  \[ + \omega^{\Psi_0} \left( \left( \langle \partial_R G, C \partial_R U
     \rangle + \frac{1}{2} Q_C (\partial_R \partial_R G)  \right) e^{- 2 U
     (\Psi_0)} \right) = 0. \]
  
\end{proof}

\subsection{Long-time behavior for small non-linearity}

We will now investigate the existence of solutions which are globally bounded
in time. In the commutative setting, non-linear equations can have globally
bounded solutions only if the non-linearity stays uniformly small or if it
shows some coercivity. Like all notions of positivity, also coercivity however
does not apply well in the Grassmann setting. The only kind of coercive term
we have identified is a linear drift with a negative sign. This is a very mild
coercive term, but it turns out to be enough provided the non-linearity is
small enough.

\

Consider the equation
\begin{equation}
  \Psi_t (v) = \Psi_0 (v) + \int_0^t (\Psi_s (A v) + \lambda \Psi_s (F (v)))
  \mathd s + B_t (v), \qquad t \geqslant 0, v \in V, \label{eq:sde-lambda}
\end{equation}
for $\Psi \in C ([0, T] ; \tmop{Hom} (V, \mathcal{A}))$, with $\lambda
\geqslant 0$ and $F \in \tmop{Hom} (V, \Lambda V)$ where $A$ is an operator
with eigenvalues having strictly negative real part less then $- \lambda_A <
0$. In this setting we introduce the notion of stationary solution to
equation~{\eqref{eq:sde-lambda}}.

\begin{definition}
  We say that $\Psi^s_t \in C^0 (\mathbb{R}, \mathcal{G} (V))$ is a stationary
  solution to equation~{\eqref{eq:sde-lambda}} of norm at most $K \in
  \mathbb{R}_+$ if the following two conditions hold:
  \begin{enumerate}
    \item $\sup_{t \in \mathbb{R}} \| \Psi^s_t \| < K$,
    
    \item $(\Psi^s_t)_{t \in \mathbb{R}}$ solves the following integral
    equation
    \begin{equation}
      \Psi^s_t = \lambda \int_{- \infty}^t \Psi^s_{\tau} (F (e^{A \tau} v))
      \mathd \tau + B^A_t \label{eq:stationary} \qquad t \in \mathbb{R}
    \end{equation}
    where \ $B^A_t (v) = \Xi (e^{A (t - \cdot)} (v))$.
  \end{enumerate}
\end{definition}

Hereafter we write
\[ L (A, C) = \sup_{t \in \mathbb{R}} \| B^A_t \|_{\mathcal{G} (V)}, \]
and we denote by $\mathfrak{K}_K \subset C^0 (\mathbb{R}, \mathcal{G} (V))$
the set of $\Psi \in C^0 (\mathbb{R}, \mathcal{G} (V))$ such that
\[ \sup_{t \in \mathbb{R}} \| \Psi_t \|_{\mathcal{G} (V)} \leqslant K_{}, \]
for some constant $K > 0$.

\begin{theorem}
  \label{theorem_stationary}Assume $K \geqslant 3 L (A, C)$ and suppose that
  all the eigenvalues of $A$ have negative real part that are less or equal
  than $- \lambda_A$, then there exists $\lambda_0 (K, L, F)$ depending on
  $K$, $L (A, C)$ and $F$ such that if $| \lambda | \leqslant 2 \lambda_0 (K,
  L, F)$ there exists a unique stationary solution of norm at most $K$ to
  equation~{\eqref{eq:sde-lambda}}.
\end{theorem}

\begin{proof}
  It is a simple application of the Banach fix point theorem by noting that
  the map
  \[ \mathcal{K}_{\lambda} (\Psi)_t \assign \lambda \int_{- \infty}^t
     \Psi_{\tau} (F (e^{A (t - \tau)} v)) \mathd \tau + B^A_t, \qquad t \in
     \mathbb{R}, \]
  is a contraction in $\mathfrak{K}_K$. Indeed, if $\Psi \in \mathfrak{K}_K$
  we have
  \[ \| \mathcal{K}_{\lambda} (\Psi)_t \|_{\mathcal{G} (V)} \leqslant |
     \lambda |  \| F \| (1 + K)^{\deg (F)} \int_{- \infty}^t e^{- \lambda_A (t
     - \tau)} \mathd \tau + L (A, C), \]
  where we use the notation $\| F \|$ for the smallest constant such that
  \[ \sup_{v : \| v \| \leqslant 1}  \| \psi (F (v)) \|_{\mathcal{A}}
     \leqslant \| F \| \deg (F) (1 + \| \psi \|_{\mathcal{G} (V)})^{\deg (F)}
     \qquad \psi \in \mathcal{G} (V) . \]
  If we request that
  \begin{equation}
    2 \lambda_0 \frac{\| F \| (1 + K)^{\deg (F)}}{\lambda_A} + 2 L < K,
    \label{eq:contraction2}
  \end{equation}
  (where $\lambda_0 \assign \lambda_0 (K, L, F)$) for $| \lambda | \leqslant 2
  \lambda_0$ small enough we have that $\mathcal{K}_{\lambda}$ maps
  $\mathfrak{K}_K$ into itself. Furthermore
  \[ \| \mathcal{K}_{\lambda} (\Psi_1)_t - \mathcal{K}_{\lambda} (\Psi_2)_t
     \|_{\mathcal{G} (V)} \leqslant \frac{| \lambda |}{\lambda_A} \|
     \partial_R F \| (1 + K)^{\deg (F) - 1} \| \Psi_{1, t} - \Psi_{2, t}
     \|_{\mathcal{G} (V)} . \]
  Therefore, if we take
  \begin{equation}
    \frac{2 \lambda_0}{\lambda_A} \| \partial_R F \| (1 + K)^{\deg (F) - 1} <
    1 \label{eq:contraction}
  \end{equation}
  the map $\mathcal{K}_{\lambda}$ is a strict contraction. The fact that
  $\mathfrak{K}_K \neq \emptyset$ will be a consequence of the Lemma
  (\ref{lemma_bound}) below.
\end{proof}

\begin{remark}
  \label{remark_lambda0}The function $\lambda_0$ in
  Theorem~\ref{theorem_stationary} can be taken to be a decreasing function of
  the parameters $L$ and $K$.
\end{remark}

We introduce an approximation for the solution to
equation~{\eqref{eq:stationary}}. Let $X \in \mathcal{G} (V)$ independent of
$B^A_t$ and consider the element $\Psi_{- T, t}^X \in C^0 (\mathbb{R},
\mathcal{G} (V))$ which is $\Psi_{- T, t}^X = X$ if $t \leqslant - T$, for
some $T > 0$, and is a solution to the SDE
\[ \Psi_{- T, t}^X = X (e^{A (t + T)} v) + \lambda \int_{- T}^t F (\Psi_{- T,
   t}^X (e^{A (t - \tau)} v)) + B^A_t - B^A_{- T} (e^{A (t + T)} v), \]
when $t > - T$.

\begin{lemma}
  \label{lemma_bound}Under the assumptions of
  Theorem~\ref{theorem_stationary}, take $\| X \|_{\mathcal{G} (V)} <
  \frac{1}{6} K$, then, for any $| \lambda | \leqslant \frac{1}{2} \lambda_0
  (K, L, F)$, we have that $\Psi_{- T, t}^X \in \mathfrak{K}_K$.
\end{lemma}

\begin{proof}
  Let $\tau \in \mathbb{R}$ then we have the bound
  \begin{eqnarray}
    \sup_{t \in (- \infty, \tau]} \| \Psi^X_{- T, t} \|_{\mathcal{G} (V)} &
    \leqslant & e^{- \lambda_A \tau} (\| X \|_{\mathcal{G} (V)} + \| B_{- T}^A
    \|_{\mathcal{G} (V)}) + \nonumber\\
    &  & + \| F \| \frac{| \lambda |}{\lambda_A} (1 + \sup_{t \in (- \infty,
    \tau]} \| \Psi_{- T, t}^X \|_{\mathcal{G} (V)})^{\deg (F)} + L (A, C) .
    \nonumber
  \end{eqnarray}
  On the other hand by definition of $\lambda_0$ we have, for any $\tau > 0$
  \[ e^{- \lambda_A \tau} \left( \frac{1}{3} K + \frac{1}{6} K \right) + \| F
     \| \frac{| \lambda |}{\lambda_A} (1 + K)^{\deg (F)} + L (A, C) < K. \]
  Let $R_{\tau} \assign \sup_{t \in (- \infty, \tau]} \| \Psi_{- T, t}^X
  \|_{\mathcal{G} (V)}$ and consider the set $Z = \left\{ \tau \in \mathbb{R}:
  R_{\tau} \leqslant \frac{2 K}{3} \right\}$. The set $Z$ is non-empty, since
  $R_{\tau} < \frac{1}{2} K$ when $\tau \leqslant - T$, open since $\tau
  \mapsto R_{\tau}$ is an increasing continuous function of $\tau$ and closed
  since $\tau \rightarrow R_{\tau}$ is continuous. Then we must have $Z
  =\mathbb{R}$ and
  \[ \sup_{t \in \mathbb{R}} \| \Psi_{- T, t}^X \|_{\mathcal{G} (V)}
     \leqslant \frac{2}{3} K, \]
  for $| \lambda | \leqslant \lambda_0$.
\end{proof}

\begin{theorem}
  \label{theorem_approximation}Under the hypotheses of Lemma~\ref{lemma_bound}
  we have that $\Psi_{- T, t}^X$ converges to $\Psi^s_t$ in $\mathcal{G} (V)$
  as $T \rightarrow + \infty$, uniformly on compact subsets of $\mathbb{R}$.
\end{theorem}

\begin{proof}
  If $t \geqslant - T$ we have that
  \begin{eqnarray}
    \| \Psi_{- T, t}^X - \Psi_t^s \|_{\mathcal{G} (V)} & \lesssim & e^{-
    \lambda_A (t + T)} (\| X \|_{\mathcal{G} (V)} + \| B_{- T}^A
    \|_{\mathcal{G} (V)}) + | \lambda | \| F \| (1 + K)^{\deg (F)} \times
    \nonumber\\
    &  & \times \int_{- \infty}^{- T} e^{- \lambda_A (t - \tau)} \mathd \tau
    + | \lambda | \| \partial_R F \| (1 + K)^{\deg (F) - 1} \times \nonumber\\
    &  & \times \int_{- T}^t e^{- \lambda_A (t - s)} \| \Psi_{- T, \tau}^X -
    \Psi_{\tau}^s \| \mathd \tau . \nonumber
  \end{eqnarray}
  By Gronwall lemma we obtain
  \begin{eqnarray}
    \| \Psi_{- T, t}^X - \Psi_t^s \|_{\mathcal{G} (V)} & \leqslant & e^{-
    \lambda_A T} \left( \| X \|_{\mathcal{G} (V)} + \frac{K}{3} + \frac{|
    \lambda | \| F \| (1 + K)^{\deg (F)}}{\lambda_A} \right) \times
    \nonumber\\
    &  & \times (1 + \exp (- (\lambda_A - | \lambda | (\| \partial_R F \|
    K)^{\deg (F) - 1}) \nobracket (t + T)) . \nonumber
  \end{eqnarray}
  Since $\lambda_A - | \lambda | (\| \partial_R F \| K)^{\deg (F) - 1} > 0$,
  by the condition~{\eqref{eq:contraction}}, we have that
  \[ \| \Psi_{- T, t}^X - \Psi_t^s \|_{\mathcal{G} (V)} \lesssim e^{-
     \lambda_A T} \]
  uniformly in $T$ when $t > P > - T$ for any such fixed $P$. This proves that
  $\Psi_{- T, t}^X$ converges to $\Psi_t^s$ uniformly on compact sets.
\end{proof}

An important consequence of Theorem~\ref{theorem_approximation} is the
following one.

\begin{theorem}
  \label{theorem_invariant2}Under the hypotheses and notations of
  Theorem~\ref{theorem_stationary} and Theorem~\ref{theorem_approximation}, if
  $\| X \|_{\mathcal{G} (V)} < \frac{1}{3} K$, $| \lambda | \leqslant
  \lambda_0 (K, L, F)$ and if $\omega (\cdot e^{- 2 U (X)})$ is an invariant
  measure for equation~{\eqref{eq:sde-lambda}} (in the sense of
  Definition~\ref{definition_invariant}), for any $t \in \mathbb{R}$ and for
  any $G \in \Lambda V$, we have that
  \[ \frac{\omega (G (X) e^{- 2 U (X)})}{\omega (e^{- 2 U (X)})} = \omega (G
     (\Psi^s_t)) . \]
\end{theorem}

\begin{proof}
  Under the previous notation by Definition~\ref{definition_invariant}, we
  have that, for any $t \geqslant - T$,
  \[ \omega (G (X) e^{- 2 U (X)}) = \omega (G (\Psi_{- T, - T}^X) e^{- 2 U
     (X)}) = \omega (G (\Psi_{- T, t}^X) e^{- 2 U (X)}) . \]
  By Theorem~\ref{theorem_approximation} we can take the limit at the right
  hand side as $T \rightarrow + \infty$, obtaining
  \[ \omega (G (X) e^{- 2 U (X)}) = \omega (G (\Psi^s_t) e^{- 2 U (X)}) . \]
  Since $X$ is independent of $B^A_t$ and since $\Psi^s_t$ is a function only
  of $B^A$ we get
  \[ \omega (G (X) e^{- 2 U (X)}) = \omega (G (\Psi^s_t)) \omega (e^{- 2 U
     (X)}), \]
  from which the thesis follows.
\end{proof}

\begin{remark}
  \label{remark_invariant2}By Theorem~\ref{th:sq-finite},
  Theorem~\ref{theorem_invariant2} applies to the case $\lambda F (v) =
  \lambda \langle C \partial_R U, v \rangle$ for $\lambda$ small enough.
\end{remark}

\section{Infinite dimensional SDEs}\label{sec:infinite-dim}

We want to study SDEs of the form
\begin{equation}
  \mathd \Psi_t = (A \Psi_t + \lambda F (\Psi_t)) \mathd t + \mathd B_t
  \label{eq:infinite5}
\end{equation}
in the case where $\Psi$ takes values in $\mathcal{G} (V)$ and $V =\mathcal{S}
(\mathbb{R}^d) \otimes \mathbb{R}^n,$ or $\mathcal{S} (\mathbb{R}^d) \otimes
\mathbb{C}^n$, or $C^{\infty} (\mathbb{T}^d) \otimes \mathbb{R}^n$ for some $n
\in \mathbb{N}_0$ and with $A$ a (possibly unbounded) linear operator on $V$.
We discuss in detail the case $V = \mathcal{S} (\mathbb{R}^d)$, but every
reasoning can be easily generalized to the other cases or more generally to
any nuclear space of smooth functions.

\

The main difference between the finite and infinite dimensional case for
Grassmann SDEs is that in general there is no natural topology on $\mathcal{G}
(V)$. The space $V = \mathcal{S} (\mathbb{R}^d)$ is a Fr{\'e}chet space. Since
not every element of $\mathcal{G} (V)$ is continuous with respect to this
topology we consider here the space $\mathcal{G}^{- \infty} (V) \subset
\mathcal{G} (V)$, the subset of $\mathcal{G} (V)$ for which $X \in
\mathcal{G}^{- \infty} (V)$ if and only if $X$ restricted to $V \subset
\Lambda V$ is a continuous linear map from $V$ to $\mathcal{A}$. Hereafter we
always consider $X \in \mathcal{G}^{- \infty} (V)$ both as a homomorphism
from the algebra $\Lambda V$ into $\mathcal{A}$, and as a continuous linear
map from $V$ into $\mathcal{A}$.

\

Since $V$ is a nuclear space whose dual is nuclear the set of continuous
linear maps from $V$ into $\mathcal{A}$ can be identified with the tensor
product $V^{\asterisk} \otimes \mathcal{A} = \mathcal{S}' (\mathbb{R}^n)
\otimes \mathcal{A}$ (where $V^{\asterisk}$ the topological dual of $V$, the
Schwartz space of tempered distribution $\mathcal{S}' (\mathbb{R}^d)$). In
other words we can look upon $\mathcal{G}^{- \infty} (V)$ as the subset of
elements of $\mathcal{S}' (V) \otimes \mathcal{A}$ that can be extended to a
homomorphism from $\Lambda V$ into $\mathcal{A}$.

\

Using a reasoning similar to the one exploited in
Lemma~\ref{lemma_completeness}, it is possible to prove that $\mathcal{G}^{-
\infty} (V)$ is a closed subset of $V^{\asterisk} \otimes \mathcal{A}$ with
respect to the strong topology (of both $V^{\asterisk} = \mathcal{S}'
(\mathbb{R}^d)$ and $\mathcal{A}$), i.e. if $(X_{\nu})_{\nu}$ is a net in
$\mathcal{G}^{- \infty} (V)$ converging to $X$ in the topology of
$V^{\asterisk} \otimes \mathcal{A}$, then $X$ can be extended to an
homomorphism from $\Lambda V$ into $\mathcal{A}$. Working on $\mathcal{G}^{-
\infty} (V)$ is not so easy since $\mathcal{S}' (V) \otimes \mathcal{A}$ is
not even a Fr{\'e}chet space. For this reason we introduce now some subsets of
$\mathcal{G}^{- \infty} (V)$ (which are analogous to Besov spaces) which \
admit a stronger metric topology.

\

We recall here only the definition of the Besov norm for functions taking
values in a Banach space, further references and details are given in Appendix
\ref{appendix_Besov} (see also {\cite{Amann1997,Amann2019}}). Let $K_i \in
\mathcal{S} (\mathbb{R}^d)$, with $i \geqslant - 1$, be the functions
corresponding to the Littlewood--Paley blocks used in the definition of Besov
spaces $B^s_{p, q} \assign B^s_{p, q} (\mathbb{R}^d, \mathbb{R})$, where $s
\in \mathbb{R}$ and $p, q \in [1, + \infty]$. If $X \in \mathcal{G}^{- \infty}
(V)$ we define
\[ \Delta_i X (x) = X (K_i (x - \cdot)), \qquad x \in \mathbb{R}^d, i
   \geqslant - 1. \]
The function $\Delta_i X : \mathbb{R}^d \rightarrow \mathcal{A}$ can be
identified with an element $\Delta_i X \in \mathcal{S} (\mathbb{R}^d) \otimes
\mathcal{A} = \mathcal{S} (\mathbb{R}^d, \mathcal{A})$ and we define
\[ \| X \|_{B^s_{p, q} (\mathbb{R}^d, \mathcal{A})} \assign \left[ \sum_{i
   \geqslant - 1} 2^{q i s} \left( \int_{\mathbb{R}^d} \| \Delta_i X (x)
   \|_{\mathcal{A}}^p \mathd x \right)^{q / p}_{} \right]^{1 / q} . \]
Given $s \in \mathbb{R}$, we say that $X \in \mathcal{G}^s (V)$ when $X \in
\mathcal{G}^{- \infty} (V)$ and $\| X \|_{B^s_{\infty, \infty} (\mathbb{R}^d,
\mathcal{A})} < + \infty$. Thus we introduce the natural metric
$d_{\mathcal{G}^s (V)} (X, Y) = \| X - Y \|_{B^s_{\infty, \infty}
(\mathbb{R}^d, \mathcal{A})}$. Hereafter we always use the notation
$\mathcal{C}^s (\mathbb{R}^d, \mathcal{A}) = B^s_{\infty, \infty}
(\mathbb{R}^d, \mathcal{A})$ and $\mathcal{C}^s = \mathcal{C}^s (\mathbb{R}^d,
\mathbb{R})$. Moreover $\mathcal{C}^{s \asterisk}$ will denote the topological
dual of $\mathcal{C}^s$.

\begin{lemma}
  The set $\mathcal{G}^s (V)$ with the metric $d_{\mathcal{G}^s (V)}$ is a
  complete metric space.
\end{lemma}

\begin{proof}
  The proof is analogous to that of Lemma~\ref{lemma_completeness}.
\end{proof}

The case $s > 0$ can be understood more easily. In this case the elements $X
\in \mathcal{G}^s (V)$ can be identified with some continuous function
$\dot{X} : \mathbb{R}^d \rightarrow \mathcal{A}$. Under this identification if
$f \in \mathcal{S} (\mathbb{R}^d)$ we have that
\[ X (f) = \int_{\mathbb{R}^d} f (x) \dot{X} (x) \mathd x, \]
where the integral is taken in Bochner sense. If $X \in \mathcal{G}^s (V)$,
with $s > 0$, we can obtain the function $\dot{X}$ in the following way. Let
$\mathfrak{a} : \mathbb{R}^d \rightarrow \mathbb{R}$ be a smooth mollifier and
define $\mathfrak{a}_{x, \varepsilon} (y) = \varepsilon^{- d} \mathfrak{a}
\left( \frac{y - x}{\varepsilon} \right)$, where $x, y \in \mathbb{R}^d$ and
$\varepsilon > 0$. Then we have that
\[ \dot{X} (x) = \lim_{\varepsilon \rightarrow 0} X
   (\mathfrak{a}_{\varepsilon, x}) = X (\delta_x), \]
where the limit is taken with respect to the strong topology of $\mathcal{A}$,
and where $\delta_x$ is the Dirac delta with unit mass at the point $x \in
\mathbb{R}^d$. More generally, if $X \in \mathcal{G}^s (V)$ and $g \in
\mathcal{C}^{s \asterisk}$, for any sequence $(g_n)_n \subset \mathcal{S}
(\mathbb{R}^d)$ converging to $g$ with respect to the natural topology of
$\mathcal{C}^{s \asterisk}$ we have that $(X (g_n))_n$ converges in
$\mathcal{A}$ to a unique element
\[ X (g) \assign \lim_{n \rightarrow + \infty} X (g_n) . \]
The map $X$ now defined on all of $\mathcal{C}^{s \asterisk}$ is continuous.
In other words, if $X \in \mathcal{G}^s (V)$ then it can be identified with an
element of $\mathcal{L} (\mathcal{C}^{s \asterisk}, \mathcal{A})$. Furthermore
it is simple to prove, using suitable $\mathcal{S} (\mathbb{R}^d)$
approximations, that if $g, h \in \mathcal{C}^{s \asterisk}$, we have
\begin{equation}
  X (g) X (h) = - X (h) X (g) . \label{eq:anticomgeneral}
\end{equation}
In particular we have
\[ \dot{X} (x_1) \dot{X} (x_2) = - \dot{X} (x_2) \dot{X} (x_1), \qquad x_1,
   x_2 \in \mathbb{R}^d . \]
The relation~{\eqref{eq:anticomgeneral}} has important consequences. For any
$n \geqslant 1$, let
\[ \begin{array}{cc}
     \Lambda^n_{\pi} \mathcal{C}^{s \asterisk} \subset &
     \underbrace{\mathcal{C}^{s \asterisk} \otimes_{\pi} \mathcal{C}^{s
     \asterisk} \otimes_{\pi} \cdots \otimes_{\pi} \mathcal{C}^{s
     \asterisk}}\\
     & \text{$n$ times}
   \end{array} \]
be the completion of $\Lambda^n \mathcal{C}^{s \asterisk}$ in $\otimes^n_{\pi}
\mathcal{C}^{s \asterisk}$. Here $\otimes_{\pi}$ is the projective tensor
product of Banach spaces whose norm is defined as follows: if $f \in \otimes^n
\mathcal{C}^{s \asterisk}$ then
\[ \| f \|_{\otimes^n_{\pi} \mathcal{C}^{s \asterisk}} \assign \inf \left\{
   \sum_{i = 1}^k \prod_{j = 1}^n \| f^i_j \|_{\mathcal{C}^{s \asterisk}},
   \tmop{where} f = \sum_{i = 1}^k f^i_1 \otimes \cdots \otimes f^i_n \right\}
   . \]
\begin{lemma}
  \label{remark_regularity}If $X \in \mathcal{G}^s (V)$ then it can be
  extended to a continuous homomorphism from $\bigoplus_{i = 1}^n
  \Lambda_{\pi}^n \mathcal{C}^{s \asterisk}$ into $\mathcal{A}$ for any $n >
  0$.
\end{lemma}

\begin{proof}
  The proof is based on the fact that $X$ can be identified with an element of
  $\mathcal{L} (\mathcal{C}^{s \asterisk}, \mathcal{A})$ and on the
  relation~{\eqref{eq:anticomgeneral}}.
\end{proof}

In the case $s > 0$, we have for example that
\[ X (\delta_{x_1} \wedge \cdots \wedge \delta_{x_n}) = \dot{X} (x_1) \cdots
   \dot{X} (x_n), \qquad x_1, \ldots, x_n \in \mathbb{R}^d \]
We want now to define a natural notion of function $F$ of an element $X \in
\mathcal{G}^s (V)$. For example, if $s > 0$, and we write $Y = F_n (X)$, we
want to be able to take $F_n$ of the form
\begin{equation}
  \dot{Y} (x) = \sum_{(k_1, \ldots, k_n) \in D} K_{(k_1, \ldots, k_n)} \dot{X}
  (x + k_1) \dot{X} (x + k_2) \cdots \dot{X} (x + k_n)
  \label{eq:examplefinite},
\end{equation}
where $D \subset \mathbb{R}^{d n}$ is some finite set and $K_{(k_1, \ldots,
k_n)} \in \mathbb{R}$. When $n$ is odd, the important property of the function
$F_n$ defined before is that $Y = F (X) \in \mathcal{C}^s (\mathbb{R}^d,
\mathcal{A})$ and furthermore for any $g, h \in \mathcal{C}^{s \asterisk}$ we
have
\begin{equation}
  \begin{array}{c}
    Y (g) Y (h) = - Y (h) Y (g)\\
    Y (g) X (h) = - X (h) Y (g) .
  \end{array} \label{eq:anticom2}
\end{equation}
We want to consider some more general setting than the one considered in
equation~{\eqref{eq:examplefinite}} that maintains the important
property~{\eqref{eq:anticom2}}.

\

Consider $F_n \in \mathcal{L} (\mathcal{S} (\mathbb{R}^d), \Lambda^n_{\pi}
\mathcal{C}^{s \asterisk})$, then $F_n$ can be seen as an element of
$\mathcal{S}' (\mathbb{R}^d) \otimes \Lambda^n_{\pi} \mathcal{C}^{s
\asterisk}$. Using the same technique as before we can introduce the Banach
space $\mathcal{C}^s (\mathbb{R}^d, \Lambda^n_{\pi} \mathcal{C}^{s \asterisk})
\subset \mathcal{S}' (\mathbb{R}^d) \otimes \Lambda^n_{\pi} \mathcal{C}^{s
\asterisk}$ equipped with its natural norm. In this case we will we use the
notation
\begin{equation}
  F_n (X) \assign X \circ F_n, \label{eq:composition}
\end{equation}
where $(X \circ F_n) (v) = X (F_n (v)) \in \mathcal{A}$ for any $v \in
\mathcal{S} (\mathbb{R}^d)$.

\begin{definition}
  \label{definition-F1}We say that the linear function $F \in \mathcal{C}^s
  \left( \mathbb{R}^d, \bigoplus_{j = 0}^k \Lambda^{2 j + 1}_{\pi}
  \mathcal{C}^{s \asterisk} \right)$ is a standard odd function of degree $2 k
  + 1$ if there is a sequence of functions $F^n \in \mathcal{C}^s \left(
  \mathbb{R}^d, \bigoplus_{j = 0}^k \Lambda^{2 j + 1} \mathcal{C}^{s
  \asterisk} \right)$ (where $\Lambda^{2 j + 1} \mathcal{C}^{s \asterisk} =
  \underbrace{\mathcal{C}^{s \asterisk} \otimes \mathcal{C}^{s \asterisk}
  \otimes \cdots \otimes \mathcal{C}^{s \asterisk}}_{\text{$2 j + 1$ times}}$
  and $\otimes$ is the algebraic tensor product) such that for any $g \in V$
  we have $F^n (g) \rightarrow F (g)$, as $n \rightarrow + \infty$, in
  $\bigoplus_{j = 0}^k \Lambda^{2 j + 1}_{\pi} \mathcal{C}^{s \asterisk}$
  equipped with its standard topology.
\end{definition}

\begin{remark}
  When $s > 0$, an example of a function $F$ satisfying
  Definition~\ref{definition-F1} is
  \[ F (X) (\delta_x) = \dot{X} (x) \int_{\mathbb{R}^d} \mathcal{G}_1 (x - y)
     \dot{X} (y) \mathd y \int_{\mathbb{R}^d} \mathcal{G}_2 (x - y) \dot{X}
     (y) \mathd y, \]
  where $\mathcal{G}_1, \mathcal{G}_2 \in \mathcal{C}^{s + \epsilon}$ are
  functions that decrease faster than any polynomial at infinity. We note that
  a finite dimensional approximation of $F (X) (\delta_x)$ is given by
  \[ F^n (X) (\delta_x) = \sum_{D_i \in \mathcal{P}_n} \sum_{D_j \in
     \mathcal{P}_n} \dot{X} (x) \dot{X} (x - x_i) \dot{X} (x - x_j) \int_{D_i}
     \mathcal{G}_1 (y) \mathd y \int_{D_j} \mathcal{G}_2 (y) \mathd y, \]
  where $\mathcal{P}_n$ is sequence of increasing partitions of $\mathbb{R}^d$
  for which $\sup_i \left| \int_{D_i} \mathcal{G}_1 (y) \mathd y \right|
  \rightarrow 0$ and $\sup_j \left| \int_{D_j} \mathcal{G}_2 (y) \mathd y
  \right| \rightarrow 0$.
\end{remark}

\begin{theorem}
  \label{theorem_F}Let $s \in \mathbb{R}$, let $F \in \mathcal{C}^s \left(
  \mathbb{R}^d, \bigoplus_{j = 0}^k \Lambda^{2 j + 1}_{\pi} \mathcal{C}^{s
  \asterisk} \right)$ be a standard odd function of degree $2 k + 1$ (see
  Definition~\ref{definition-F1}), and consider $X \in \mathcal{G}^s (V)$.
  Then $F (X) \in \mathcal{C}^s (\mathbb{R}^d, \mathcal{A})$, furthermore for
  any $v_1, v_2 \in \mathcal{C}^{s \asterisk}$ we have
  \begin{equation}
    \begin{array}{c}
      F (X) (v_1) F (X) (v_2) = - F (X) (v_2) F (X) (v_1)\\
      F (X) (v_1) X (v_2) = - X (v_2) F (X) (v_1) .
    \end{array} \label{eq:anticommutation}
  \end{equation}
\end{theorem}

\begin{proof}
  The fact that $F (X) \in \mathcal{C}^s (\mathbb{R}^d, \mathcal{A})$ follows
  from Remark~\ref{remark_regularity}. The
  relations~{\eqref{eq:anticommutation}} are obvious when $F \in \mathcal{C}^s
  \left( \mathbb{R}^d, \bigoplus_{j = 0}^k \Lambda^{2 j + 1} \mathcal{C}^{s
  \asterisk} \right)$. Using a finite dimensional approximation $F^n$
  converging to $F$ the relations~{\eqref{eq:anticommutation}} are extended to
  the more general case.
\end{proof}

\begin{remark}
  An important consequence of Theorem~\ref{theorem_F} is that if $X \in
  \mathcal{G}^s (V)$ and $F$ satisfies the hypotheses of
  Theorem~\ref{theorem_F}, $F (X)$ can be extended to a Grassmann random
  variable (which in the following we denote by the same symbol $F (X)$) in
  $\mathcal{G}^s (V)$. Furthermore $F (X)$ and $X$ are compatible (in the
  sense of the corresponding definition of Section \ref{sec:grassmann-rv}).
\end{remark}

We are \ now in the position to give a precise meaning to
equation~{\eqref{eq:infinite5}} and to the various hypotheses on its data. Let
the operator $A$ be the (generally unbounded) generator of a strongly
continuous and exponentially contractive semigroup $(e^{A t})_{t \geqslant 0}$
on $L^2 (\mathbb{R}^d)$. Moreover consider the stationary Ornstein-Uhlenbeck
motion
\begin{equation}
  B^A_t (v) = \Xi (\mathbbm{1}_{(- \infty, t]} \otimes e^{A (t - \cdot)} v),
  \qquad t \in \mathbb{R}, v \in \mathcal{S} (\mathbb{R}^d), \label{eq:ba}
\end{equation}
where $\Xi$ is a Grassmann Gaussian noise on $\mathcal{G} (C^{\infty}
(\mathbb{R}_+) \otimes \mathcal{S} (\mathbb{R}^d))$ with covariance
$\mathbbm{1} \otimes C$ (we consider on $C^{\infty}_c (\mathbb{R}_+) \otimes
\mathcal{S} (\mathbb{R}^d)$ the pre-Hilbert space structure given by the
natural inclusion of $C^{\infty}_c (\mathbb{R}_+) \otimes \mathcal{S}
(\mathbb{R}^d)$ in $L^2 (\mathbb{R}_+ \times \mathbb{R}^d)$).

\begin{definition}
  \label{definition_F2}We say that the process $\Psi \in C^0 ([0, T],
  \mathcal{G}^s (V))$ satisfies equation~{\eqref{eq:infinite5}} with initial
  condition $\Psi_0$ if $F$ is \ a standard odd function of finite degree,
  $\Psi_0 \in \mathcal{G}^s (V)$ is independent of $\Xi$, and for any $t \in
  [0, T]$ and any $v \in V$ we have
  \begin{equation}
    \Psi_t (v) = \Psi_0 (e^{A t} v) + \lambda \int_0^t F (\Psi_s) (e^{A (t -
    s)} v) \mathd s + B^A_t (v) - B_0^A (e^{A t} v) \label{eq:infinite2},
  \end{equation}
  where the integral in the variable $s$ is taken in Bochner sense with
  respect the norm of $\mathcal{A}$.
\end{definition}

We introduce some conditions on $F$, $A$ and $\Xi$ in order to prove a result
on existence and uniqueness of solutions to~{\eqref{eq:infinite2}}.

\begin{definition}
  \label{definition_drift}Let $F$ be a standard odd function of degree $2 k +
  1$. We say that $F$ is a regular map on $\mathcal{C}^s (\mathbb{R}^d,
  \mathcal{A})$ if it is a locally bounded and Lipschitz function, i.e. it
  satisfies the following conditions:
  \begin{enumerateroman}
    \item There is a continuous non-decreasing function $f : \mathbb{R}_+
    \rightarrow \mathbb{R}_+$ such that for every $X \in \mathcal{G}^s (V)$ we
    have
    \begin{equation}
      \| F (X) \|_{\mathcal{C}^s (\mathbb{R}^d, \mathcal{A})} \leqslant f (\|
      X \|_{\mathcal{C}^s (\mathbb{R}^d, \mathcal{A})}) . \label{eq:bound1}
    \end{equation}
    \item The map $F$ is locally Lipschitz continuous and there is a
    continuous non-decreasing function $g : \mathbb{R}_+ \rightarrow
    \mathbb{R}_+$ such that, for any $X, Y \in \mathcal{G}^s (V)$:
    \begin{equation}
      d_{\mathcal{G}^s (V)} (F (X), F (Y)) \leqslant g (\max (\| X
      \|_{\mathcal{C}^s (\mathbb{R}^d, \mathcal{A})}, \| Y \|_{\mathcal{C}^s
      (\mathbb{R}^d, \mathcal{A})})) d_{\mathcal{G}^s (V)} (X, Y) .
      \label{eq:bound2}
    \end{equation}
  \end{enumerateroman}
\end{definition}

\begin{definition}
  \label{definition_infinitelinear}We say that the unbounded operator $A$ on
  $L^2 (\mathbb{R}^d)$, and the Gaussian noise $\Xi$ with covariance $C$ are
  adapted to the space $\mathcal{G}^s (V)$, if
  \begin{enumerateroman}
    \item The operator $A$ generates a contraction semigroup on $\mathcal{C}^s
    (\mathbb{R}^d)$ such that $\| e^{A t} \|_{\mathcal{L} (\mathcal{C}^s
    (\mathbb{R}^d), \mathcal{C}^s (\mathbb{R}^d))} \lesssim e^{- \lambda_A t}$
    for some strictly positive constant $\lambda_A > 0$.
    
    \item The Gaussian process $B^A$ defined by~{\eqref{eq:ba}} belongs to
    $C^0 (\mathbb{R}, \mathcal{G}^s (V))$ and furthermore we have $\sup_{t \in
    \mathbb{R}} \| B_t^A \|_{\mathcal{C}^s (\mathbb{R}^d, \mathcal{A})} < +
    \infty$.
  \end{enumerateroman}
\end{definition}

\begin{theorem}
  \label{theorem_existence2}Suppose that $F$ satisfies
  Definition~\ref{definition_drift} and $A, \Xi$ satisfy
  Definition~\ref{definition_infinitelinear} and let $X$ be a variable in
  $\mathcal{G}^s (V)$ independent of the Grassmann Gaussian noise $\Xi$. Then
  there exists a $T > 0$ such that there is a unique solution $\Psi \in C^0
  ([0, T], \mathcal{G}^s (V))$ to equation~{\eqref{eq:infinite2}}. Furthermore
  $\Psi_t$ is compatible with $(\Psi_0, \Xi, \Psi_s)$ for every $t, s \in [0,
  T]$.
\end{theorem}

\begin{proof}
  The proof is a straightforward generalization of that of
  Theorem~\ref{theorem_picard} for the case of our infinite dimensional
  setting.
\end{proof}

In order to obtain solutions for all times it is not possible to extend
directly Theorem~\ref{th:global-sol}, since its proof was essentially based on
the hypothesis that $V$ is finite dimensional. On the other hand the proofs of
Theorem~\ref{theorem_stationary}, Lemma~\ref{lemma_bound} and
Theorem~\ref{theorem_approximation} do not depend on the dimension. For this
reason we introduce the space $\mathfrak{K}_K \subset C^0 (\mathbb{R},
\mathcal{G}^s (V))$ defined as the set of functions $\Psi_t \in C^0
(\mathbb{R}, \mathcal{G}^s (V))$ such that
\[ \sup_{t \in \mathbb{R}} \| \Psi_t \|_{\mathcal{C}^s (\mathbb{R}^d,
   \mathcal{A})} < K, \]
for some constant $K > 0$, the function $L$ being given by
\[ L (A, C) = \sup_{t \in \mathbb{R}} \| B^A_t \|_{\mathcal{C}^s
   (\mathbb{R}^d, \mathcal{A})}, \]
and the equation for stationary solutions corresponding to
{\eqref{eq:infinite2}} as
\begin{equation}
  \Psi_t^s = \lambda \int_{- \infty}^t F (\Psi_{\tau}^s (e^{A (t - \tau)}))
  \mathd \tau + B^A_t \label{eq:stationary2} \qquad t \in \mathbb{R}.
\end{equation}
We also denote by $\Psi_{- T, t}^X$ the Grassmann process defined by $\Psi_{-
T, t}^X = X$ if $t \leqslant - T$ and such that $\Psi_{- T, t}^X$ is the
solution to equation~{\eqref{eq:infinite2}} with initial condition $X$ at time
$t = - T$. \

\begin{theorem}
  \label{theorem_longtime2}Suppose that $F$ satisfies
  Definition~\ref{definition_drift} and $A, \Xi$ satisfy
  Definition~\ref{definition_infinitelinear} and let $K > 2 L (A, C)$. Then
  there exists a function $\lambda_0 (K, L, F)$ such that for any $| \lambda |
  \leqslant 2 \lambda_0 (K, L, F)$ there exists a unique (stationary) solution
  to equation~{\eqref{eq:stationary2}} in $\mathfrak{K}_K$. Furthermore for
  any $| \lambda | \leqslant \lambda_0 (K, L, F)$ and for any random variable
  $X \in \mathcal{G}^s_{\mathcal{}} (V)$, independent of $B^A_t$ and such that
  $\| X \|_{\mathcal{C}^s (\mathbb{R}^d, \mathcal{A})} < \frac{1}{3} K$, we
  have that $\Psi^X_{- T, t}$ converges, as $T \rightarrow + \infty$, in
  $\mathcal{G}^s (V)$ to $\Psi^s_t$ uniformly on compact subsets of
  $\mathbb{R}$.
\end{theorem}

\begin{proof}
  Due to our appropriate infinite dimensional setting the proof is a
  straightforward generalization of Theorem~\ref{theorem_stationary},
  Lemma~\ref{lemma_bound} and Theorem~\ref{theorem_approximation}.
\end{proof}

\begin{remark}
  The function $\lambda_0$ defined in Theorem~\ref{theorem_longtime2} has to
  satisfy the following inequalities
  \begin{equation}
    \frac{2 \lambda_0 (K, L, F) f_F (K)}{\lambda_A} + 2 L < K
    \label{eq:contraction3},
  \end{equation}
  \begin{equation}
    \frac{2 \lambda_0 (K, L, F)}{\lambda_A} g_F (K) < 1
    \label{eq:contraction4}
  \end{equation}
  where $f_F$, $g_F$ and $\lambda_A$ are the quantities introduced in
  Definition~\ref{definition_drift} and
  Definition~\ref{definition_infinitelinear} respectively.
\end{remark}

\section{The Yukawa$_2$ model}\label{sec:yukawa}

We are now in position to study the stochastic quantization for an Euclidean
fermionic QFT. We consider the Euclidean counterpart of the Yukawa model
describing the interaction of a scalar field with a spin $1 / 2$ field in a $1
+ 1$ Minkowski space-time. However we will not attempt here to remove the UV
cutoff so the dimension of the space will not play any fundamental role. The
Yukawa$_2$ model contains already the main features of a large class of
Euclidean fermionic theories. It will be clear from the considerations below
that more general results are easily derivable at the price of heavier
notation. For the sake of readability we choose to stick to a definite model.

\subsection{Osterwalder--Schrader fields}

We define the Euclidean Yukawa$_2$ quantum field theory following Osterwalder
and Schrader~{\cite{osterwalder_euclidean_1973-2}}. In the language of this
paper, the Euclidean free Fermions introduced by Osterwalder and Schrader
(compare with~{\cite{osterwalder_euclidean_1973-2}} page~282 and
equation~(3.3) page ~284) are given by a quadruplet $\tilde{\psi} = (\psi,
\bar{\psi}) \in \mathcal{G} (\mathcal{S} (\mathbb{R}^2) \otimes \mathbb{C}^4)$
of complex Grassmann Gaussian fields on the complex pre-Hilbert space
$\mathcal{S} (\mathbb{R}^2) \otimes \mathbb{C}^4 \subseteq L^2 (\mathbb{R}^2)
\otimes \mathbb{C}^4$ with an Hermitian scalar product denoted $(\cdot,
\cdot)$, a standard bilinear form $\langle \cdot, \cdot \rangle_{\mathbb{C}}$
(defined by $\langle f, g \rangle \assign (\bar{f}, g)$, for $f, g \in
\mathcal{G} (\mathcal{S} (\mathbb{R}^2) \otimes \mathbb{C}^4)$) and the
correlations
\begin{equation}
  \omega (\tilde{\psi} (f) \tilde{\psi} (g)) = \langle f, K g
  \rangle_{\mathbb{C}}, \label{eq:free-corrrelation} \qquad f, g \in
  \mathcal{S} (\mathbb{R}^2) \otimes \mathbb{C}^4,
\end{equation}
where $K$ is the bounded operator on $L^2 (\mathbb{R}^2) \otimes \mathbb{C}^4$
given by
\[ K = (m_f^2 - \Delta)^{- 1} \left(\begin{array}{cc}
     0 & \left( \not{\nabla} + m_f \right)\\
     - \left( \not{\nabla} + m_f \right)^{\Tau} & 0
   \end{array}\right), \]
where $\Delta$ is the Laplacian on $\mathbb{R}^2$, $m_f > 0$ is a constant
interpreted as the mass of the fermionic particle associated with the
fermionic fields $(\psi, \bar{\psi})$. As before, $[\cdots]^{\Tau}$ denotes
real transposition, that is $[\cdots]^{\Tau} = \overline{[\cdots]}^{\ast}$.
Moreover $\not{\nabla} \assign \gamma_1 \nabla_1 + \gamma_2 \nabla_2$ with
$\nabla_i$ the partial derivatives with respect to $x_i$, $i = 1, 2$, and with
Euclidean Dirac matrices $\gamma_1, \gamma_2 : \mathbb{C}^2 \rightarrow
\mathbb{C}^2$ given by
\[ \gamma_1 = \left(\begin{array}{cc}
     0 & 1\\
     1 & 0
   \end{array}\right), \quad \gamma_2 = \left(\begin{array}{cc}
     0 & - \mathi\\
     \mathi & 0
   \end{array}\right), \]
and such that
\[ \gamma_a \gamma_b + \gamma_b \gamma_a = 2 \delta_{a b}, \qquad a, b = 1, 2.
\]
Note that $K$ is complex anti-symmetric: $K^{\Tau} = - K$, but not Hermitian
nor anti-Hermitian. Indeed $K = (m_f^2 - \Delta)^{- 1 / 2} U$ with $U$ an
unitary operator in $L^2 (\mathbb{R}^2) \otimes \mathbb{C}^4$. Hence the
correlation matrix of the field $\tilde{\psi}$ can also be written as
\[ \omega (\tilde{\psi} (f) \tilde{\psi} (g)) = \langle (m_f^2 - \Delta)^{- 1
   / 4} f, U (m_f^2 - \Delta)^{- 1 / 4} g \rangle_{\mathbb{C}} . \]
Therefore we can realize $\tilde{\psi}$ as
\[ \tilde{\psi} (f) \assign \zeta ((m_f^2 - \Delta)^{- 1 / 4} f), \qquad f \in
   \mathcal{S} (\mathbb{R}^2) \otimes \mathbb{C}^4, \]
where $\tilde{\zeta} \in \mathcal{G} (\mathcal{S} (\mathbb{R}^2) \otimes
\mathbb{C}^4)$ is the complex Grassmann Gaussian field with the bounded
covariance
\[ \omega (\tilde{\zeta} (f) \tilde{\zeta} (g)) = \langle f, U g
   \rangle_{\mathbb{C}} . \]
As a consequence, in this specific realization, we have
\begin{equation}
  \| \tilde{\psi} (f) \| = \| \tilde{\zeta} ((m_f^2 - \Delta)^{- 1 / 4} f) \|
  \lesssim \| (m_f^2 - \Delta)^{- 1 / 4} f \|_{L^2 (\mathbb{R}^2) \otimes
  \mathbb{C}^4}, \label{eq:dirac-reg}
\end{equation}
where the constant, in the inequality, depends only on $m_f$.

\

In order to describe the Yukawa$_2$ model we need also to introduce an
Euclidean boson field $\varphi$ of mass $m_b > 0$: it is the centered Gaussian
field on $\mathbb{R}^2$ with covariance $(m^2_b - \Delta)^{- 1}$ where $m_b >
0$ is looked upon as a fixed parameter and $\Delta$ is the Laplacian on
$\mathbb{R}^2$. We can realize $\varphi$ as the canonical multiplication
operator on the Hilbert space $\mathcal{H}_{\varphi} = L^2 \left(
\mathcal{C}^{- \kappa}_{\text{loc}} (\mathbb{R}^2), \nu \right)$ where $\nu$
is the centered Gaussian measure on $\mathcal{C}^{- \kappa}_{\text{loc}}
(\mathbb{R}^2)$ (the local version of the $\mathcal{C}^{- \kappa}$ Besov
space) with covariance $(m^2_b - \Delta)^{- 1}$ and $\kappa > 0$ is any small
positive quantity. We shall denote by $\omega_{\varphi}$ the state associated
with $\nu$ and $\mathcal{C}^{- \kappa}_{\text{loc}} (\mathbb{R}^2)$.

\

In addition, large scale and small scale cutoffs are introduced as follows.
Fix a smooth positive functions $h \in C^{\infty} (\mathbb{R}^2)$ with compact
support and let $\mathfrak{a}$ be a regular enough (see below for the precise
regularity) continuous function with compact support such that $\mathfrak{a}
(0) = 1$. Take $\varepsilon > 0$ and let $\mathfrak{a}_{\varepsilon} (x) =
\mathfrak{a} (x / \varepsilon) \varepsilon^{- 2}$. Then we define new
Grassmann variables
\[ \tilde{\psi}_{\varepsilon} (x) \assign \tilde{\psi}
   (\mathfrak{a}_{\varepsilon} (\cdot - x)), \qquad x \in \mathbb{R}^2 . \]
Then $\tilde{\psi}_{\varepsilon} \in \mathcal{G}^s (\mathcal{S} (\mathbb{R}^2)
\otimes \mathbb{C}^4)$ for some $s > 0$ (in general depending on the
regularity of $\mathfrak{a}$ see below).

\

Following the argument of Section~6 of~{\cite{osterwalder_euclidean_1973-2}}
we define the Euclidean Yukawa$_2$ model as follows.

\begin{definition}
  \label{def:yukawa2}The approximate fermionic Schwinger functions
  $(\rho^n_{\varepsilon, h})_n$ of the Yukawa$_2$ model on $\mathbb{R}^2$ with
  infrared (IR) cutoff $h$ and ultraviolet (UV) fermion cutoff $\varepsilon$
  are given by the Euclidean averages
  \begin{equation}
    \rho^n_{\varepsilon, h} (f_1, \ldots, f_n) = \frac{\omega \otimes
    \omega_{\varphi} (\tilde{\psi}_{\varepsilon} (f_1) \cdots
    \tilde{\psi}_{\varepsilon} (f_n) e^{- V_{\varepsilon, h}
    (\varphi)})}{\omega \otimes \omega_{\varphi} (e^{- V_{\varepsilon}})}
    \label{eq:euclidean-yukawa-main}
  \end{equation}
  with potential given by the unbounded (non-self-adjoint) operator on
  $\mathcal{H} \otimes \mathcal{H}_{\varphi}$, where $\mathcal{H}$ is the
  fermionic Hilbert space,
  \begin{equation}
    V_{\varepsilon, h} (\varphi) = \lambda \int_{\mathbb{R}^2} \mathd x h (x)
    (\bar{\psi}_{\varepsilon} (x) \cdot \psi_{\varepsilon} (x)) \varphi (x),
    \label{eq:yukawa-potential-boson}
  \end{equation}
  and $\lambda \in \mathbb{R}$ is a coupling constant.
\end{definition}

The potential $V_{\varepsilon, h} (\varphi)$ can be expressed in terms of the
fields $\tilde{\psi}_{\varepsilon}$ as follows
\[ V_{\varepsilon, h} (\varphi) = \frac{\lambda}{2} \int_{\mathbb{R}^2} \mathd
   x h (x) (\tilde{\psi}_{\varepsilon} (x) \cdot J
   \widetilde{\psi}_{\varepsilon} (x)) \varphi (x), \]
where
\[ J \assign \left( \begin{array}{cc}
     0 & -\mathbbm{1}_{}\\
     \mathbbm{1} & 0
   \end{array} \right) . \]
$\text{}$

Note that there is no difficulty in making sense
of~{\eqref{eq:euclidean-yukawa-main}}, due to the following observations (see
also~{\cite{osterwalder_euclidean_1973-2}}). The map $x \in \mathbb{R}^2
\mapsto \tilde{\psi}_{\varepsilon} (x) \in \mathcal{A}= \mathcal{L}
(\mathcal{H})$ is $\mathcal{C}^s$ due to the presence of the regularization
$\varepsilon > 0$ and the regularity of the Euclidean field $\tilde{\psi}$
given in~{\eqref{eq:dirac-reg}}. The massive Gaussian free field $\varphi$ has
(local) regularity $\mathcal{C}^{- \kappa}_{\tmop{loc}}$ for arbitrarily small
$\kappa$, therefore, taking $s > \kappa$, by a standard Besov estimates for
Banach valued functions, we have, for any \tmtextit{fixed} $\varphi \in
\mathcal{C}^{- \kappa}_{\tmop{loc}} (\mathbb{R}^2)$
\[ \| V_{\varepsilon, h} (\varphi) \|_{\mathcal{A}} \lesssim \lambda \|
   \tilde{\psi}_{\varepsilon} \|^2_{\mathcal{C}^{\kappa} (\mathcal{A})} \| h
   \varphi \|_{B^{- \kappa}_{1, 1} (\mathbb{R}^2)} \lesssim \lambda \| h
   \varphi \|_{\mathcal{C}^{- \kappa} (\mathbb{R}^2)} . \]
It is also not difficult to see that the map $V_{\varepsilon, h} :
\mathcal{C}^{- \kappa} \rightarrow \mathcal{A}$ given by $V_{\varepsilon, h} :
\varphi \mapsto V_{\varepsilon, h} (\varphi)$ is linear and continuous in
$\varphi$, therefore measurable and that
\[ \omega \otimes \omega_{\varphi} (\tilde{\psi} (f_1) \cdots \tilde{\psi}
   (f_n) e^{- V_{\varepsilon, h} (\varphi)}) = \omega (\tilde{\psi} (f_1)
   \cdots \tilde{\psi} (f_n) \omega_{\varphi} (e^{- V_{\varepsilon, h}
   (\varphi)})) \]
e.g. by proving this first for a piecewise approximations of $\varphi$ and
then taking the limit. Moreover we have (by the properties of the Bochner
integral)
\[ \| \omega_{\varphi} (e^{- V_{\varepsilon, h} (\varphi)}) \|_{\mathcal{A}}
   \leqslant \omega_{\varphi} (\exp (C | \lambda | \| h \varphi
   \|_{\mathcal{C}^{- \kappa}})) < \infty, \]
for all $\lambda \in \mathbb{R}$. We conclude that the averages
in~{\eqref{eq:euclidean-yukawa-main}} are well defined since
\[ | (\omega \otimes \omega_{\varphi}) (\tilde{\psi} (f_1) \cdots
   \tilde{\psi} (f_n) e^{- V_{\varepsilon, h}}) | \leqslant \| \tilde{\psi}
   (f_1) \cdots \tilde{\psi} (f_n) \omega_{\varphi} (e^{- V_{\varepsilon, h}
   (\varphi)}) \|_{\mathcal{A}} < \infty . \]
By expanding the quantity $\omega_{\varphi} (e^{- V_{\varepsilon, h}
(\varphi)}) \in \mathcal{A}$ in power series and integrating away the Gaussian
field we obtain
\[ \omega_{\varphi} (e^{- V_{\varepsilon, h} (\varphi)}) =
   e^{\mathcal{V}_{\varepsilon, h}}, \]
with a purely fermionic potential given by
\[ \mathcal{V}_{\varepsilon, h} (\tilde{\psi}_{\varepsilon}) = \lambda^2
   \int_{(\mathbb{R}^2)^2} (h \bar{\psi}_{\varepsilon} \cdot
   \psi_{\varepsilon}) (x) \mathcal{G} (x - y) (h \bar{\psi}_{\varepsilon}
   \cdot \psi_{\varepsilon}) (y) \mathd x \mathd y. \]
Here $\mathcal{G} = (m^2_b - \Delta)^{- 1}$ is the Green function of the
operator $(m^2_b - \Delta)$ and $m_b$ is the boson mass. This gives us the
purely fermionic representation of the Schwinger functions
{\eqref{eq:euclidean-yukawa-main}} given by:
\begin{equation}
  \rho^n_{\varepsilon, h} (f_1, \ldots, f_n) = \frac{\omega (\tilde{\psi}
  (f_1) \cdots \tilde{\psi} (f_n) e^{\mathcal{V}_{\varepsilon, h}})}{\omega
  (e^{\mathcal{V}_{\varepsilon, h}})} . \label{eq:euclidean-yukawa-fermion}
\end{equation}
Perturbation theory suggests that the boson field requires a mass
renormalization $a_{\varepsilon}$ which depends on the fermion cutoff. In
theory one would need to modify the Gaussian field covariance to include a
mass renormalization and let $\varphi$ have covariance $(m^2_b - \Delta +
a_{\varepsilon})^{- 1}$. However in this paper we will not discuss the removal
of the UV cutoff and keep $\varepsilon > 0$ fixed. Therefore, for simplicity,
we ignore this mass renormalization.

\subsection{Stochastic quantization}

The Grassmann SPDEs associated to the measure $\tilde{\omega} (\cdot) = \omega
(\cdot e^{\mathcal{V}_{\varepsilon, h}})$ is formally
\begin{equation}
  \mathd \tilde{\Psi}_{\varepsilon, t}^h = ((m_f^2 - \Delta)
  \tilde{\Psi}_{\epsilon, t}^h + \lambda^2 F_{\varepsilon, h, Y}
  (\tilde{\Psi}_{\varepsilon, t}^h)) \mathd t + \mathd \tilde{B}_{\varepsilon,
  t} \label{eq:stochasticquantization}
\end{equation}
where $\tilde{\Psi}_{\varepsilon, t}^h = (\Psi_{\varepsilon, t}^h,
\bar{\Psi}_{\varepsilon, t}^h)$,
\[ F_{\varepsilon, h, Y} (\tilde{\Psi}_{\varepsilon, t}^h) = C_f \left(
   \begin{array}{l}
     \mathfrak{a}_{\varepsilon}^{\asterisk 2} \ast \left( h
     \bar{\Psi}_{\varepsilon, t}^h \int \mathcal{G} (y - \cdot) h (y)
     \Psi_{\varepsilon, t}^h (y) \bar{\Psi}_{\varepsilon, t}^h (y) \mathd y
     \right)\\
     \mathfrak{a}_{\varepsilon}^{\asterisk 2} \ast \left( h \Psi_{\varepsilon,
     t}^h \int \mathcal{G} (y - \cdot) h (y) \Psi_{\varepsilon, t}^h (y)
     \bar{\Psi}_{\varepsilon, t}^h (y) \mathd y \right)
   \end{array} \right), \]
with
\[ C_f = \left(\begin{array}{cc}
     0 & \left( \not{\nabla} + m_f \right)\\
     - \left( \not{\nabla} + m_f \right)^{\Tau} & 0
   \end{array}\right), \]
and where $\mathfrak{a}_{\varepsilon} (x) = \varepsilon^{- 2} \mathfrak{a}
\left( \frac{x}{\varepsilon} \right)$, $\varepsilon > 0$, $\mathfrak{a} :
\mathbb{R}^2 \rightarrow \mathbb{R}^2$ is the regular enough function with
compact support (see below for the more precise requirements) considered
before, $\mathfrak{a}_{\varepsilon}^{\asterisk 2} = \mathfrak{a}_{\varepsilon}
\asterisk \mathfrak{a}_{\varepsilon}$, $h$ is a smooth function (in the
beginning with compact support, later on we will consider $h \equiv 1$).

The Brownian motion $\tilde{B}$ is described as follows. Let $\tilde{\Xi}$ be
the white noise on the pre-Hilbert space $L^2 (\mathbb{R}_+) \otimes
\mathcal{S} (\mathbb{R}^2) \otimes \mathbb{C}^4 \subseteq L^2 (\mathbb{R}_+)
\otimes B_{2, 2}^{- 1 / 2} (\mathbb{R}^2) \otimes \mathbb{C}^4$ equipped with
the scalar product
\begin{equation}
  ((f, g)) = ((m_f^2 - \Delta)^{- 1 / 4} f, (m_f^2 - \Delta)^{- 1 / 4} g)_{L^2
  (\mathbb{R}_+ \times \mathbb{R}^2) \otimes \mathbb{C}^4}
  \label{eq:scalarproduct} \qquad f, g \in L^2 (\mathbb{R}_+) \otimes
  \mathcal{S} (\mathbb{R}^2) \otimes \mathbb{C}^4,
\end{equation}
and having bounded covariance
\[ \tilde{C}_f = (m^2_f - \Delta)^{- 1 / 2} C_f . \]
The (two-sided) Brownian motion $\tilde{B}$ on $\mathcal{S} (\mathbb{R}^2)
\otimes \mathbb{C}^4$ is defined by
\[ \tilde{B}_t (v) = \Xi ((\mathbbm{1}_{t > 0} -\mathbbm{1}_{t < 0}) \otimes v),
   \qquad t \in \mathbb{R}, v \in \mathcal{S} (\mathbb{R}^2) \otimes
   \mathbb{C}^4 . \]
We set $\tilde{B}_{\varepsilon} = \mathfrak{a}_{\varepsilon} \asterisk
\tilde{B}$.

\

According to Definition~\ref{definition_F2} the stationary version of
equation~{\eqref{eq:stochasticquantization}} is the solution to the following
integral problem
\begin{equation}
  \tilde{\Psi}^{h, s}_{\varepsilon, t} (v) = \lambda^2 \int_{- \infty}^t
  F_{\varepsilon, h, Y} (\tilde{\Psi}_{\varepsilon, \tau}^{h, s}) (e^{A (t -
  \tau)} v) \mathd \tau + \tilde{B}^A_t (v) - \tilde{B}^A_0 (e^{A t} v),
  \label{eq:yukawas} \qquad t \in \mathbb{R}, v \in \mathbb{C}^4,
\end{equation}
where we take $A = \Delta - m_f^2$ and
\[ \tilde{B}^A_t (v) = \Xi (\mathbbm{1}_{[0, t]} \otimes e^{A t} v), \qquad t
   \in \mathbb{R}, v \in \mathcal{S} (\mathbb{R}^2) \otimes \mathbb{C}^4 . \]


From now on we will make the following assumption. We fix $\delta > 0$ and
$\gamma > 0$, let \ $\alpha = \gamma + 1 / 2 + 2 \delta$ and take
$\mathfrak{a} \in B^{\alpha}_{1, \infty} (\mathbb{R}^2) \cap
\mathcal{C}^{\alpha} (\mathbb{R}^2)$. Again we recall that $\varepsilon > 0$
will be fixed in all our study.

\

The main theorem of this section concerns the construction of a stochastic
quantization for~{\eqref{eq:euclidean-yukawa-fermion}} and the removal of the
IR cutoff $h$ for $\lambda$ small.

\begin{theorem}
  \label{theorem_main1}If $| \lambda |$ is small enough and $h$ is a smooth
  function with compact support, then, for any $k \geqslant 1$, $f_1, \ldots,
  f_k \in C^{\infty}_c (\mathbb{R}^2) \otimes \mathbb{C}^4$ and $t \in
  \mathbb{R}$, we have
  \begin{equation}
    \frac{\omega (\tilde{\psi} (f_1) \cdots \tilde{\psi} (f_k)
    e^{\mathcal{V}_{\varepsilon, h}})}{\omega (e^{\mathcal{V}_{\varepsilon,
    h}})} = \omega (\tilde{\Psi}_{\varepsilon, t}^{h, s} (f_1) \cdots
    \tilde{\Psi}_{\varepsilon, t}^{h, s} (f_k)) .
  \end{equation}
\end{theorem}

We are also able to obtain the following result on the removal of the IR
cutoff.

\begin{theorem}
  \label{theorem_main2}If $| \lambda |$ is small enough and if we denote by
  $\tilde{\Psi}_{\varepsilon, t}^s$ the solution to
  equation~{\eqref{eq:yukawas}} with $h \equiv 1$, we have that for any $f_1,
  \ldots, f_k \in C^{\infty}_c (\mathbb{R}^2) \otimes \mathbb{C}^4$ and for
  any $t \in \mathbb{R}$
  \begin{equation}
    \lim_{h_n \rightarrow 1} \frac{\omega \left( \tilde{\psi} (f_1) \cdots
    \tilde{\psi} (f_k) e^{\mathcal{V}_{\varepsilon, h_n}} \right)}{\omega
    \left( e^{\mathcal{V}_{\varepsilon, h_n}} \right)} = \omega
    (\tilde{\Psi}_{\varepsilon, t}^s (f_1) \cdots \tilde{\Psi}_{\varepsilon,
    t}^s (f_k)) .
  \end{equation}
  where $(h_n)_{n \geqslant 1}$ is any sequence of smooth functions with
  compact support converging to $1$ uniformly in $\mathcal{C}^{\gamma +
  \delta}$ (see below the definition of $\delta$) on compact subsets of
  $\mathbb{R}^2$.
\end{theorem}

The rest of the section focus on the proof of these two theorems. For the
moment we record here two technical lemmas.

\begin{lemma}
  \label{lemma_noise}We have $\tilde{B}^A_{\varepsilon} \in
  \mathcal{C}^{\delta} (\mathbb{R}, \mathcal{C}^{\gamma} (\mathbb{R}^2,
  \mathcal{A}) \otimes \mathbb{C}^4)$ and
  \[ \sup_{t \in \mathbb{R}} \| \tilde{B}^A_{\varepsilon, t}
     \|_{\mathcal{C}^{\gamma + 2 \delta} (\mathbb{R}^2, \mathcal{A}) \otimes
     \mathbb{C}^4} < + \infty ._{} \]
\end{lemma}

\begin{proof}
  It is sufficient to prove the lemma for $\varepsilon = 0$ and $\gamma = -
  \frac{1}{2} - 2 \delta$, since for $\varepsilon > 0$ we have
  $\mathfrak{a}_{\varepsilon} \asterisk : \mathcal{C}^s \rightarrow
  \mathcal{C}^{s + \alpha}$ by Theorem~\ref{theorem_convolution} in Appendix
  \ref{appendix_Besov}. By definition of Grassmann Gaussian random variables,
  if $K_i \in \mathcal{S} (\mathbb{R}^2)$ is the function corresponding to the
  Littlewood--Paley block $\Delta_i$ (see Appendix \ref{appendix_Besov} for
  the definition), we have
  \[ \sup_{x \in \mathbb{R}^2} \| \Delta_i \tilde{B}_{0, t} - \Delta_i
     \tilde{B}_{0, s} \|_{\mathcal{A} \otimes \mathbb{C}^4} \lesssim \left.
     \left\| \int_{- \infty}^t e^{(\Delta - m_f^2) (t - \tau)} (- \Delta +
     m^2)^{1 / 4} (K_i) \mathd \tau + \right. \right. \]
  \[ \left. \left. - \int_{- \infty}^s e^{(\Delta - m_f^2) (s - \tau)} (m^2_f
     - \Delta)^{1 / 4} (K_i) \mathd \tau \right\| \right._{L^2 (\mathbb{R}^2)}
     \lesssim \]
  \[ \lesssim \left\| \int_s^t e^{(\Delta + m_f^2) (t - \tau)} (m^2_f -
     \Delta)^{1 / 4} (K_i) \mathd \tau \right\| + \| (e^{(\Delta - m_f^2) (t -
     s)} - 1) (m^2_f - \Delta)^{1 / 4} (K_i) \|_{L^2} \lesssim \]
  \[ \lesssim \left( \int_{\mathbb{R}^2} \left( \phi_i (p) \left( \int_0^{t -
     s} | p |^2 e^{- k (| p |^2 + m^2_f)} \mathd k \right) \right)^2 | p |
     \mathd p \right) \lesssim \]
  \[ \lesssim \left( \int_{2^{i - 1} \lesssim | p | \lesssim 2^i} \left(
     \int_0^{t - s} \frac{\mathd k}{k^{(1 - \delta)}} \right)^2 | p |^{1 + 4
     \delta} \mathd p \right)^{1 / 2} \lesssim 2^{i \left( \frac{1}{2} + 2
     \delta \right)} | t - s |^{\delta} . \]
  From this the thesis follows from the definition of the norm of the Besov
  space $\mathcal{C}^{\gamma} (\mathbb{R}^2, \mathcal{A})$.
\end{proof}

\begin{lemma}
  \label{lemma_drift}For any $s > 0$, the function $F_{\varepsilon, h, Y}$ is
  well defined on $\mathcal{C}^s (\mathbb{R}^d, \mathcal{A})$ into itself, it
  satisfies Definition~\ref{definition-F1} and
  Definition~\ref{definition_drift}. Furthermore the functions $f_F, g_F$ in
  equations~{\eqref{eq:bound1}} and~{\eqref{eq:bound2}} (of
  Definition~\ref{definition_drift}) can be chosen to be independent of $h$
  when $h \in \mathcal{C}^s$ has the norm bounded by some constant $C > 0.$
\end{lemma}

\begin{proof}
  Since for $s > 0$ we have $\mathcal{C}^s (\mathbb{R}^d, \mathcal{A}) \subset
  C^0_b (\mathbb{R}^d, \mathcal{A})$ we have that the function
  \[ x \mapsto \int \mathcal{G} (y - x) h (y) \Psi_{\varepsilon, t}^h (y)
     \bar{\Psi}_{\varepsilon, t}^h (y) \mathd y \]
  is well defined and continuous. Since $\mathcal{C}^s (\mathbb{R}^d,
  \mathcal{A})$ is an algebra, this proves that $F_{\varepsilon, h, Y}$ is
  well defined, continuous and locally bounded. This proves that
  $F_{\varepsilon, h, Y}$ satisfies Definition~\ref{definition_drift}. The
  bounds required on $F_{\varepsilon, h, Y}$ are uniform in $h$ indeed
  \[ \| F_{\varepsilon, h, Y} \|_{\mathcal{C}^s} \lesssim \| \tilde{\Psi}^h
     \|_{\mathcal{C}^s} \left\| \int \mathcal{G} (x - y) h (y) \Psi^h (y)
     \bar{\Psi}^h (y) \mathd y \right\|_{\mathcal{C}^s} \lesssim \|
     \tilde{\Psi}^h \|_{\mathcal{C}^s} \| \Psi^h \bar{\Psi}^h
     \|_{\mathcal{C}^s} \lesssim \| \tilde{\Psi}^h \|_{\mathcal{C}^s}^3 . \]
  Since $F_{\varepsilon, h, Y}$ is of third degree, we need only to prove that
  it can be approximated with finite dimensional operators. This can be done
  using the fact that when $s > 0$, $\tilde{\Psi}_{\varepsilon, t}$ is
  H{\"o}lder continuous in space and that $\mathcal{G}$ decreases
  exponentially at infinity. This means that if $\mathcal{P}_n$ is a sequence
  of increasing (finite) partition of the plane we can approximate the
  integral by the sum
  \[ S_n = \sum_{D_i \in \mathcal{P}_n} h (x - x_i) \Psi^h (x - x_i)
     \bar{\Psi}^h (x - x_i) \int_{D_i} \mathcal{G} (y) \mathd y, \]
  where $x_i \in D_i$ is any point in $D_i$. If we choose the sequence of
  partitions $\mathcal{P}_n$ in such a way that $\lim_{n \rightarrow + \infty}
  \left\{ \sup_i \left( \int_{D_i} \mathcal{G} (y) \mathd y \right) \right\} =
  0$, we have that
  \[ S_n \rightarrow \int \mathcal{G} (y - x) h (y) \Psi_{\varepsilon, t}^h
     (y) \bar{\Psi}_{\varepsilon, t}^h (y) \mathd y \]
  in $\mathcal{C}^{s'} (\mathbb{R}^d, \mathcal{A})$, where $s' < s$.
\end{proof}

Using the previous two lemmas we can prove the following theorem.

\begin{theorem}
  \label{theorem_main3}Equation~{\eqref{eq:stochasticquantization}} admits a
  unique (local in time) solution, starting from $0$ in $\mathcal{G}^{\gamma}
  (V)$ ($\gamma$ being defined before Theorem~\ref{theorem_main1} and
  Theorem~\ref{theorem_main2}). Furthermore if $| \lambda |$ is small enough,
  there exists a unique solution to the stationary
  equation~{\eqref{eq:yukawas}} (associated with
  {\eqref{eq:stochasticquantization}}) which is uniformly bounded in
  $\mathcal{C}^{\gamma} (\mathbb{R}^d, \mathcal{A}) \otimes \mathbb{C}^4$ by
  some constant $K > 0$ (that can be chosen independent of $h$ when $\| h
  \|_{\mathcal{C}^{\gamma}} < C$ for some $C > 0$).
\end{theorem}

\begin{proof}
  The proof is based on Theorem~\ref{theorem_existence2} and
  Theorem~\ref{theorem_longtime2} whose hypotheses are satisfied in this
  setting because of Lemma~\ref{lemma_noise} and Lemma~\ref{lemma_drift}.
\end{proof}

\subsection{Finite dimensional approximations in finite volume}

The first step in order to prove Theorem~\ref{theorem_main1} and
Theorem~\ref{theorem_main2} is to introduce finite dimensional approximations
of the Yukawa model. Consider the torus $\mathbb{T}_R^2 =\mathbb{R}^2
\backslash (2 \pi L\mathbb{Z})^2$ for $R > 0$ and the periodic version
$\tilde{\psi}_R$ of the Fermion field $\tilde{\psi}$ defined by
\[ \tilde{\psi}_R (f) \assign \sum_{n \in \mathbb{Z}^2} \tilde{\psi} (f (\cdot
   - 2 \pi R n)), \]
for all $f \in \mathcal{S} (\mathbb{R}^2) \otimes \mathbb{C}^4$. Then letting
$\tilde{\psi}_{R, \varepsilon} (x) \assign \tilde{\psi}_R
(\mathfrak{a}_{\varepsilon} (\cdot - x))$ we have that the function $x \in
\mathbb{T}^2_R \mapsto \tilde{\psi}_{R, \varepsilon} (x)$ is well defined and
$\mathcal{C}^{\gamma + 2 \delta} (\mathbb{T}^2, \mathcal{A})$. Moreover
$\tilde{\psi}_{R, \varepsilon} (x)$ is uniformly bounded in $\mathcal{A} =
\mathcal{L} (\mathcal{H})$ and
\[ \| \tilde{\psi}_{R, \varepsilon} (x) - \tilde{\psi}_{\varepsilon} (x)
   \|_{\mathcal{C}^{\gamma + 2 \delta} (K, \mathcal{A} \otimes \mathbb{C}^4)}
   \rightarrow 0, \]
uniformly on compacts as $R \rightarrow \infty$. As a consequence we have
\begin{equation}
  \frac{\omega (\tilde{\psi} (f_1) \cdots \tilde{\psi} (f_n)
  e^{\mathcal{V}_{\varepsilon, h}})}{\omega (e^{\mathcal{V}_{\varepsilon,
  h}})} = \lim_{R \rightarrow \infty} \frac{\omega (\tilde{\psi}_R (f_1)
  \cdots \tilde{\psi}_R (f_n) e^{\mathcal{V}_{R, \varepsilon, h}})}{\omega
  (e^{\mathcal{V}_{R, \varepsilon, h}})}
  \label{eq:euclidean-yukawa-fermion-periodic}
\end{equation}
with
\[ \mathcal{V}_{R, \varepsilon, h} \assign \lambda^2 \int_{\mathbb{R}^4}
   \mathd x (h \bar{\psi}_{R, \varepsilon} \cdot \psi_{R, \varepsilon}) (x)
   \mathcal{G} (x - y) (h \bar{\psi}_{R, \varepsilon} \cdot \psi_{R,
   \varepsilon}) (y) \mathd x \mathd y, \]
for all compactly supported $f_1, \ldots, f_n, h \in \mathcal{S}
(\mathbb{R}^2)$.

\

Let now $(e_k)_{k \in \mathbb{Z}^2}$ be the orthonormal Fourier basis of (the
complex) $L^2 (\mathbb{T}_R^2)$ and $P_N : L^2 (\mathbb{T}_R^2) \rightarrow
L^2 (\mathbb{T}_R^2)$ the orthogonal projection to the closed subspace
generated by $\{ e_k : k \in \mathbb{Z}^2, | k | \leqslant N \}$. Let
$\tilde{\psi}_{N, R, \varepsilon} = P_N \tilde{\psi}_{R, \varepsilon}$ and
observe again that $\| \tilde{\psi}_{N, R, \varepsilon} (x) - \tilde{\psi}_{R,
\varepsilon} (x) \|_{\mathcal{C}^{\alpha} (\mathbb{R}^2)} \rightarrow 0$ as $N
\rightarrow \infty$, for any $\alpha > 0$. \ Defining
\[ \mathcal{V}_{N, R, \varepsilon, h} \assign \lambda^2 \int_{\mathbb{R}^4} (h
   \bar{\psi}_{N, R, \varepsilon} \cdot \psi_{N, R, \varepsilon}) (x)
   \mathcal{G}_{\varepsilon} (x - y) (h \bar{\psi}_{N, R, \varepsilon} \cdot
   \psi_{N, R, \varepsilon}) (y) \mathd x \mathd y \]
and we easily conclude that
\begin{equation}
  \frac{\omega (\tilde{\psi} (f_1) \cdots \tilde{\psi} (f_n)
  e^{\mathcal{V}_{\varepsilon, h}})}{\omega (e^{\mathcal{V}_{\varepsilon,
  h}})} = \lim_{R \rightarrow \infty} \lim_{N \rightarrow \infty} \frac{\omega
  (\tilde{\psi}_{N, R} (f_1) \cdots \tilde{\psi}_{N, R} (f_n)
  e^{\mathcal{V}_{N, R, \varepsilon, h}})}{\omega (e^{\mathcal{V}_{N, R,
  \varepsilon, h}})} \label{eq:euclidean-yukawa-fermion-periodic-galerkin} .
\end{equation}
Now consider the finite set $\Lambda_N = \{ k \in \mathbb{Z}^2 : | k |
\leqslant N \}$ and let
\[ \tilde{\theta} (k) \assign \int_{\mathbb{T}_R^2} \overline{e_k (x)}
   \tilde{\psi}_{N, R} (x) \mathd x, \qquad k \in \Lambda_N . \]
We can show easily that $(\tilde{\theta} (k))_{k \in \Lambda_N}$ is a
Grassmann Gaussian vector with correlation given by
\begin{equation}
  \omega (\tilde{\theta}_{\alpha} (k) \tilde{\theta}_{\beta} (\ell)) =
  \delta_{k + \ell = 0} \left(\begin{array}{cc}
    0 & \left( i \not{k} + m_f \right)^{- 1}\\
    - \left( i \not{k} + m_f \right)^{- 1\mathrm{T}} & 0
  \end{array}\right)_{\alpha, \beta}
  \label{eq:free-corrrelation-fourier-galerkin}
\end{equation}
where $\alpha, \beta = 1, \ldots, 4$ denote the canonical coordinates in
$\mathbb{C}^4$ and $k, \ell \in \Lambda_N$. Moreover
\[ \tilde{\psi}_{N, R} (x) = \sum_{k \in \Lambda_N} e_k (x) \tilde{\theta}
   (k), \qquad x \in \mathbb{T}^2_R \]
and $\tilde{\psi}_{N, R, \varepsilon} = \mathfrak{a}_{\varepsilon}
\ast_{\mathbb{T}_R^2} \tilde{\psi}_{N, R}$ where the convolution is done on
$\mathbb{T}^2_R$ as indicated.

\

At this point we have reduced via two subsequent approximations the number of
degrees of freedom to a finite amount, described by the Grassmann random
variable $\tilde{\theta} \in \mathcal{G} (\Lambda_N \otimes \mathbb{C}^4)$.

\

We introduce the following equation
\begin{equation}
  \tilde{\Psi}^{h_R, s}_{\varepsilon, N, R, t} (v) = \lambda^2 \int_{-
  \infty}^t P_N (F_{\varepsilon, h_R, Y} (P_N (\tilde{\Psi}_{\varepsilon, N,
  R, \tau}^{h_R, s}))) (e^{(\Delta - m^2_f) (t - \tau)} v) \mathd \tau +
  \tilde{B}_{\varepsilon, R, t}^A (v) \label{eq:finitedimensional}
\end{equation}
where $(\tilde{B}_{\varepsilon, R, t}^A)_{t \in \mathbb{R}}$ is the
Ornstein--Uhlenbeck motion on $\mathbb{T}^2_R$ obtained from
$(\tilde{B}_{\varepsilon, t}^A)_{t \in \mathbb{R}}$ via the restriction from
$\mathbb{R}^2$ to $\mathbb{T}^2$ of the periodicization
\[ \tilde{B}_{\varepsilon, R, t}^A (f) \assign \sum_{n \in \mathbb{Z}^2}
   \tilde{B}_{\varepsilon, t}^A (f (\cdot - 2 \pi R n)), \qquad f \in
   \mathcal{S} (\mathbb{R}^2), \]
and
\[ h_R (z) \assign \sum_{n \in \mathbb{Z}^2} h (z - 2 \pi R n) . \]
Obviously we have that
\[ (I - P_N) \tilde{\Psi}_{\varepsilon, N, R, t}^{h_R, s} = (I - P_N)
   \tilde{B}^A_{\varepsilon, R, t} . \]
This implies that $P_N (\tilde{\Psi}_{\varepsilon, N, R, t}^{h_R, s})$ solves
a finite dimensional Grassmann SDE. Furthermore the drifts $P_N
(F_{\varepsilon, h_R, Y} (P_N (\tilde{\Psi}_{\varepsilon, N, R, t}^{h_R,
s})))$ can be seen as the differentials of the function
\[ U (P_N (\tilde{\Psi})) = \int_{\mathbb{T}_L^4} h_L (x) h_L (y) (P_N (\Psi)
   P_N (\bar{\Psi})) (x) \mathcal{G} (x - y) (P_N (\Psi) P_N (\bar{\Psi})) (y)
   \mathd x \mathd y. \]
times the matrix $P_N \mathfrak{a}_{\varepsilon}^{\asterisk 2} \asterisk C_f$
(that is the covariance of the finite dimensional Brownian motion $P_N
(\tilde{B}_{\varepsilon, R, t})$). This implies the following key lemma.

\begin{lemma}
  \label{lemma_finite1}For any $k \geqslant 1$ and $f_1, \ldots, f_k$ some
  linear combination of $e_1, \ldots, e_N$ then, for $| \lambda |$ small
  enough (that can be chosen in a way independent of $h$) we have for any $t
  \in \mathbb{R}$
  \[ \frac{\omega (\tilde{\psi}_{N, R} (f_1) \cdots \tilde{\psi}_{N, R} (f_k)
     e^{\mathcal{V}_{N, R, \varepsilon, h}})}{\omega (e^{\mathcal{V}_{N, R,
     \varepsilon, h}})} = \omega (\tilde{\Psi}_{\varepsilon, N, R, t}^{h_R, s}
     (f_1) \cdots \tilde{\Psi}_{\varepsilon, N, R, t}^{h_R, s} (f_k)) . \]
  
\end{lemma}

\begin{proof}
  The proof is based on the fact that for any linear combination $f_i$ of
  $e_1, \ldots, e_N$ we have $\tilde{\Psi}_{\varepsilon, N, R, + \infty}^{h_L,
  s} (f_i) = P_N (\tilde{\Psi}_{\varepsilon, N, R, + \infty}^{h_R, s}) (f_i)$.
  Thus the thesis follows from Theorem~\ref{theorem_invariant2} and Remark
  \ref{remark_invariant2} by noting that $\mathcal{V}_{N, R, \varepsilon, h}
  (\tilde{\psi}) = \lambda^2 U (\tilde{\psi})$.
\end{proof}

\subsection{The infinite dimensional equation in finite volume}

Consider the equation
\begin{equation}
  \tilde{\Psi}_{\varepsilon, R, t}^{h_R, s} (v) = \lambda^2 \int_{- \infty}^t
  F_{\varepsilon, h_R, Y} (\tilde{\Psi}_{\varepsilon, R, \tau}^{h_R, s})
  (e^{(\Delta - m^2_f)} v) \mathd t + \tilde{B}^A_{\varepsilon, R, t}
  \label{eq:infinite1} .
\end{equation}
This is a modification of {\eqref{eq:finitedimensional}} where the finite
dimensional projections $P_N$ are removed from the non-linearity
$F_{\varepsilon, h_R, Y}$.

\begin{remark}
  Lemma~\ref{lemma_drift} holds for the drift of
  equation~{\eqref{eq:infinite1}}, and it is easy to generalize
  Lemma~\ref{lemma_noise} to the OU motion $(\tilde{B}_{\varepsilon, R,
  t}^A)_{t \in \mathbb{R}}$. This implies that it is possible to generalize
  Theorem~\ref{theorem_main3} to equation~{\eqref{eq:infinite1}} proving
  existence and uniqueness of its global solution (when $| \lambda |$ is small
  enough). Finally, it is important to note that the bound $K$ on the solution
  $\tilde{\Psi}_{\varepsilon, R, t}^{h_R, s}$ to
  equation~{\eqref{eq:infinite1}}, and the bounds on the constant $| \lambda
  |$ can be chosen independent of $h$ when $\| h \|_{\mathcal{C}^{\gamma}} <
  C$ for some constant $C > 0$.
\end{remark}

Equation~{\eqref{eq:infinite1}} is infinite dimensional and in order to
generalize Lemma~\ref{lemma_finite1} we cannot apply directly the finite
dimensional results contained in Theorem~\ref{theorem_invariant2} and
Remark~\ref{remark_invariant2}. For this reason we need to prove that the
solution to equation~{\eqref{eq:finitedimensional}} converges, as $N
\rightarrow + \infty$, to the solution to equation~{\eqref{eq:infinite1}}.

\begin{theorem}
  \label{theorem_convergence1}For any $0 < q < \gamma$ we have that
  $\tilde{\Psi}_{\varepsilon, N, R, t}^{h_R, s} \rightarrow
  \tilde{\Psi}_{\varepsilon, R, t}^{h_R, s}$, in $\mathcal{G}^q (C^{\infty}
  (\mathbb{T}^d) \otimes \mathbb{C}^4)$ when $N \rightarrow + \infty$ and
  uniformly in $t$.
\end{theorem}

\begin{proof}
  We have that the difference $\Psi_{\varepsilon, N, R, t}^{h_R, s} -
  \Psi_{\varepsilon, R, t}^{h_R, s} $ satisfies the following inequality
  \begin{eqnarray}
    \frac{\| \tilde{\Psi}_{\varepsilon, N, R, t}^{h_R, s} -
    \tilde{\Psi}_{\varepsilon, R, t}^{h_R, s} \|_{\mathcal{C}^q}}{\lambda^2} &
    \lesssim & \int_{- \infty}^t e^{- \lambda_A (t - \tau)} \| (I - P_N)
    F_{\varepsilon, h_R, Y} (P_N (\tilde{\Psi}_{\varepsilon, N, R, \tau}^{h_R,
    s})) \|_{\mathcal{C}^q} \mathd \tau + \nonumber\\
    &  & + \int_{- \infty}^t e^{- \lambda_A (t - \tau)} g_F (\max (\|
    \tilde{\Psi}_{\varepsilon, N, R, \tau}^{h_R, s} \|_{\mathcal{C}^q}, \| P_N
    \tilde{\Psi}_{\varepsilon, N, R, \tau}^{h_R, s} \|_{\mathcal{C}^q}))
    \times \nonumber\\
    &  & \times (\| \tilde{\Psi}_{\varepsilon, N, R, \tau}^{h_R, s} - P_N
    \tilde{\Psi}_{\varepsilon, N, R, \tau}^{h_R, s} \|_{\mathcal{C}^q}) \mathd
    \tau + \nonumber\\
    &  & + \int_{- \infty}^t e^{- \lambda_A (t - \tau)} g_F (\max (\|
    \tilde{\Psi}_{\varepsilon, N, R, \tau}^{h_R, s} \|_{\mathcal{C}^q}, \|
    \tilde{\Psi}_{\varepsilon, R, \tau}^{h_R, s} \|_{\mathcal{C}^q})) \times
    \nonumber\\
    &  & \times \| \tilde{\Psi}_{\varepsilon, N, R, \tau}^{h_R, s} -
    \tilde{\Psi}_{\varepsilon, R, \tau}^{h_R, s} \|_{\mathcal{C}^q} \mathd
    \tau . \nonumber
  \end{eqnarray}
  Using Gronwall inequality we obtain that
  \[ \| \tilde{\Psi}_{\varepsilon, N, R, t}^{h_R, s} -
     \tilde{\Psi}_{\varepsilon, R, t}^{h_R, s} \|_{\mathcal{C}^q} \lesssim (\|
     I - P_N \|_{\mathcal{L} (\mathcal{C}^{\gamma}, \mathcal{C}^q)}) \times \]
  \[ \times \left( 1 + \int_{- \infty}^t \exp (- (\lambda_A - \lambda^2 g_F
     (K)) (t - \tau)) \mathd \tau \right) . \]
  Since $\lambda$ is small we have that $\lambda_A - \lambda^2 g_F (K) > 0$
  (this is exactly the request~{\eqref{eq:contraction4}} for $| \lambda | <
  \lambda_0 (K, L, F)$). This means that
  \[ \| \tilde{\Psi}_{\varepsilon, N, R, t}^{h_R, s} -
     \tilde{\Psi}_{\varepsilon, R, t}^{h_R, s} \|_{\mathcal{C}^q} \lesssim (\|
     I - P_N \|_{\mathcal{L} (\mathcal{C}^{\gamma}, \mathcal{C}^q)}) . \]
  On the other hand, since, \ $\gamma > q$ by assumption, we have $\| I - P_N
  \|_{\mathcal{L} (\mathcal{C}^{\gamma}, \mathcal{C}^q)} \rightarrow 0$ as $N
  \rightarrow + \infty$. The thesis is proven noting that the convergence
  $\tilde{\Psi}_{\varepsilon, N, R, t}^{h_R, s} \rightarrow
  \tilde{\Psi}_{\varepsilon, R, t}^{h_R, s}$ is uniform in $t$.
\end{proof}

An easy corollary of Theorem~\ref{theorem_convergence1},
Lemma~\ref{lemma_finite1}, and the
convergence~{\eqref{eq:euclidean-yukawa-fermion-periodic-galerkin}} is the
following.

\begin{corollary}
  \label{corollary_formula}For $| \lambda |$ small enough (that can be chosen
  independent of $h$ when $\| h \|_{\mathcal{C}^{\gamma}} < C$ where $C > 0$),
  for any $k \geqslant 1$, $f_1, \ldots, f_k \in C^{\infty} (\mathbb{T}^d_L)
  \otimes \mathbb{C}^4$ and any $t \in \mathbb{R}$ we have
  \[ \frac{\omega (\tilde{\psi}_R (f_1) \cdots \tilde{\psi}_R (f_k)
     e^{\mathcal{V}_{R, \varepsilon, h}})}{\omega (e^{\mathcal{V}_{R,
     \varepsilon, h}})} = \omega (\tilde{\Psi}_{\varepsilon, R, t}^{h_R, s}
     (f_1) \cdots \tilde{\Psi}_{\varepsilon, R, t}^{h_R, s} (f_k)) . \]
\end{corollary}

\subsection{The infinite volume limit}

In order to tackle the infinite volume limit and prove
Theorem~\ref{theorem_main1} and Theorem~\ref{theorem_main2} we have to
introduce weighted Besov spaces (see the Appendix~\ref{appendix_Besov}
and~{\cite{TriebelIII}} for the details). Let $\ell \in \mathbb{R}$ and
consider $\rho_{\ell} : \mathbb{R}^d \rightarrow \mathbb{R}_+$ defined by
$\rho_{\ell} (x) = (1 + | x |^2)^{(- \ell / 2)}$. Consider the following norm
on $\mathcal{S}' (\mathbb{R}^d) \otimes \mathcal{A}$ (where as before
$\mathcal{A} = \mathcal{L} (\mathcal{H})$ and $\mathcal{H}$ is the fermionic
Hilbert space)
\[ \| u \|_{\mathcal{C}^s_{\ell} (\mathbb{R}^d, \mathcal{A})} = \sup_{j
   \geqslant - 1} \sup_{x \in \mathbb{R}^d} 2^{j s} \| \Delta_j u (x)
   \rho_{\ell} (x) \|_{\mathcal{A}} . \]
Define the space $\mathcal{C}^s_{\ell} (\mathbb{R}^d, \mathcal{A})$ as the
subspace of $\mathcal{S}' (\mathbb{R}^d) \otimes \mathcal{A}$ where $\| \cdot
\|_{\mathcal{C}^s_{\ell} (\mathbb{R}^d, \mathcal{A})}$ is finite. We define
also $\mathcal{G}^s_{\ell} (V) = \mathcal{C}^s_{\ell} (\mathbb{R}^d,
\mathcal{A}) \cap \mathcal{G}^{- \infty} (V)$ (recall that $V =\mathcal{S}
(\mathbb{R}^d) \otimes \mathbb{R}^n$ or $\mathcal{S} (\mathbb{R}^d) \otimes
\mathbb{C}^n$ for some $n \geqslant 1$) and the distance
$d_{\mathcal{G}^s_{\ell} (V)} (X - Y) = \| X - Y \|_{\mathcal{C}^s_{\ell}
(\mathbb{R}^d, \mathcal{A})}$.

\begin{remark}
  \label{remark_weight}It is important to note that, by
  Theorem~\ref{theorem_weight},
  \[ \sup_{x \in \mathbb{R}^d} \| \Delta_j u (x) \rho_{\ell} (x)
     \|_{\mathcal{A}} \sim \sup_{x \in \mathbb{R}^d} \| \Delta_j (\rho_{\ell}
     u) (x) \|_{\mathcal{A}}, \]
  where $\sim$ means that both quantities considered can be bounded from above
  and from below by some positive constant times the other quantity.
\end{remark}

\begin{lemma}
  \label{lemma_convergenceB}Let $\gamma, \delta$ be as stated just before
  Theorem~\ref{theorem_main1}, for any $s < \gamma + 2 \delta$ and any $\ell >
  0$ we have that
  \[ \sup_{t \in \mathbb{R}} d_{\mathcal{G}^s_{\ell} (V)}
     (\tilde{B}_{\varepsilon, t}^A, \tilde{B}^A_{\varepsilon, R, t})
     \rightarrow 0 \]
  when $R \rightarrow + \infty$.
\end{lemma}

\begin{proof}
  Let $v$ be a function on $\mathbb{R}^2$ decreasing faster then any
  polynomial at infinity. Then we have
  \[ \tilde{B}_{\varepsilon, t}^A (v) - \tilde{B}^A_{\varepsilon, L, t} (v) =
     \tilde{\Xi}_{\varepsilon} (\mathbb{I}_{| x |_{\infty} > R} (\cdot) e^{A
     (t - \cdot)} v) - \tilde{\Xi}_{\varepsilon, L} (\mathbb{I}_{| x
     |_{\infty} > R} (\cdot) e^{A (t - \cdot)} v) . \]
  This means that $\tilde{B}^A_{\varepsilon, t}$ and
  $\tilde{B}^A_{\varepsilon, R, t}$ coincide on the set $\{ | x |_{\infty} < R
  \} \subset \mathbb{R}^2$. Using Remark~\ref{remark_weight}, this implies
  that
  \[ \| \Delta_i (\tilde{B}_{\varepsilon, t}^A - \tilde{B}^A_{\varepsilon, R,
     t}) \rho_{\ell} \|_{\mathcal{A} \otimes \mathbb{C}^4} \lesssim \|
     \Delta_i (\rho_{\ell} (\tilde{B}_{\varepsilon, t}^A -
     \tilde{B}^A_{\varepsilon, R, t})) \|_{\mathcal{A} \otimes \mathbb{C}^4}
     \lesssim \]
  \[ \lesssim R^{- \ell} \| \Delta_i (\tilde{B}_{\varepsilon, t}^A -
     \tilde{B}^A_{\varepsilon, R, t}) \|_{\mathcal{A} \otimes \mathbb{C}^4}
     \lesssim R^{- \ell} 2^{j (- 2 \delta - \gamma)}, \]
  which implies that $\sup_{t \in \mathbb{R}} \| \tilde{B}_{\varepsilon, t}^A
  - \tilde{B}^A_{\varepsilon, R, t} \|_{\mathcal{C}^s_{\ell} (\mathbb{R}^2,
  \mathcal{A} \otimes \mathbb{C}^4)} \lesssim R^{- \ell} \rightarrow 0$, when
  $\ell > 0$ as $R \rightarrow + \infty$.
\end{proof}

We now prove the analogue of Theorem~\ref{theorem_convergence1} in this case.

\begin{theorem}
  \label{theorem_torusplane}Let \ $0 < q < \gamma$, then
  $\tilde{\Psi}_{\varepsilon, R, t}^{h_R, s} \rightarrow
  \tilde{\Psi}_{\varepsilon, t}^{h, s}$, in $\mathcal{G}^q_{\ell} (\mathcal{S}
  (\mathbb{R}^2) \otimes \mathbb{C}^4)$ uniformly in $t \in \mathbb{R}$, when
  $R \rightarrow + \infty$.
\end{theorem}

\begin{proof}
  The proof is based on the fact that for any $\tilde{\Psi}_1, \tilde{\Psi}_2
  \in \mathcal{} \mathcal{G}^{\gamma}_0 (\mathcal{S} (\mathbb{R}^2) \otimes
  \mathbb{C}^4)$ (it is important to note that in these spaces we take the
  constant weight) we have, by Theorem~\ref{theorem_multiplication} in
  Appendix \ref{appendix_Besov},
  \[ \| F_{\varepsilon, h_R, Y} (\tilde{\Psi}_1) - F_{\varepsilon, h, Y}
     (\tilde{\Psi}_2) \|_{\mathcal{C}^q_{\ell} (\mathbb{R}^2, \mathcal{A}
     \otimes \mathbb{C}^4)} \lesssim g_F (\max (\| h_R \tilde{\Psi}_1
     \|_{\mathcal{C}^q}, \| h \tilde{\Psi}_1 \|_{\mathcal{C}^q}, \| h
     \tilde{\Psi}_2 \|_{\mathcal{C}^q})) \times \]
  \begin{equation}
    \times (\| h \|_{\mathcal{C}^q} \| \tilde{\Psi}_1 - \tilde{\Psi}_2
    \|_{\mathcal{C}^q_{\ell} (\mathbb{R}^2, \mathcal{A} \otimes \mathbb{C}^4)}
    + \| \Psi_1 \|_{\mathcal{C}^q_0 (\mathbb{R}^2, \mathcal{A} \otimes
    \mathbb{C}^4)} \| h_R - h \|_{\mathcal{C}^q_{\ell}})
    \label{eq:inequalityh},
  \end{equation}
  where $g_F$ is the function in Definition~\ref{definition_drift} and being
  $h$ the IR cut-off in Definition~\ref{def:yukawa2}. Using a reasoning
  similar to the one exploited in the proof of
  Theorem~\ref{theorem_convergence1} we get
  \[ \| \tilde{\Psi}_{\varepsilon, R, t}^{h_R, s} -
     \tilde{\Psi}_{\varepsilon, t}^{h, s} \|_{\mathcal{C}^q_{\ell}
     (\mathbb{R}^2, \mathcal{A} \otimes \mathbb{C}^4)} \lesssim (\| h_R - h
     \|_{\mathcal{C}^q_{\ell}} + \sup_{\tau \in \mathbb{R}} \|
     \tilde{B}_{\varepsilon, \tau}^A - \tilde{B}^A_{\varepsilon, R, \tau}
     \|_{\mathcal{C}^q_{\ell} (\mathbb{R}^2, \mathcal{A} \otimes
     \mathbb{C}^4)}) \times \]
  \[ \times \left( 1 + \int_{- \infty}^t \exp (- (\lambda_A - \lambda^2 g_F
     (K)) (t - \tau)) \mathd \tau \right) . \]
  By $| \lambda |$ small enough, this converges to $0$ as $R \rightarrow +
  \infty$ uniformly in $t$, since, by Lemma~\ref{lemma_convergenceB},
  \[ \lim_{R \rightarrow + \infty} \sup_{\tau \in \mathbb{R}} \|
     \tilde{B}_{\varepsilon, \tau}^A - \tilde{B}^A_{\varepsilon, R, \tau}
     \|_{\mathcal{C}^q_{\ell} (\mathbb{R}^2, \mathcal{A} \otimes
     \mathbb{C}^4)} \rightarrow 0 \]
  and $\| h_R - h \|_{\mathcal{C}^q_{\ell}} \rightarrow 0$, as $R \rightarrow
  + \infty$, by the definition of $h_R$.
\end{proof}

Hereafter we denote by $\tilde{\Psi}_{\varepsilon, t}^s$ the solution to
equation~{\eqref{eq:stochasticquantization}} with $h \equiv 1$ (so with the
spatial cut-off $h$ removed).

\begin{theorem}
  \label{theorem_cutoffh}Using the previous notations and hypotheses, let
  $(h_n : \mathbb{R}^2 \rightarrow \mathbb{R}^{}_{})_{n \geqslant 1}$ be a
  sequence of smooth functions with compact supports converging to $1$
  uniformly on compacts in the norm of $\mathcal{C}^{\gamma}$, $\gamma > 0$,
  and with uniformly bound in $\mathcal{C}^{\gamma}$. For any $0 < q < \gamma$
  we have
  \[ \tilde{\Psi}_{\varepsilon, t}^{h_n, s} \rightarrow
     \tilde{\Psi}_{\varepsilon, t}^s \]
  in $\mathcal{G}^q_{\ell} (\mathcal{S} (\mathbb{R}^2) \otimes \mathbb{C}^4)$,
  for any $\ell > 0$ and uniformly in $t \in \mathbb{R}$ when $n \rightarrow +
  \infty$.
\end{theorem}

\begin{proof}
  The proof is analogous to the one of Theorem~\ref{theorem_torusplane} where
  in inequality~{\eqref{eq:inequalityh}} we have $\| h_n - 1
  \|_{\mathcal{C}^q_{\ell}}$ (instead of $\| h_R - h
  \|_{\mathcal{C}^q_{\ell}}$) which converges to $0$, as $n \rightarrow +
  \infty$, by hypothesis when $\ell > 0$.
\end{proof}

\begin{proof*}{Proof of Theorem~\ref{theorem_main1} and
Theorem~\ref{theorem_main2}}
  The proof is analogous to the one of Corollary~\ref{corollary_formula} where
  Theorem~\ref{theorem_convergence1} is replaced by
  Theorem~\ref{theorem_torusplane} and Theorem~\ref{theorem_cutoffh}
  respectively.
\end{proof*}

\appendix\section{Besov spaces of Banach algebras}\label{appendix_Besov}

In this section we want to recall some results about Besov spaces of functions
(or distributions) from $\mathbb{R}^d$ taking values in a Banach algebra
$\mathcal{A}$. All the results of this section can be found
in~{\cite{Amann1997,Amann2019}} for the theory of Besov spaces taking values
in a Banach space and~{\cite{TriebelIII}} for weighted Besov spaces. We
present here only the $\mathbb{R}^d$ case. The definition of Besov spaces on
$\mathbb{T}^d$ is similar, and all the results announced here also hold on
$\mathbb{T}^d$.

\

We denote by $\mathcal{S} (\mathbb{R}^d)$ the set of smooth functions $f \in
\mathcal{S} (\mathbb{R}^d)$ such that
\[ \| f \|_{\ell, \alpha} \assign \| (1 + | x |)^{\ell} | D^{\alpha} f |
   \|_{L^{\infty} (\mathbb{R}^d)} < + \infty \]
where $\ell \in \mathbb{R}_+$ and $\alpha \in \mathbb{N}^d$. We denote by
$\mathcal{S}' (\mathbb{R}^d)$ the strong dual of $\mathcal{S} (\mathbb{R}^d)$
(equipped with the topology of the semi-norms $\| \cdot \|_{\ell, \alpha}$)
and by $\mathcal{O}_M (\mathbb{R}^d)$ the set of smooth functions $f \in
\mathcal{O}_M (\mathbb{R}^d)$ such that for any $\alpha \in \mathbb{N}^d$
there exists $\ell_{\alpha} \in \mathbb{R}_+$ for which
\[ \| (1 + | x |)^{- \ell_{\alpha}} | D^{\alpha} f | \|_{L^{\infty}
   (\mathbb{R}^d)} < + \infty . \]
We use the notations
\[ \mathcal{S} (\mathbb{R}^d, \mathcal{A}) = \mathcal{S} (\mathbb{R}^d)
   \otimes \mathcal{A}, \quad \mathcal{S}' (\mathbb{R}^d, \mathcal{A}) =
   \mathcal{S}' (\mathbb{R}^d) \otimes \mathcal{A}, \quad \mathcal{O}_M
   (\mathbb{R}^d, \mathcal{A}) = \mathcal{O}_M (\mathbb{R}^d) \otimes
   \mathcal{A}, \]
where $\mathcal{A}$ is a Banach algebra. It is important to note that
\[ \mathcal{S} (\mathbb{R}^d, \mathcal{A}) \overset{d}{\hookrightarrow}
   \mathcal{O}_M (\mathbb{R}^d, \mathcal{A}) \overset{d}{\hookrightarrow}
   \mathcal{S}' (\mathbb{R}^d, \mathcal{A}) \]
where the arrows mean that each space is continuously embedded and dense in
the following one.

\begin{theorem}
  It is possible to define uniquely $\langle \cdot, \cdot \rangle :
  \mathcal{S} (\mathbb{R}^d, \mathcal{A}) \times \mathcal{S}' (\mathbb{R}^d,
  \mathcal{A}) \rightarrow \mathcal{A}$, $\cdot : \mathcal{S} (\mathbb{R}^d,
  \mathcal{A}) \times \mathcal{S}' (\mathbb{R}^d, \mathcal{A}) \rightarrow
  \mathcal{S}' (\mathbb{R}^d, \mathcal{A})$, $\cdot : \mathcal{S}
  (\mathbb{R}^d) \times \mathcal{S}' (\mathbb{R}^d, \mathcal{A}) \rightarrow
  \mathcal{S}' (\mathbb{R}^d, \mathcal{A})$, $\asterisk : \mathcal{S}
  (\mathbb{R}^d) \times \mathcal{S}' (\mathbb{R}^d, \mathcal{A}) \rightarrow
  \mathcal{O}_M (\mathbb{R}^d, \mathcal{A})$ and $D^{\alpha} : \mathcal{S}'
  (\mathbb{R}^d, \mathcal{A}) \rightarrow \mathcal{S}' (\mathbb{R}^d,
  \mathcal{A})$ (where $\alpha \in \mathbb{N}^d$) which extend in a
  epicontinuous way the following operations: any $f \in \mathcal{S}
  (\mathbb{R}^d)$, $u \in \mathcal{S}' (\mathbb{R}^d)$ and $a_1, a_2 \in
  \mathcal{A}$ we have
  \[ \begin{array}{rll}
       \langle f \otimes a_1, u \otimes a_2 \rangle & = & \langle f, u \rangle
       a_1 a_2\\
       (f \otimes a_1) \cdot (u \otimes a_2) & = & (f u) \otimes (a_1 a_2)\\
       f \cdot (u \otimes a_2) & = & (f u) \otimes a_1\\
       f \asterisk (u \otimes a_2) & = & (f \asterisk u) \otimes a_1\\
       D^{\alpha} (u \otimes a_1) & = & (D^{\alpha} u) \otimes a_1
     \end{array} \]
  where $\langle f, u \rangle$ is the normal pairing in $\mathcal{S}
  (\mathbb{R}^d) \times \mathcal{S}' (\mathbb{R}^d)$, $(f u)$ is the product
  in $\mathcal{S} (\mathbb{R}^d) \times \mathcal{S}' (\mathbb{R}^d)$, $(f
  \asterisk u)$ is the convolution in $\mathcal{S} (\mathbb{R}^d) \times
  \mathcal{S}' (\mathbb{R}^d)$ and $D^{\alpha}$ is the $\alpha$ derivatives in
  $\mathcal{S}' (\mathbb{R}^d)$.
\end{theorem}

\begin{proof}
  The proof of this theorem can be found in~{\cite{Amann2019}} Appendix~1.
\end{proof}

We recall the definition of Littlewood--Paley blocks: let $\chi, \varphi$ be
smooth non-negative functions from $\mathbb{R}^n$ into $\mathbb{R}$ such that
\begin{itemize}
  \item $\tmop{supp} (\chi) \subset B_{\frac{4}{3}} (0)$ and $\tmop{supp}
  (\varphi) \subset B_{\frac{8}{3}} (0) \setminus B_{\frac{3}{4}} (0)$,
  
  \item $\chi, \varphi \leq 1$ and $\chi (y) + \sum_{j \geq 0} \varphi (2^{-
  j} y) = 1$ for any $y \in \mathbb{R}^n$,
  
  \item $\tmop{supp} (\chi) \cap \tmop{supp} (\varphi (2^{- i} \cdot)) =
  \emptyset$ for $i \geqslant 1$,
  
  \item $\tmop{supp} (\varphi (2^{- j} \cdot)) \cap \tmop{supp} (\varphi (2^{-
  i} \cdot)) = \emptyset$ if $| i - j | > 1$,
\end{itemize}
where by $B_r (x)$ we denote the ball of center $x \in \mathbb{R}^n$ and
radius $r > 0$. We introduce the following notations: $\varphi_{- 1} = \chi$,
$\varphi_j (\cdot) = \varphi (2^{- j} \cdot)$, $K_j = \hat{\varphi}_j$ and,
for any $\ell > 0$, we define $\rho_{\ell} (y) = (1 + | x |^2)^{- \ell / 2}$.

\

If $v \in \mathcal{S}' (\mathbb{R}^d, \mathcal{A})$ and if $i \in
\mathbb{Z}$, $i \geqslant - 1$ we define $i$-th Littlewood--Paley block as
follows
\[ \Delta_i v = K_i \asterisk v \in \mathcal{O}_M (\mathbb{R}^d, \mathcal{A})
   . \]
Then, if $s \in \mathbb{R}$, $p, q \in [1, + \infty]$ and $\ell \in
\mathbb{R}$ we define the function
\[ \| v \|_{B^s_{p, q, \ell} (\mathbb{R}^d, \mathcal{A})} = \left( \sum_{j = -
   1}^{+ \infty} 2^{j s q} \| \Delta_j v \|_{L^p_{\ell} (\mathbb{R}^d,
   \mathcal{A})}^q \right)^{1 / q}, \]
when $q \in [1, + \infty)$ and $\| v \|_{B^s_{p, + \infty, \ell}
(\mathbb{R}^d, \mathcal{A})} = \sup_j (2^{j s} \| \Delta_j v \|_{L^p_{\ell}
(\mathbb{R}^d, \mathcal{A})})$, where $\| \cdot \|_{L^p_{\ell} (\mathbb{R}^d,
\mathcal{A})}$ is the norm in the space $L_{\ell}^p (\mathbb{R}^d,
\mathcal{A})$ that is
\[ \| f \|_{L^p_{\ell} (\mathbb{R}^d, \mathcal{A})} = \left(
   \int_{\mathbb{R}^n} (\| f (y) \|_{\mathcal{A}} \rho_{\ell} (y))^p \mathd y
   \right)^{1 / p}, \]
if $p \in [1, + \infty)$ and
\[ \| f \|_{L^{\infty}_{\ell} (\mathbb{R}^d, \mathcal{A})} = \sup_{y \in
   \mathbb{R}^d} (\| f (y) \|_{\mathcal{A}} \rho_{\ell} (y)), \]
if $p = + \infty$. For any $v \in \mathcal{S} (\mathbb{R}^d, \mathcal{A})$ the
norm $\| v \|_{B^s_{p, q, \ell} (\mathbb{R}^d, \mathcal{A})} < + \infty$ is
finite. Then we look at $B^s_{p, q, \ell} (\mathbb{R}^d, \mathcal{A})$ as the
closure of $\mathcal{S} (\mathbb{R}^d, \mathcal{A})$ in $\mathcal{S}'
(\mathbb{R}^d, \mathcal{A})$ with respect to the norm $\| \cdot \|_{B^s_{p, q,
\ell} (\mathbb{R}^d, \mathcal{A})}$. Hereafter, if $s \in \mathbb{R}$, $p, q
\in [1, + \infty]$ and $\ell \in \mathbb{R}$, we use the following notations
$\mathcal{C}^s_{\ell} (\mathbb{R}^d, \mathcal{A}) \assign B^s_{\infty, \infty,
\ell} (\mathbb{R}^d, \mathcal{A})$, $B^s_{p, q} (\mathbb{R}^d, \mathcal{A}) =
B^s_{p, q, 0} (\mathbb{R}^d, \mathcal{A})$, $\mathcal{C}^s (\mathbb{R}^d,
\mathcal{A}) \assign B^s_{\infty, \infty, 0} (\mathbb{R}^d, \mathcal{A})$,
$B^s_{p, q, \ell} \assign B^s_{p, q, \ell} (\mathbb{R}^d, \mathbb{R})$ etc.

\

In this paper we need only the next three results.

\begin{theorem}
  \label{theorem_weight}For any $s, \ell \in \mathbb{R}$ and $v \in
  \mathcal{C}^s_{\ell} (\mathbb{R}^d, \mathcal{A})$ we have that $\rho_{\ell}
  v \in \mathcal{C}^s (\mathbb{R}^d, \mathcal{A})$ and
  \[ \| v \|_{\mathcal{C}^s_{\ell} (\mathbb{R}^d, \mathcal{A})} \sim \|
     \rho_{\ell} v \|_{\mathcal{C}^s (\mathbb{R}^d, \mathcal{A})} . \]
\end{theorem}

\begin{proof}
  By Theorem~6.5 in~{\cite{TriebelIII}} for weighted Besov spaces on
  $\mathbb{R}^d$ with polynomial-like weights, the theorem for a general
  Banach algebra $\mathcal{A}$ is a straightforward generalization.
\end{proof}

\begin{theorem}
  \label{theorem_convolution}Let $\alpha > 0$ and consider $\mathfrak{a} \in
  B^{\alpha}_{1, \infty}$ then we have that the convolution $\mathfrak{a}
  \asterisk$ can be extended in a unique continuous way from an operator from
  $\mathcal{S} (\mathbb{R}^d, \mathcal{A})$ into itself, into an operator from
  $\mathcal{C}^s (\mathbb{R}^d, \mathcal{A})$ into $\mathcal{C}^{s + \alpha}
  (\mathbb{R}^d, \mathcal{A})$.
\end{theorem}

\begin{proof}
  Using the notations of~{\cite{Amann1997}}, if $\mathfrak{a \in}
  B^{\alpha}_{1, \infty}$ then $((- \Delta + 1)^{\alpha / 2} (\mathfrak{a}))
  \in \mathcal{F} L_1$, then the thesis follows from Corollary~6.4
  in~{\cite{Amann1997}}.
\end{proof}

\begin{theorem}
  \label{theorem_multiplication}Let $s > 0$ and $\ell_1, \ell_2 \in
  \mathbb{R}$ we have that the product $\cdot : \mathcal{S} (\mathbb{R}^d,
  \mathcal{A}) \times \mathcal{S} (\mathbb{R}^d, \mathcal{A}) \rightarrow
  \mathcal{S} (\mathbb{R}^d, \mathcal{A})$ can be (uniquely) extended in a
  continuous way from $\mathcal{C}^s_{\ell_1} (\mathbb{R}^d, \mathcal{A})
  \times \mathcal{C}^s (\mathbb{R}^d, \mathcal{A})$ into
  $\mathcal{C}^s_{\ell_1 + \ell_2} (\mathbb{R}^d, \mathcal{A})$, furthermore
  for any $v_1 \in \mathcal{C}^s_{\ell_1} (\mathbb{R}^d, \mathcal{A})$ and
  $v_2 \in \mathcal{C}^s_{\ell_2} (\mathbb{R}^d, \mathcal{A})$ we have
  \[ \| v_1 \cdot v_2 \|_{\mathcal{C}^s_{\ell_1 + \ell_2} (\mathbb{R}^d,
     \mathcal{A})} \lesssim \| v_1 \|_{\mathcal{C}^s_{\ell_1} (\mathbb{R}^d,
     \mathcal{A})}  \| v_2 \|_{\mathcal{C}^s_{\ell_2} (\mathbb{R}^d,
     \mathcal{A})}, \]
\end{theorem}

\begin{proof}
  The proof can be found in~{\cite{Amann2019}} in the unweighted case. The
  proof in the Besov spaces with polynomial weights $\rho_{\ell}$ is similar.
\end{proof}

\section{Series expansion of solutions}\label{app:series}

In this appendix we sketch some implications of the exclusion principle for
the existence of global solutions to finite-dimensional Grassmann SDEs. In
principle some of these considerations also apply to some more realistic
models like the Yukawa$_2$ model however, to our surprise, stochastic
quantization, at least in the perturbative regime, can be carried on without
establishing them.

\

The solutions to the non-linear equation~{\eqref{eq:sde-lambda}} can be
represented via a series. We investigate now the properties of this series and
its absolute convergence. The explicit control of the series is not needed for
the development of the theory of stochastic quantization as we have
shown.\quad For simplicity let us assume that $(v_{\alpha})_{\alpha = 1,
\ldots, N}$ is a finite basis of $V$ and that $F \in \tmop{Hom} (V, \Lambda
V)$ is given by a sum of cubic monomials
\[ F (v_{\alpha}) = \sum_{\alpha_1, \alpha_2, \alpha_3}
   \lambda^{\alpha}_{\alpha_1 \alpha_2 \alpha_3} \psi (v_{\alpha_1}) \psi
   (v_{\alpha_2}) \psi (v_{\alpha_2}) \]
where $\psi \in \tmop{Hom} (V, \Lambda V)$ is the canonical injection of $V$
into $\Lambda V$ and $(\lambda^{\alpha}_{\alpha_1 \alpha_2 \alpha_3})$ a
family of coefficients in $\mathbb{R}$. Moreover we take $A = -\mathbb{I}$.
equation~{\eqref{eq:sde-lambda}} has then the integral formulation given by
\begin{equation}
  \Psi_t (v) = \Phi_t (v) + \int_0^t \Psi_s (e^{- (t - s)} F (v)) \mathd s,
  \qquad t \geqslant 0, v \in V. \label{eq:mild}
\end{equation}
where
\[ \Phi_t (v) \assign \Psi_0 (e^{- t} v) + B^A_t (v) - B^A_0 (e^{- t} v),
   \qquad t \geqslant 0, v \in V. \]
By iteratively expanding $\Psi_s$ on the right hand side of~{\eqref{eq:mild}}
we obtain a series expansion for $\Psi_t$ of the form
\[ \Psi_t = \sum_{\tau} J_{\tau} (\Phi) (t) \]
\[ = J_{\bullet} (\Phi) (t) + J_{[\bullet \bullet \bullet]} (\Phi) (t) +
   J_{[[\bullet \bullet \bullet] \bullet \bullet]} (\Phi) (t) + \cdots +
   J_{[[\bullet \bullet \bullet] [\bullet [\bullet \bullet [\bullet \bullet
   \bullet]] \bullet] \bullet]} (\Phi) (t) + \cdots \]
The series is indexed by (planar) trees $\tau$ which have branches of order 3
and where $J$ is a multilinear integral operator such that
\[ J_{\bullet} (\Phi) (t)^{\alpha} = \Phi_t (v_{\alpha}) \]
\[ J_{[\tau_1 \tau_2 \tau_3]} (\Phi) (t)^{\alpha} = \sum_{\alpha_1, \alpha_2,
   \alpha_3} \int_0^t e^{- (t - s)} \lambda^{\alpha}_{\alpha_1, \alpha_2,
   \alpha_3} J_{\tau_1} (\Phi) (s)^{\alpha_1} J_{\tau_2} (\Phi) (s)^{\alpha_2}
   J_{\tau_3} (\Phi) (s)^{\alpha_3} \mathd s \]
where $\bullet$ denotes the simple tree and $[\tau_1, \ldots, \tau_3]$ the
tree with branches $\tau_1, \ldots, \tau_3$. Our goal is to prove that the
above series converges for all times and derive estimates for the norm $\|
\Psi_t \|$ of the solution $\Psi$. Expanding the expression for $J_{\tau}
(\Phi)$ we have
\[ J_{\tau} (\Phi) = \int \left( \prod_{p \in I (\tau)}
   \lambda^{\alpha_p}_{\alpha_{p_1} \alpha_{p_2} \alpha_{p_3}} \right) \left(
   \prod_{p \in L (\tau)} \Phi^{\alpha_p} (s_{p_-}) \right) \left( \prod_{p
   \in I (\tau)} \mathd s_p \right) \]
where $I (\tau)$ denotes the internal nodes, $L (\tau)$ the leaves, $p$ the
parent node and $p_1, p_2, p_3$ the three children of every internal node and
$\Phi^{\alpha} (t) = \Phi_t (v_{\alpha})$. We have $| \tau | = | L (\tau) | +
| I (\tau) |$, $| \tau | = 1 + 3 | I (\tau) |$, $| L (\tau) | = 1 + 3 | I
(\tau) | - | I (\tau) | = 1 + 2 | I (\tau) |$.

\

Note that $(\Phi^{\alpha}_t)^2 = 0$ and that the increments $\Phi^{\alpha}
(t, s) \assign \Phi^{\alpha}_t - \Phi^{\alpha}_s$ have the norm bound
\[ \| \Phi^{\alpha} (t, s) \|^2 \leqslant \int_s^t e^{- (t - r)} \mathd r + |
   e^{- (t - s)} - 1 |^2 \leqslant 2 (1 - e^{- (t - s)}) \backassign H (t - s)
\]
which can be estimated as
\[ \| \Phi^{\alpha} (t, s) \| \leqslant | t - s |^{1 / 2} . \]
The following Lemma~uses $(\Phi^{\alpha}_t)^2 = 0$ (Pauli exclusion principle)
to derive good estimates for products of fields.

\begin{lemma}
  Assume $m \geqslant 1$. For any $n \geqslant 1$ and any $t_1, \ldots, t_n
  \in [0, T]$, $\alpha_1, \ldots, \alpha_n \in \{ 1, \ldots, N \}$ we have
  \[ \| \Phi^{\alpha_1}_{t_1} \cdots \Phi^{\alpha_n}_{t_n} \| \leqslant
     \frac{C^{n + 1} T^{n / 2}}{(n!)^{1 / 2}} \]
  where $C$ is a universal constant depending only on $N$.
\end{lemma}

\begin{proof}
  Partition $[0, T]$ in $r = n / (2 N)$ equal intervals $(I_k)_k$ and let
  $(s_k)_k$ be the centers of those intervals. Now in the product replace each
  $\Phi^{\alpha_i}_{t_1}$ by $\Phi^{\alpha_i}_{t_1} = \Phi^{\alpha_i}_{s_k} +
  \Phi^{\alpha_i} (t_i, s_k)$ where $s_k$ is the nearest to $t_i$ of the
  centers. By doing so we rewrite $\Phi^{\alpha_1}_{t_1} \cdots
  \Phi^{\alpha_n}_{t_n}$ as a sum $S$ of $2^n$ products of fields which are
  either $\Phi^{\alpha_i}_{s_k}$ or \ $\Phi^{\alpha_i} (t_i, s_k)$, moreover
  for a given $I_k$ there cannot be more than $N$ factors
  $\Phi^{\alpha_i}_{s_k}$ by the exclusion principle. Denote by $n_k$ the
  number of $(t_i)_i$ in $I_k$. Then $n = \sum_k n_k$ and using the estimate
  $\| \Phi^{\alpha_i} (t_i, s_k) \| \leqslant | t_i - s_k | \leqslant | I_k |
  = (2 N T) / n$ and the fact that in each cell we have at most $n_k - N$
  increments, we have
  \[ \| S \| \leqslant 2^n \prod_{k = 1}^r | I_k |^{n_k - N} \leqslant 2^n
     \left( \frac{2 N T}{n} \right)^{n - N r} = 2^n \left( \frac{2 N T}{n}
     \right)^{n / 2} \leqslant \frac{C^{n + 1} T^{n / 2}}{(n!)^{1 / 2}} . \]
\end{proof}

Using this lemma above we can estimate
\[ \| J_{\tau} (\Phi) (t) \| \leqslant \int \left( \prod_{p \in I (\tau)}
   \left| \lambda^{\alpha_p}_{\alpha_{p_1} \alpha_{p_2} \alpha_{p_3}} \right|
   \right) \left\| \prod_{p \in L (\tau)} \Phi^{\alpha_p}_{s_{p_-}} \right\|
   \left( \prod_{p \in I (\tau)} \mathd s_p \right) \]
\[ \lesssim | \lambda |^{| I (\tau) |} N^{| L (\tau) |} \frac{C^{| L (\tau) |
   + 1} t^{| L (\tau) | / 2}}{(| L (\tau) | !)^{1 / 2}} \]
where we estimated all the integrals by constants uniformly in $t$ due to the
presence of the exponential factors. The number of trees $\tau$ with a given
number of nodes $| \tau | = n$ is no more than $D^n$ for some $n$ and $D > 0$
and $| L (\tau) | = 1 + 2 | I (\tau) | = 1 + (2 / 3) (| \tau | - 1) = 1 / 3 +
2 | \tau | / 3$ so we obtain at the end for the solution $\Psi_t$ of
{\eqref{eq:mild}}
\[ \| \Psi_t \| \leqslant \sum_{\tau} \| J_{\tau} (\Phi) (t) \| \lesssim
   \sum_{n \geqslant 0} D^n | \lambda |^{(2 / 3) (n - 1)} N^{2 n / 3}
   \frac{C^{2 n / 3} t^{n / 3 + 1 / 6}}{(n / 3 + 1 / 6) !} \]
which is a series which converges for all $t$ (with some stretched exponential
behavior).

\begin{theorem}
  There exists an increasing function $E (t)$ depending on $N, | \lambda |, m$
  such that for all $m > 0$,
  \[ \sup_{s \leqslant t} \| \Psi_s \| \leqslant E (t) . \]
\end{theorem}

In particular this shows again that explosion is not possible for this kind of
SDEs, unlike their bosonic analogs.

\

Our goal now is to remove the time dependence on our estimate for the
solution exploiting the exponential decay. Let us partition the interval $[0,
\infty]$ into intervals $(I_h)_{h \geqslant 0}$ of size $1$. In each of these
intervals we will use the finer partition of the previous lemma to estimate
products, which will then not depend on $T$ anymore. Of course now the problem
is that we have a bound of the form
\[ \| \Phi^{\alpha_1} (t_1) \cdots \Phi^{\alpha_n} (t_n) \| \leqslant \prod_{h
   \geqslant 0} \frac{C^{n_h + 1}}{(n_h !)^{1 / 2}} \]
where $n_h$ is the number of time variables in $I_h$ and $\sum_{h \geqslant 0}
n_h = n$, and so far there is no useful upper bound on this (since we observe
that we can have the situation where $n_h = 1$ for all $h$ and so we loose the
factorial contribution). Let $Q$ be the number of intervals with $n_h > 0$. In
the expression for $J_{\tau} (\Phi) (t)$ there are at least $Q$ factors of the
form $e^{- (s - s')}$, $s' < s$, for which the sum of the quantities $s - s'$
is at least $Q$. So we have the bounds
\[ \| J_{\tau} (\Phi) (t) \| \leqslant | \lambda |^{(2 / 3) (n - 1)} N^n
   \sum_{Q \geqslant 1} e^{- Q / 2} \prod_{h : n_h > 0} \frac{C^{n_h +
   1}}{(n_h !)^{1 / 2}} \leqslant \frac{C^n | \lambda |^{(2 / 3) (n -
   1)}}{(n!)^{1 / 2}} \sum_{Q \geqslant 1} e^{- Q / 2} Q^{n / 2} \]
where we used that
\[ \prod_{h : n_h > 0} \frac{1}{n_h !} \leqslant \frac{1}{n!}
   \sum_{\tmscript{\begin{array}{c}
     k_1, \ldots, k_Q \geqslant 1\\
     k_1 + \cdots + k_Q = n
   \end{array}}} \frac{n!}{k_1 ! \cdots k_Q !} = \frac{1}{n!} Q^n . \]
Now, using the bound $\sum_{Q \geqslant 1} e^{- Q / 2} Q^{n / 2} \lesssim c^n
(n!)^{1 / 2}$ we obtain that
\[ \| J_{\tau} (\phi) (t) \| \lesssim | \lambda |^{(2 / 3) (n - 1)} C^n \]
so provided $| \lambda |$ is small enough (depending on our choice of $N$) we
conclude that the uniform bound $\sup_{t \geqslant 0} \| \Psi_t \| \lesssim 1$
holds.

\begin{thebibliography}{10}
  \bibitem[1]{accardi_outline_2015}L.~Accardi. {\newblock}An outline of
  quantum probability. {\newblock}2015. {\newblock}Preprint
  \url{http://rgdoi.net/10.13140/RG.2.1.3078.3844}.{\newblock}
  
  \bibitem[2]{accardi_quantum_1982}L.~Accardi, A.~Frigerio, and  J.~T.~Lewis.
  {\newblock}Quantum stochastic processes. {\newblock}\tmtextit{Kyoto
  University. Research Institute for Mathematical Sciences. Publications},
  18(1):97--133, 1982.{\newblock}
  
  \bibitem[3]{albeverio_elliptic_2018}S.~Albeverio, F.~C.~De Vecchi, and 
  M.~Gubinelli. {\newblock}Elliptic stochastic quantization.
  {\newblock}\tmtextit{ArXiv:1812.04422 [math-ph]}, dec 2018.
  {\newblock}ArXiv: 1812.04422.{\newblock}
  
  \bibitem[4]{albeverio_elliptic_2019}S.~Albeverio, F.~C.~De Vecchi, and 
  M.~Gubinelli. {\newblock}The elliptic stochastic quantization of some two
  dimensional Euclidean QFTs. {\newblock}\tmtextit{ArXiv:1906.11187
  [math-ph]}, jun 2019. {\newblock}ArXiv: 1906.11187.{\newblock}
  
  \bibitem[5]{albeverio_invariant_2017}S.~Albeverio  and  S.~Kusuoka.
  {\newblock}The invariant measure and the flow associated to the
  $\Phi^4_3$-quantum field model. {\newblock}\tmtextit{Annali della Scuola
  Normale di Pisa - Classe di Scienze}, 2018.
  {\newblock}\href{Https://doi.org/10.2422/2036-2145.201809_008}{}10.2422/2036-2145.201809\_008.{\newblock}
  
  \bibitem[6]{Amann1997}H.~Amann. {\newblock}Operator-valued Fourier
  multipliers, vector-valued Besov spaces, and applications.
  {\newblock}\tmtextit{Math. Nachr.}, 186:5--56, 1997.{\newblock}
  
  \bibitem[7]{Amann2019}H.~Amann. {\newblock}\tmtextit{Linear and quasilinear
  parabolic problems. Vol. II},  volume  106  of \tmtextit{Monographs in
  Mathematics}. {\newblock}Birkh{\"a}user/Springer, Cham, 2019.
  {\newblock}Function spaces.{\newblock}
  
  \bibitem[8]{Applebaum2}D.~B.~Applebaum  and  R.~L.~Hudson.
  {\newblock}Fermion It{\^o}'s formula and stochastic evolutions.
  {\newblock}\tmtextit{Comm. Math. Phys.}, 96(4):473--496, 1984.{\newblock}
  
  \bibitem[9]{Applebaum1}D.~Applebaum  and  R.~L.~Hudson. {\newblock}Fermion
  diffusions. {\newblock}\tmtextit{J. Math. Phys.}, 25(4):858--861,
  1984.{\newblock}
  
  \bibitem[10]{Segalbook}J.~C.~Baez, I.~E.~Segal, and  Z.-F.~Zhou.
  {\newblock}\tmtextit{Introduction to algebraic and constructive quantum
  field theory}. {\newblock}Princeton Series in Physics. Princeton University
  Press, Princeton, NJ, 1992.{\newblock}
  
  \bibitem[11]{Doplicher}D.~Bahns, S.~Doplicher, G.~Morsella, and 
  G.~Piacitelli. {\newblock}Quantum spacetime and algebraic quantum field
  theory. {\newblock}In \tmtextit{Advances in algebraic quantum field theory},
  Math. Phys. Stud.,  pages  289--329. Springer, Cham, 2015.{\newblock}
  
  \bibitem[12]{Streater3}C.~Barnett, R.~F.~Streater, and  I.~F.~Wilde.
  {\newblock}The It{\^o}-Clifford integral. {\newblock}\tmtextit{J. Functional
  Analysis}, 48(2):172--212, 1982.{\newblock}
  
  \bibitem[13]{Streater2}C.~Barnett, R.~F.~Streater, and  I.~F.~Wilde.
  {\newblock}The It{\^o}-Clifford integral. II. Stochastic differential
  equations. {\newblock}\tmtextit{J. London Math. Soc. (2)}, 27(2):373--384,
  1983.{\newblock}
  
  \bibitem[14]{Streater1}C.~Barnett, R.~F.~Streater, and  I.~F.~Wilde.
  {\newblock}The It{\^o}-Clifford integral. III. The Markov property of
  solutions to stochastic differential equations. {\newblock}\tmtextit{Comm.
  Math. Phys.}, 89(1):13--17, 1983.{\newblock}
  
  \bibitem[15]{Baumgartel1}H.~Baumg{\"a}rtel. {\newblock}\tmtextit{Operator
  algebraic methods in quantum field theory}. {\newblock}Akademie Verlag,
  Berlin, 1995. {\newblock}A series of lectures.{\newblock}
  
  \bibitem[16]{Baumgartel2}H.~Baumg{\"a}rtel  and  M.~Wollenberg.
  {\newblock}\tmtextit{Causal nets of operator algebras},  volume~80  of
  \tmtextit{Mathematische Lehrb{\"u}cher und Monographien, II. Abteilung:
  Mathematische Monographien [Mathematical Textbooks and Monographs, Part II:
  Mathematical Monographs]}. {\newblock}Akademie-Verlag, Berlin, 1992.
  {\newblock}Mathematical aspects of algebraic quantum field
  theory.{\newblock}
  
  \bibitem[17]{Belavkin}V.~P.~Belavkin. {\newblock}Quantum stochastic calculus
  and quantum nonlinear filtering. {\newblock}\tmtextit{J. Multivariate
  Anal.}, 42(2):171--201, 1992.{\newblock}
  
  \bibitem[18]{Berezin1}F.~A.~Berezin. {\newblock}\tmtextit{The method of
  second quantization}. {\newblock}Translated from the Russian by Nobumichi
  Mugibayashi and Alan Jeffrey. Pure and Applied Physics, Vol. 24. Academic
  Press, New York-London, 1966.{\newblock}
  
  \bibitem[19]{kirillov_introduction_1987}F.~A.~Berezin.
  {\newblock}\tmtextit{Introduction to Superanalysis}. {\newblock}Mathematical
  Physics and Applied Mathematics. Springer Netherlands, 1987.{\newblock}
  
  \bibitem[20]{Berezin2}F.~A.~Berezin. {\newblock}\tmtextit{Introduction to
  superanalysis},  volume~9  of \tmtextit{Mathematical Physics and Applied
  Mathematics}. {\newblock}D. Reidel Publishing Co., Dordrecht, 1987.
  {\newblock}Edited and with a foreword by A. A. Kirillov, With an appendix by
  V. I. Ogievetsky, Translated from the Russian by J. Niederle and R.
  Koteck{\'y}, Translation edited by Dimitri Le\u{\i}tes.{\newblock}
  
  \bibitem[21]{Bratteli1}O.~Bratteli  and  D.~W.~Robinson.
  {\newblock}\tmtextit{Operator algebras and quantum statistical mechanics.
  1}. {\newblock}Texts and Monographs in Physics. Springer-Verlag, New York,
  Second  edition, 1987. {\newblock}$C^{\ast}$- and $W^{\ast}$-algebras,
  symmetry groups, decomposition of states.{\newblock}
  
  \bibitem[22]{Bratteli2}O.~Bratteli  and  D.~W.~Robinson.
  {\newblock}\tmtextit{Operator algebras and quantum statistical mechanics.
  2}. {\newblock}Texts and Monographs in Physics. Springer-Verlag, Berlin,
  Second  edition, 1997. {\newblock}Equilibrium states. Models in quantum
  statistical mechanics.{\newblock}
  
  \bibitem[23]{Carlen1}E.~Carlen  and  P.~Kr{\'e}e. {\newblock}$L^p$ estimates
  on iterated stochastic integrals. {\newblock}\tmtextit{Ann. Probab.},
  19(1):354--368, 1991.{\newblock}
  
  \bibitem[24]{catellier_paracontrolled_2013}R.~Catellier  and  K.~Chouk.
  {\newblock}Paracontrolled distributions and the 3-dimensional stochastic
  quantization equation. {\newblock}\tmtextit{The Annals of Probability},
  46(5):2621--2679, 2018.
  {\newblock}\href{Https://doi.org/10.1214/17-AOP1235}{}10.1214/17-AOP1235.{\newblock}
  
  \bibitem[25]{Cipriani}F.~Cipriani. {\newblock}Noncommutative potential
  theory: a survey. {\newblock}\tmtextit{J. Geom. Phys.}, 105:25--59,
  2016.{\newblock}
  
  \bibitem[26]{damgaard_stochastic_1984}P.~Damgaard  and  K.~Tsokos.
  {\newblock}Stochastic quantization with fermions.
  {\newblock}\tmtextit{Nuclear Physics B}, 235(1):75--92, may 1984.{\newblock}
  
  \bibitem[27]{de_vecchi_note_2019}F.~C.~De Vecchi  and  M.~Gubinelli.
  {\newblock}A note on supersymmetry and stochastic differential equations.
  {\newblock}\tmtextit{ArXiv:1912.04830 [math-ph]}, dec 2019.
  {\newblock}ArXiv: 1912.04830.{\newblock}
  
  \bibitem[28]{DeWitt}B.~DeWitt. {\newblock}\tmtextit{Supermanifolds}.
  {\newblock}Cambridge Monographs on Mathematical Physics. Cambridge
  University Press, Cambridge, Second  edition, 1992.{\newblock}
  
  \bibitem[29]{dirksen_noncommutative_2010}S.~Dirksen.
  {\newblock}Noncommutative stochastic integration through decoupling.
  {\newblock}\tmtextit{Journal of Mathematical Analysis and Applications},
  370(1):200--223, oct 2010.{\newblock}
  
  \bibitem[30]{disertori_continuous_1998}M.~Disertori  and  V.~Rivasseau.
  {\newblock}Continuous Constructive Fermionic Renormalization.
  {\newblock}\tmtextit{ArXiv:hep-th/9802145}, apr 1998.{\newblock}
  
  \bibitem[31]{efremov_stochastic_2019}A.~N.~Efremov. {\newblock}Stochastic
  Quantization of Massive Fermions. {\newblock}\tmtextit{International Journal
  of Theoretical Physics}, 58(4):1150--1156, apr 2019.{\newblock}
  
  \bibitem[32]{Emch}G.~G.~Emch. {\newblock}\tmtextit{Mathematical and
  conceptual foundations of 20th-century physics},  volume  100  of
  \tmtextit{North-Holland Mathematics Studies}. {\newblock}North-Holland
  Publishing Co., Amsterdam, 1984. {\newblock}Notas de Matem{\'a}tica
  [Mathematical Notes], 100.{\newblock}
  
  \bibitem[33]{MR1905424}J.~Feldman, H.~Kn{\"o}rrer, and  E.~Trubowitz.
  {\newblock}\tmtextit{Fermionic functional integrals and the renormalization
  group},  volume~16  of \tmtextit{CRM Monograph Series}. {\newblock}American
  Mathematical Society, Providence, RI, 2002.{\newblock}
  
  \bibitem[34]{feldman_renormalizable_1986}J.~Feldman, J.~Magnen,
  V.~Rivasseau, and  R.~S{\'e}n{\'e}or. {\newblock}A renormalizable field
  theory: The massive Gross-Neveu model in two dimensions.
  {\newblock}\tmtextit{Commun.Math. Phys.}, 103(1):67--103, mar
  1986.{\newblock}
  
  \bibitem[35]{fukai_stochastic_1983}T.~Fukai, H.~Nakazato, I.~Ohba, K.~Okano,
  and  Y.~Yamanaka. {\newblock}Stochastic Quantization Method of Fermion
  Fields. {\newblock}\tmtextit{Progress of Theoretical Physics},
  69(5):1600--1616, may 1983.{\newblock}
  
  \bibitem[36]{Gelfand}I.~Gelfand  and  M.~Neumark. {\newblock}On the
  imbedding of normed rings into the ring of operators in Hilbert space.
  {\newblock}\tmtextit{Rec. Math. [Mat. Sbornik] N.S.}, 12(54):197--213,
  1943.{\newblock}
  
  \bibitem[37]{Gordina}M.~Gordina. {\newblock}Stochastic differential
  equations on noncommutative $L^2$. {\newblock}In \tmtextit{Finite and
  infinite dimensional analysis in honor of Leonard Gross (New Orleans, LA,
  2001)},  volume  317  of \tmtextit{Contemp. Math.},  pages  87--98. Amer.
  Math. Soc., Providence, RI, 2003.{\newblock}
  
  \bibitem[38]{Gross1967}L.~Gross. {\newblock}Abstract Wiener spaces.
  {\newblock}In \tmtextit{Proc. Fifth Berkeley Sympos. Math. Statist. and
  Probability (Berkeley, Calif., 1965/66), Vol. II: Contributions to
  Probability Theory, Part 1},  pages  31--42. Univ. California Press,
  Berkeley, Calif., 1967.{\newblock}
  
  \bibitem[39]{gubinelli_pde_2018}M.~Gubinelli  and  M.~Hofmanov{\'a}.
  {\newblock}A PDE construction of the Euclidean \$\textbackslash
  Phi{\textasciicircum}4\_3\$ quantum field theory.
  {\newblock}\tmtextit{ArXiv:1810.01700 [math-ph]}, oct 2018.
  {\newblock}ArXiv: 1810.01700.{\newblock}
  
  \bibitem[40]{gubinelli_global_2019}M.~Gubinelli  and  M.~Hofmanov{\'a}.
  {\newblock}Global Solutions to Elliptic and Parabolic $\Phi^4$ Models in
  Euclidean Space. {\newblock}\tmtextit{Communications in Mathematical
  Physics}, 368(3):1201--1266, 2019.{\newblock}
  
  \bibitem[41]{Haag1}R.~Haag. {\newblock}\tmtextit{Local quantum physics}.
  {\newblock}Texts and Monographs in Physics. Springer-Verlag, Berlin, Second 
  edition, 1996. {\newblock}Fields, particles, algebras.{\newblock}
  
  \bibitem[42]{hairer_theory_2014}M.~Hairer. {\newblock}A theory of regularity
  structures. {\newblock}\tmtextit{Inventiones mathematicae}, 198(2):269--504,
  2014.
  {\newblock}\href{Https://doi.org/10.1007/s00222-014-0505-4}{}10.1007/s00222-014-0505-4.{\newblock}
  
  \bibitem[43]{Holevo}A.~Holevo. {\newblock}\tmtextit{Probabilistic and
  statistical aspects of quantum theory},  volume~1  of
  \tmtextit{Quaderni/Monographs}. {\newblock}Edizioni della Normale, Pisa,
  Second  edition, 2011. {\newblock}With a foreword from the second Russian
  edition by K. A. Valiev.{\newblock}
  
  \bibitem[44]{hsu_stochastic_2002}E.~P.~Hsu. {\newblock}\tmtextit{Stochastic
  Analysis on Manifolds}. {\newblock}Graduate Studies in Mathematics 38.
  American Mathematical Society, 2002.{\newblock}
  
  \bibitem[45]{ignatyuk_convergence_1993}I.~A.~Ignatyuk, V.~A.~Malyshev, and 
  V.~Sidoravi{\v c}ius. {\newblock}Convergence of a Method of the Stochastic
  Quantization II. {\newblock}\tmtextit{Theory of Probability \& Its
  Applications}, 37(4):599--620, jan 1993.{\newblock}
  
  \bibitem[46]{Itzykson}C.~Itzykson  and  J.~B.~Zuber.
  {\newblock}\tmtextit{Quantum field theory}. {\newblock}McGraw-Hill
  International Book Co., New York, 1980. {\newblock}International Series in
  Pure and Applied Physics.{\newblock}
  
  \bibitem[47]{janson_gaussian_1997}S.~Janson. {\newblock}\tmtextit{Gaussian
  Hilbert Spaces}. {\newblock}Cambridge University Press, jun 1997.{\newblock}
  
  \bibitem[48]{jona_lasinio_stochastic_1985}G.~Jona-Lasinio  and 
  P.~K.~Mitter. {\newblock}On the stochastic quantization of field theory.
  {\newblock}\tmtextit{Communications in Mathematical Physics (1965-1997)},
  101(3):409--436, 1985.{\newblock}
  
  \bibitem[49]{Kadison2}R.~V.~Kadison  and  J.~R.~Ringrose.
  {\newblock}\tmtextit{Fundamentals of the theory of operator algebras. Vol.
  I},  volume~15  of \tmtextit{Graduate Studies in Mathematics}.
  {\newblock}American Mathematical Society, Providence, RI, 1997.
  {\newblock}Elementary theory, Reprint of the 1983 original.{\newblock}
  
  \bibitem[50]{Kadison1}R.~V.~Kadison  and  J.~R.~Ringrose.
  {\newblock}\tmtextit{Fundamentals of the theory of operator algebras. Vol.
  II},  volume~16  of \tmtextit{Graduate Studies in Mathematics}.
  {\newblock}American Mathematical Society, Providence, RI, 1997.
  {\newblock}Advanced theory, Corrected reprint of the 1986
  original.{\newblock}
  
  \bibitem[51]{kakudo1983gauge}Y.~Kakudo, Y.~Taguchi, A.~Tanaka, and 
  K.~Yamamoto. {\newblock}Gauge-independent calculation of s-matrix elements
  in quantum electrodynamics. {\newblock}\tmtextit{Progress of Theoretical
  Physics (Kyoto)}, 69(4):1225--1235, 1983.{\newblock}
  
  \bibitem[52]{klein_supersymmetry_1985}A.~Klein. {\newblock}Supersymmetry and
  a two-dimensional reduction in random phenomena. {\newblock}In
  \tmtextit{Quantum probability and applications, II (Heidelberg, 1984)}, 
  volume  1136  of \tmtextit{Lecture Notes in Math.},  pages  306--317.
  Springer, Berlin, 1985.{\newblock}
  
  \bibitem[53]{klein_supersymmetry_1983}A.~Klein  and  J.~Fernando Perez.
  {\newblock}Supersymmetry and dimensional reduction: A non-perturbative
  proof. {\newblock}\tmtextit{Physics Letters B}, 125(6):473--475, jun
  1983.{\newblock}
  
  \bibitem[54]{Kummer}B.~K{\"u}mmerer. {\newblock}Survey on a theory of
  noncommutative stationary Markov processes. {\newblock}In \tmtextit{Quantum
  probability and applications, III (Oberwolfach, 1987)},  volume  1303  of
  \tmtextit{Lecture Notes in Math.},  pages  154--182. Springer, Berlin,
  1988.{\newblock}
  
  \bibitem[55]{kupiainen_renormalization_2016}A.~Kupiainen.
  {\newblock}Renormalization Group and Stochastic PDEs.
  {\newblock}\tmtextit{Annales Henri Poincar{\'e}}, 17(3):497--535, 2016.
  {\newblock}\href{Https://doi.org/10.1007/s00023-015-0408-y}{}10.1007/s00023-015-0408-y.{\newblock}
  
  \bibitem[56]{Leites}D.~A.~Le\u{\i}tes. {\newblock}Introduction to the theory of
  supermanifolds. {\newblock}\tmtextit{Uspekhi Mat. Nauk}, 35(1(211)):3--57,
  1980.{\newblock}
  
  \bibitem[57]{mastropietro_non-perturbative_2008}V.~Mastropietro.
  {\newblock}\tmtextit{Non-Perturbative Renormalization}. {\newblock}World
  Scientific Publishing Co Pte Ltd, Hackensack, NJ, 2008.{\newblock}
  
  \bibitem[58]{Meyerbook}P.-A.~Meyer. {\newblock}\tmtextit{Quantum probability
  for probabilists},  volume  1538  of \tmtextit{Lecture Notes in
  Mathematics}. {\newblock}Springer-Verlag, Berlin, 1993.{\newblock}
  
  \bibitem[59]{moinat_space_time_2018}A.~Moinat  and  H.~Weber.
  {\newblock}Space-time localisation for the dynamic $\Phi^4_3$ model.
  {\newblock}\tmtextit{ArXiv:1811.05764}, nov 2018. {\newblock}ArXiv:
  1811.05764.{\newblock}
  
  \bibitem[60]{MWcomedown}J.-C.~Mourrat  and  H.~Weber. {\newblock}The dynamic
  $\Phi^4_3$ model comes down from infinity. {\newblock}\tmtextit{Comm. Math.
  Phys.}, 356(3):673--753, 2017.{\newblock}
  
  \bibitem[61]{Murray1}F.~J.~Murray  and  J.~Von Neumann. {\newblock}On rings
  of operators. {\newblock}\tmtextit{Ann. of Math. (2)}, 37(1):116--229,
  1936.{\newblock}
  
  \bibitem[62]{Murray2}F.~J.~Murray  and  J.~von~Neumann. {\newblock}On rings
  of operators. II. {\newblock}\tmtextit{Trans. Amer. Math. Soc.},
  41(2):208--248, 1937.{\newblock}
  
  \bibitem[63]{Murray3}F.~J.~Murray  and  J.~von~Neumann. {\newblock}On rings
  of operators. IV. {\newblock}\tmtextit{Ann. of Math. (2)}, 44:716--808,
  1943.{\newblock}
  
  \bibitem[64]{Naimark}M.~A.~Na\u{\i}mark. {\newblock}\tmtextit{Normed algebras}.
  {\newblock}Wolters-Noordhoff Publishing, Groningen, Third  edition, 1972.
  {\newblock}Translated from the second Russian edition by Leo F. Boron,
  Wolters-Noordhoff Series of Monographs and Textbooks on Pure and Applied
  Mathematics.{\newblock}
  
  \bibitem[65]{osterwalder_euclidean_1973}K.~Osterwalder. {\newblock}Euclidean
  fermi fields. {\newblock}In \tmtextit{Constructive Quantum Field Theory},
  Lecture Notes in Physics,  pages  326--331. Springer, Berlin, Heidelberg,
  1973.{\newblock}
  
  \bibitem[66]{osterwalder_euclidean_1973-2}K.~Osterwalder  and  R.~Schrader.
  {\newblock}Euclidean Fermi fields and a Feynman-Kac formula for
  Boson-Fermions models. {\newblock}\tmtextit{Helvetica Physica Acta},
  46:277--302, 1973.{\newblock}
  
  \bibitem[67]{parisi_perturbation_1981}G.~Parisi  and  Y.~S.~Wu.
  {\newblock}Perturbation theory without gauge fixing.
  {\newblock}\tmtextit{Scientia Sinica. Zhongguo Kexue}, 24(4):483--496,
  1981.{\newblock}
  
  \bibitem[68]{rogers_graded_1986}A.~Rogers. {\newblock}Graded manifolds,
  supermanifolds and infinite-dimensional Grassmann algebras.
  {\newblock}\tmtextit{Communications in Mathematical Physics},
  105(3):375--384, sep 1986.{\newblock}
  
  \bibitem[69]{Rogers2007}A.~Rogers. {\newblock}\tmtextit{Supermanifolds}.
  {\newblock}World Scientific Publishing Co. Pte. Ltd., Hackensack, NJ, 2007.
  {\newblock}Theory and applications.{\newblock}
  
  \bibitem[70]{salmhofer_renormalization}M.~Salmhofer.
  {\newblock}\tmtextit{Renormalization: An Introduction}. {\newblock}Springer,
  Berlin ; New York, 1st Corrected ed. 1999, Corr. 2nd printing 2007 edition 
  edition, sep 2007.{\newblock}
  
  \bibitem[71]{Segal1}I.~E.~Segal. {\newblock}A non-commutative extension of
  abstract integration. {\newblock}\tmtextit{Ann. of Math. (2)}, 57:401--457,
  1953.{\newblock}
  
  \bibitem[72]{Segalbook2}I.~E.~Segal. {\newblock}\tmtextit{Mathematical
  problems of relativistic physics},  volume  1960  of \tmtextit{With an
  appendix by George W. Mackey. Lectures in Applied Mathematics (proceedings
  of the Summer Seminar, Boulder, Colorado}. {\newblock}American Mathematical
  Society, Providence, R.I., 1963.{\newblock}
  
  \bibitem[73]{shcherbakov_elements_1993-2}V.~V.~Shcherbakov.
  {\newblock}Elements of stochastic analysis for the case of Grassmann
  variables. I. Grassmann stochastic integrals and random processes.
  {\newblock}\tmtextit{Theoretical and Mathematical Physics}, 96(1):792--800,
  jul 1993.{\newblock}
  
  \bibitem[74]{shcherbakov_elements_1993-1}V.~V.~Shcherbakov.
  {\newblock}Elements of stochastic analysis for the case of Grassmann
  variables. II. Stochastic partial differential equations for Grassmann
  processes. {\newblock}\tmtextit{Theoretical and Mathematical Physics},
  97(2):1229--1235, nov 1993.{\newblock}
  
  \bibitem[75]{shcherbakov_elements_1993}V.~V.~Shcherbakov.
  {\newblock}Elements of stochastic analysis for the case of Grassmann
  variables. III. Correlation functions. {\newblock}\tmtextit{Theoretical and
  Mathematical Physics}, 97(3):1323--1332, dec 1993.{\newblock}
  
  \bibitem[76]{she-sheng_stochastic_1985}X.~She-Sheng  and  H.~Ting-chang.
  {\newblock}Stochastic quantization of fermions on lattice.
  {\newblock}\tmtextit{Chinese Physics Letters}, 2(10):474--476, oct
  1985.{\newblock}
  
  \bibitem[77]{Sinha}K.~B.~Sinha  and  D.~Goswami.
  {\newblock}\tmtextit{Quantum stochastic processes and noncommutative
  geometry},  volume  169  of \tmtextit{Cambridge Tracts in Mathematics}.
  {\newblock}Cambridge University Press, Cambridge, 2007.{\newblock}
  
  \bibitem[78]{summers2012perspective}S.~J.~Summers. {\newblock}A perspective
  on constructive quantum field theory. {\newblock}\tmtextit{ArXiv preprint
  arXiv:1203.3991}, 2012.{\newblock}
  
  \bibitem[79]{TriebelIII}H.~Triebel. {\newblock}\tmtextit{Theory of function
  spaces. III},  volume  100  of \tmtextit{Monographs in Mathematics}.
  {\newblock}Birkh{\"a}user Verlag, Basel, 2006.{\newblock}
  
  \bibitem[80]{van2003algebra1}B.~van der~Waerden.
  {\newblock}\tmtextit{Algebra: Volume I}. {\newblock}Algebra : based in part
  on lectures by E. Artin and E. Noether. Springer New York, 2003.{\newblock}
  
  \bibitem[81]{van2003algebra2}B.~van der~Waerden.
  {\newblock}\tmtextit{Algebra: Volume II}. {\newblock}Algebra : based in part
  on lectures by E. Artin and E. Noether. Springer New York, 2003.{\newblock}
  
  \bibitem[82]{VonNeumann}J.~von~Neumann. {\newblock}On rings of operators.
  III. {\newblock}\tmtextit{Ann. of Math. (2)}, 41:94--161, 1940.{\newblock}
  
  \bibitem[83]{MR0448079}W.~{\.Z}elazko. {\newblock}\tmtextit{Banach
  algebras}. {\newblock}Elsevier Publishing Co., Amsterdam-London-New York;
  PWN--Polish Scientific Publishers, Warsaw, 1973. {\newblock}Translated from
  the Polish by Marcin E. Kuczma.{\newblock}
  
  \bibitem[84]{ZinnJustin}J.~Zinn-Justin. {\newblock}\tmtextit{Quantum field
  theory and critical phenomena},  volume~85  of \tmtextit{International
  Series of Monographs on Physics}. {\newblock}The Clarendon Press, Oxford
  University Press, New York, Second  edition, 1993. {\newblock}Oxford Science
  Publications.{\newblock}
\end{thebibliography}

\end{document}
